% Nombres complexes : approche géométrique — Mathématiques Tle C — Cours signé OnBuch+
% Programme officiel MINESEC. Compiler avec : tectonic M13-nombres-complexes-approche-geometrique.tex
\newcommand{\DOCMATIERE}{Mathématiques}
\newcommand{\DOCNIVEAU}{Tle C}
\newcommand{\DOCMODULE}{Configurations et transformations élémentaires du plan}
\newcommand{\DOCLECON}{13}
\newcommand{\DOCTITRE}{Nombres complexes : approche géométrique}
% OnBuch+ — préambule « cours signés » ENRICHI (compilé avec tectonic / XeLaTeX)
% Système visuel « Pop » d'OnBuch+ : fond crème (jamais blanc), bordures encre
% épaisses, ombres « dures » décalées, titres Archivo Black en capitales.
% Version MATHÉMATIQUES : graphes (pgfplots/tikz), boîtes pédagogiques
% (définition, propriété, méthode, attention, le savais-tu, exercice résolu,
% exercice, TP, bilan), page de garde et signature OnBuch+ ancrée en bas.
%
% Macros attendues dans main.tex AVANT \input{preamble} :
%   \DOCMATIERE  \DOCNIVEAU  \DOCMODULE  \DOCLECON (numéro)  \DOCTITRE
\documentclass[11pt]{article}

\usepackage{fontspec}
\usepackage[french]{babel}
\usepackage[a4paper,margin=2cm,top=3.4cm,bottom=2.8cm,headheight=1.4cm,footskip=1.3cm]{geometry}
\usepackage{amsmath,amssymb,amsfonts}
\usepackage{mathtools} % \xmapsto, \coloneqq... souvent utilisés par les modèles
\usepackage{cancel} % simplifications barrées
% \abs{z} (module, valeur absolue) et \norm{u} (norme), utilisés par les modèles
\providecommand{\abs}{}\let\abs\relax\DeclarePairedDelimiter{\abs}{\lvert}{\rvert}
\providecommand{\norm}{}\let\norm\relax\DeclarePairedDelimiter{\norm}{\lVert}{\rVert}
\usepackage[table]{xcolor} % [table] : fournit aussi \rowcolors (alternance de lignes)
\usepackage{pagecolor}
\usepackage{tikz}
\usetikzlibrary{arrows.meta,positioning,calc,shapes.geometric,shapes.misc,shapes.arrows,decorations.pathmorphing,decorations.markings,patterns,shadows,mindmap,trees,fit,backgrounds,matrix,intersections}
\usepackage{pgfplots}
\usepackage{pgf-pie} % diagrammes circulaires \pie (statistiques)
\pgfplotsset{compat=1.18}
\usepgfplotslibrary{fillbetween,groupplots} % aires entre courbes (calcul intégral)
\usepackage{enumitem}
\usepackage{fancyhdr}
% [most] charge toutes les bibliothèques tcolorbox (listings, minted, raster,
% documentation...) et gonfle inutilement la mémoire TeX sur un document long
% (~300 boîtes) : on ne charge que ce qui est réellement utilisé.
\usepackage[skins,breakable]{tcolorbox}
\usepackage{graphicx}
\usepackage{colortbl}
\usepackage{booktabs}
\usepackage{tabularx}
\usepackage{diagbox} % case d'en-tête diagonale des tableaux à double entrée
% Colonnes extensibles centrée / gauche / droite (C, L, R), souvent utilisées par les modèles
\newcolumntype{C}{>{\centering\arraybackslash}X}
\newcolumntype{L}{>{\raggedright\arraybackslash}X}
\newcolumntype{R}{>{\raggedleft\arraybackslash}X}
\usepackage{array}
\usepackage{multicol}
\usepackage{siunitx}
% parse-numbers=false : accepte aussi \SI{2,5 \times 10^{-3}}{...}
\sisetup{locale=FR,output-decimal-marker={,},group-minimum-digits=4,parse-numbers=false}
\DeclareSIUnit\ohm{\ensuremath{\symup{\Omega}}} % U+2126 absent de la police
\usepackage{qrcode}
% \degree : commande fréquemment utilisée par les modèles pour un angle ou une
% température en dehors de \SI{}{} (non fournie par les paquets chargés).
\providecommand{\degree}{^{\circ}}
% \qed (fin de démonstration), utilisé par les modèles sans amsthm
\providecommand{\qed}{\hfill$\square$}
% \barre : double barre des valeurs interdites dans un tableau de variations (tkz-tab)
\providecommand{\barre}{\ensuremath{\|}}
% environnement proof (amsthm non chargé) utilisé par les modèles
\ifdefined\proof\else\newenvironment{proof}[1][Démonstration]{\par\noindent\textit{#1.}\ }{\qed\par}\fi
\usepackage{lastpage}
\usepackage{hyperref}
\hypersetup{hidelinks}

% ============================================================
% Palette « Pop » OnBuch+ — mêmes hex que la classe P (plus_ui.dart)
% ============================================================
\definecolor{poporange}{HTML}{F59321}
\definecolor{popdark}{HTML}{E07A0C}
\definecolor{popcream}{HTML}{FFF6E8}
\definecolor{popink}{HTML}{1C1714}
\definecolor{poppurple}{HTML}{7A5AE0}
\definecolor{popgreen}{HTML}{2E9E6A}
\definecolor{poppink}{HTML}{E0457B}
\definecolor{popblue}{HTML}{2F7DE1}
\definecolor{popgold}{HTML}{C98A12}
\definecolor{popmuted}{HTML}{8A7F76}
\definecolor{popline}{HTML}{EADFD2}
\definecolor{popgray}{HTML}{8A7F76}
\colorlet{popgrey}{popgray}
% Alias tolérants (le modèle nomme parfois les couleurs différemment)
\colorlet{popwhite}{white}
\colorlet{popblack}{popink}
\colorlet{popred}{poppink}
\colorlet{popyellow}{popgold}
% Teintes claires pour fonds de tableaux / zones
\colorlet{popblueL}{popblue!10!white}
\colorlet{poporangeL}{poporange!14!white}
\colorlet{popgreenL}{popgreen!12!white}
\colorlet{poppurpleL}{poppurple!12!white}
\colorlet{poppinkL}{poppink!12!white}
\colorlet{popgoldL}{popgold!14!white}

\pagecolor{popcream}

% ============================================================
% Typographie
% ============================================================
\defaultfontfeatures{Ligatures=TeX}
\setmainfont{PlusJakartaSans-Medium.ttf}[Path=../../../fonts/,BoldFont=PlusJakartaSans-Bold.ttf]
% Maths en sans-serif (Fira Math) avec chiffres et lettres droites de Plus
% Jakarta, pour que les formules se fondent dans le texte.
\usepackage[math-style=ISO,bold-style=ISO]{unicode-math}
\setmathfont{FiraMath-Regular.otf}[Path=../../../fonts/]
\setmathfont{PlusJakartaSans-Medium.ttf}[Path=../../../fonts/,range={up/num,up/{Latin,latin},\mathup}]
\usepackage{icomma} % virgule décimale sans espace en mode math
% Caractères mathématiques Unicode absents des polices de texte (Plus Jakarta,
% Archivo Black) : rendus par leur équivalent mathématique, en texte comme en
% formule — sinon ils s'affichent en carré vide (ex. « Arithmétique dans ℤ »).
\usepackage{newunicodechar}
\newunicodechar{ℝ}{\ensuremath{\mathbb{R}}}
\newunicodechar{ℤ}{\ensuremath{\mathbb{Z}}}
\newunicodechar{ℂ}{\ensuremath{\mathbb{C}}}
\newunicodechar{ℕ}{\ensuremath{\mathbb{N}}}
\newunicodechar{ℚ}{\ensuremath{\mathbb{Q}}}
\newunicodechar{𝔻}{\ensuremath{\mathbb{D}}}
\newunicodechar{α}{\ensuremath{\alpha}}
\newunicodechar{β}{\ensuremath{\beta}}
\newunicodechar{γ}{\ensuremath{\gamma}}
\newunicodechar{δ}{\ensuremath{\delta}}
\newunicodechar{ε}{\ensuremath{\varepsilon}}
\newunicodechar{θ}{\ensuremath{\theta}}
\newunicodechar{λ}{\ensuremath{\lambda}}
\newunicodechar{μ}{\ensuremath{\mu}}
\newunicodechar{σ}{\ensuremath{\sigma}}
\newunicodechar{τ}{\ensuremath{\tau}}
\newunicodechar{φ}{\ensuremath{\varphi}}
\newunicodechar{ψ}{\ensuremath{\psi}}
\newunicodechar{ω}{\ensuremath{\omega}}
\newunicodechar{Σ}{\ensuremath{\Sigma}}
% majuscules produites par \MakeUppercase dans les titres (α-aminés -> Α)
\newunicodechar{Α}{\ensuremath{\alpha}}
\newunicodechar{Β}{\ensuremath{\beta}}
\newunicodechar{Γ}{\ensuremath{\Gamma}}
\newunicodechar{Δ}{\ensuremath{\Delta}}
\newunicodechar{Ε}{\ensuremath{\varepsilon}}
\newunicodechar{Θ}{\ensuremath{\Theta}}
\newunicodechar{Λ}{\ensuremath{\Lambda}}
\newunicodechar{Μ}{\ensuremath{\mu}}
\newunicodechar{Τ}{\ensuremath{\tau}}
\newunicodechar{Φ}{\ensuremath{\Phi}}
\newunicodechar{Ψ}{\ensuremath{\Psi}}
\newunicodechar{Ω}{\ensuremath{\Omega}}
\newunicodechar{↦}{\ensuremath{\mapsto}}
\newunicodechar{∈}{\ensuremath{\in}}
\newunicodechar{∉}{\ensuremath{\notin}}
\newunicodechar{∧}{\ensuremath{\wedge}}
\newunicodechar{⇔}{\ensuremath{\Leftrightarrow}}
\newunicodechar{⟺}{\ensuremath{\Longleftrightarrow}}
\newunicodechar{⇒}{\ensuremath{\Rightarrow}}
\newunicodechar{⟹}{\ensuremath{\Longrightarrow}}
\newunicodechar{∘}{\ensuremath{\circ}}
\newunicodechar{∪}{\ensuremath{\cup}}
\newunicodechar{∩}{\ensuremath{\cap}}
\newunicodechar{≡}{\ensuremath{\equiv}}
\newunicodechar{⊂}{\ensuremath{\subset}}
\newunicodechar{⊥}{\ensuremath{\perp}}
\newunicodechar{∃}{\ensuremath{\exists}}
\newunicodechar{∀}{\ensuremath{\forall}}
\newunicodechar{∛}{\ensuremath{\sqrt[3]{\,}}}
\newunicodechar{∜}{\ensuremath{\sqrt[4]{\,}}}
\newunicodechar{⃗}{\ensuremath{{}^{\to}}}
\newunicodechar{ⁿ}{\ifmmode^{n}\else\textsuperscript{\ensuremath{n}}\fi}
\newunicodechar{ˣ}{\ifmmode^{x}\else\textsuperscript{\ensuremath{x}}\fi}
\newunicodechar{ᵇ}{\ifmmode^{b}\else\textsuperscript{\ensuremath{b}}\fi}
\newunicodechar{ʳ}{\ifmmode^{r}\else\textsuperscript{\ensuremath{r}}\fi}
\newunicodechar{ᵘ}{\ifmmode^{u}\else\textsuperscript{\ensuremath{u}}\fi}
\newunicodechar{ʸ}{\ifmmode^{y}\else\textsuperscript{\ensuremath{y}}\fi}
\newunicodechar{ᵃ}{\ifmmode^{a}\else\textsuperscript{\ensuremath{a}}\fi}
\newunicodechar{ᵉ}{\ifmmode^{e}\else\textsuperscript{\ensuremath{e}}\fi}
\newunicodechar{ᵏ}{\ifmmode^{k}\else\textsuperscript{\ensuremath{k}}\fi}
\newunicodechar{ᵖ}{\ifmmode^{p}\else\textsuperscript{\ensuremath{p}}\fi}
\newunicodechar{ᴺ}{\ifmmode^{N}\else\textsuperscript{\ensuremath{N}}\fi}
\newunicodechar{ᶜ}{\ifmmode^{c}\else\textsuperscript{\ensuremath{c}}\fi}
\newunicodechar{ʰ}{\ifmmode^{h}\else\textsuperscript{\ensuremath{h}}\fi}
\newunicodechar{ᵠ}{\ifmmode^{q}\else\textsuperscript{\ensuremath{q}}\fi}
\newunicodechar{ₙ}{\ifmmode_{n}\else\textsubscript{\ensuremath{n}}\fi}
\newunicodechar{ₐ}{\ifmmode_{a}\else\textsubscript{\ensuremath{a}}\fi}
\newunicodechar{ᵢ}{\ifmmode_{i}\else\textsubscript{\ensuremath{i}}\fi}
\newunicodechar{ᵣ}{\ifmmode_{r}\else\textsubscript{\ensuremath{r}}\fi}
\newunicodechar{ₖ}{\ifmmode_{k}\else\textsubscript{\ensuremath{k}}\fi}
\newunicodechar{ᵦ}{\ifmmode_{\beta}\else\textsubscript{\ensuremath{\beta}}\fi}
\newunicodechar{ₓ}{\ifmmode_{x}\else\textsubscript{\ensuremath{x}}\fi}
\newfontfamily\popdisplay{ArchivoBlack-Regular.ttf}[Path=../../../fonts/]
\newfontfamily\poplogo{SpaceGrotesk-Bold.ttf}[Path=../../../fonts/]
\newfontfamily\popsemi{PlusJakartaSans-SemiBold.ttf}[Path=../../../fonts/]
\newfontfamily\popxbold{PlusJakartaSans-ExtraBold.ttf}[Path=../../../fonts/]
\newfontfamily\popxboldls{PlusJakartaSans-ExtraBold.ttf}[Path=../../../fonts/,LetterSpace=3.0]

\setlist[enumerate]{leftmargin=1.7em,itemsep=5pt,topsep=4pt}
\setlist[itemize]{leftmargin=1.7em,itemsep=4pt,topsep=4pt,label={\color{poporange}\textbullet}}
\setlength{\parindent}{0pt}
\setlength{\parskip}{5pt}
\renewcommand{\arraystretch}{1.25}

% pgfplots : style Pop par défaut
\pgfplotsset{
  every axis/.append style={axis line style={popink,line width=1pt},tick style={popink},
    grid style={popline,line width=0.5pt},label style={font=\small},tick label style={font=\footnotesize}},
  popcurve/.style={poporange,line width=1.8pt,no markers,smooth},
}
% Même style disponible dans un \draw TikZ simple (hors pgfplots)
\tikzset{popcurve/.style={poporange,line width=1.8pt}}

% ============================================================
% Pill / squiggle
% ============================================================
\newcommand{\poppill}[3]{%
  \tikz[baseline=-0.55ex]\node[draw=popink,line width=1.2pt,rounded corners=2.6pt,%
    fill=#2,inner xsep=6.5pt,inner ysep=3pt]{%
    \popxboldls\fontsize{7}{7}\selectfont\color{#3}\MakeUppercase{#1}};%
}
\newcommand{\popsquiggle}[1]{%
  \tikz[baseline=-0.4ex]{
    \draw[#1,line width=2.6pt,line cap=round]
      plot [smooth] coordinates {(0,0.09) (0.34,0.01) (0.68,0.21) (1.02,0.01) (1.36,0.10)};
  }%
}
% Mot-clé surligné (marqueur orange sous le texte)
\newcommand{\cle}[1]{{\popxbold\color{popdark}#1}}

% ============================================================
% En-tête / pied
% ============================================================
\pagestyle{fancy}
\fancyhf{}
\renewcommand{\headrulewidth}{0pt}
\renewcommand{\footrulewidth}{0pt}
\lhead{\raisebox{-0.15ex}{\poplogo\fontsize{13}{13}\selectfont\color{popink}OnBuch\color{poporange}+}}
\rhead{\poppill{\DOCMATIERE\ \textbullet\ \DOCNIVEAU\ \textbullet\ Leçon \DOCLECON}{white}{popink}}
\lfoot{\popxbold\fontsize{7.6}{7.6}\selectfont\color{popmuted}\MakeUppercase{Cours signé OnBuch+}\ \textbullet\ {\popsemi\DOCTITRE}}
\rfoot{\tikz[baseline=-0.6ex]\node[rounded corners=8.5pt,fill=popink,minimum height=17pt,minimum width=17pt,inner xsep=5pt,inner sep=0]{\popxbold\fontsize{8}{8}\selectfont\color{popcream}\thepage\,/\,\pageref*{LastPage}};}

% ============================================================
% Titres
% ============================================================
\newcommand{\chaptitre}[2]{%
  \begin{tcolorbox}[enhanced,colback=poporange,colframe=popink,boxrule=2.6pt,arc=11pt,
      drop shadow=popink,left=15pt,right=15pt,top=12pt,bottom=11pt,before skip=3pt,after skip=18pt]
    \poppill{#2}{popink}{popcream}\par\vspace{9pt}
    {\raggedright\hyphenpenalty=10000\exhyphenpenalty=10000\popdisplay\fontsize{22}{24}\selectfont\color{popcream}\MakeUppercase{#1}\par}
  \end{tcolorbox}
}
% Section numérotée : pastille orange avec le numéro + titre Archivo
\newcounter{popsec}
\newcommand{\coursec}[1]{%
  \stepcounter{popsec}\setcounter{popsub}{0}%
  \par\vspace{16pt}\noindent
  \begin{minipage}{\linewidth}
  \tikz[baseline=-0.7ex]\node[draw=popink,line width=1.6pt,fill=poporange,rounded corners=4pt,minimum size=22pt,inner sep=0]{\popdisplay\fontsize{12}{12}\selectfont\color{popcream}\thepopsec};%
  \hspace{8pt}{\popdisplay\fontsize{15}{17}\selectfont\color{popink}\MakeUppercase{#1}}\par
  \vspace{4pt}\noindent{\color{poporange}\rule{36pt}{3.6pt}}
  \end{minipage}\par\nopagebreak\vspace{8pt}
}
\newcounter{popsub}
\newcommand{\courssub}[1]{%
  \stepcounter{popsub}%
  \par\vspace{9pt}\noindent{\popxbold\fontsize{11.5}{13}\selectfont\color{popink}{\color{poporange}\thepopsec.\thepopsub}\hspace{6pt}#1}\par\nopagebreak\vspace{4pt}
}

% ============================================================
% Boîtes pédagogiques — même recette : fond blanc, bordure épaisse colorée,
% ombre dure encre, badge plein. Titre optionnel [#1].
% ============================================================
\tcbset{popbase/.style={breakable,enhanced,colback=white,boxrule=2.3pt,arc=9pt,
  shadow={3pt}{-3pt}{0pt}{popink},left=14pt,right=14pt,top=11pt,bottom=11pt,before skip=11pt,after skip=15pt,
  fontupper=\popsemi\fontsize{10.6}{14.5}\selectfont\color{popink},
  fontlower=\popsemi\fontsize{10.6}{14.5}\selectfont\color{popink}}}
% Coche dessinée (le glyphe ✓ n'existe pas dans les polices du document)
\newcommand{\popcheck}[1]{\tikz[baseline=-0.5ex]\draw[#1,line width=1.6pt,line cap=round,line join=round](-0.13,0.0)--(-0.04,-0.09)--(0.13,0.1);}
\newcommand{\popboxhead}[3]{\poppill{#1}{#2}{white}\ifx\relax#3\relax\else\hspace{6pt}{\popxbold\fontsize{10.6}{12}\selectfont\color{popink}#3}\fi\par\vspace{7pt}}

\newtcolorbox{aretenir}[1][]{popbase,colframe=popgreen,before upper={\popboxhead{À retenir}{popgreen}{#1}}}
\newtcolorbox{exemplebox}[1][]{popbase,colframe=poporange,before upper={\popboxhead{Exemple}{poporange}{#1}}}
\newtcolorbox{definition}[1][]{popbase,colframe=popblue,before upper={\popboxhead{Définition}{popblue}{#1}}}
\newtcolorbox{propriete}[1][]{popbase,colframe=poppurple,before upper={\popboxhead{Propriété}{poppurple}{#1}}}
\newtcolorbox{methode}[1][]{popbase,colframe=popgold,before upper={\popboxhead{Méthode}{popgold}{#1}}}
\newtcolorbox{attention}[1][]{popbase,colframe=poppink,before upper={\popboxhead{Attention}{poppink}{#1}}}
\newtcolorbox{experience}[1][]{popbase,colframe=popblue,colback=popblueL,before upper={\popboxhead{Expérience}{popblue}{#1}}}
% « Le savais-tu ? » : carte violette pleine (contexte, culture, Cameroun)
\newtcolorbox{savaistu}[1][]{popbase,colback=poppurple,colframe=popink,
  fontupper=\popsemi\fontsize{10.4}{14.2}\selectfont\color{white},
  before upper={\poppill{Le savais-tu\,?}{white}{poppurple}\ifx\relax#1\relax\else\hspace{6pt}{\popxbold\color{white}#1}\fi\par\vspace{7pt}}}
% Exercice résolu : énoncé en haut, \tcblower puis la solution sur fond crème
\newcounter{popexo}\newcounter{popexr}
\newtcolorbox{exoresolu}[1][]{popbase,colframe=poporange,colbacklower=popcream,
  segmentation style={popink,line width=1.2pt,dashed},
  before upper={\stepcounter{popexr}\popboxhead{Exercice résolu \thepopexr}{poporange}{#1}},
  before lower={\poppill{Solution}{popink}{popcream}\par\vspace{6pt}}}
% Exercice (non résolu en place) — la correction est dans la partie Corrigés
\newtcolorbox{exercice}[1][]{popbase,colframe=popink,
  before upper={\stepcounter{popexo}\popboxhead{Exercice \thepopexo}{popink}{#1}}}
% Corrigé
\newtcolorbox{corrige}[1][]{popbase,colframe=popgreen,colback=popgreenL,
  before upper={\popboxhead{Corrigé}{popgreen}{#1}}}
% Objectifs / prérequis / bilan
\newtcolorbox{objectifs}{popbase,colframe=popink,colback=poporangeL,
  before upper={\popboxhead{Objectifs de la leçon}{popink}{}}}
\newtcolorbox{prerequis}{popbase,colframe=popink,colback=popblueL,
  before upper={\popboxhead{Prérequis}{popblue}{}}}
\newtcolorbox{bilan}{popbase,colframe=popink,colback=white,boxrule=2.8pt,
  before upper={\popboxhead{Fiche bilan}{poporange}{L'essentiel à connaître pour l'examen}}}
% Figure encadrée avec légende
\newtcolorbox{popfigure}[1][]{enhanced,breakable=false,colback=white,colframe=popline,boxrule=1.4pt,arc=8pt,
  left=10pt,right=10pt,top=10pt,bottom=8pt,before skip=10pt,after skip=12pt,halign=center,
  lower separated=false,halign lower=center,
  fontlower=\popsemi\fontsize{9}{11.5}\selectfont\color{popmuted},
  #1}
\newcommand{\legende}[1]{\par\vspace{6pt}{\popsemi\fontsize{9}{11.5}\selectfont\color{popmuted}#1}}

% ============================================================
% Signature OnBuch+
% ============================================================
\newcommand{\poplogoinline}[1]{{\poplogo\fontsize{#1}{#1}\selectfont\color{popink}OnBuch\color{poporange}+}}
\newcommand{\signatureonbuch}{%
  \begin{tcolorbox}[enhanced,colback=popink,colframe=popink,boxrule=2.2pt,arc=10pt,
      drop shadow=poporange,left=14pt,right=14pt,top=10pt,bottom=10pt,before skip=0pt,after skip=0pt]
    \tikz[baseline=-0.7ex]\node[circle,fill=popgreen,draw=popcream,line width=1.3pt,minimum size=22pt,inner sep=0]{\popcheck{white}};%
    \hspace{9pt}\begin{minipage}[c]{0.62\linewidth}
      {\popxbold\fontsize{12}{13}\selectfont\color{popcream}Signé\ \textbullet\ {\poplogo OnBuch\color{poporange}+}}\par\vspace{2pt}
      {\raggedright\popsemi\fontsize{8}{10.5}\selectfont\color{popline}\MakeUppercase{Équipe pédagogique OnBuch+}\\\MakeUppercase{Conforme au programme officiel MINESEC}\par}
    \end{minipage}\hfill
    {\poplogo\fontsize{22}{22}\selectfont\color{popcream}OnBuch\color{poporange}+}
  \end{tcolorbox}%
}

% ============================================================
% Page de garde
% #1 titre · #2 matière · #3 niveau · #4 module · #5 n° leçon · #6 durée ·
% #7 description / objectifs (texte ou itemize)
% ============================================================
\newcommand{\pagegarde}[7]{%
  \thispagestyle{empty}
  \newgeometry{a4paper,margin=1.7cm,top=1.6cm,bottom=1.4cm}%
  \begin{tikzpicture}[remember picture,overlay]
    \fill[poporange!18] ([xshift=-3.2cm,yshift=-3.6cm]current page.north east) circle (4.6cm);
    \fill[poppurple!14] ([xshift=2.4cm,yshift=7.6cm]current page.south west) circle (3.8cm);
  \end{tikzpicture}%
  {\poplogo\fontsize{30}{30}\selectfont\color{popink}OnBuch\color{poporange}+}\hfill
  \raisebox{0.6ex}{\poppill{Cours signé \textbullet\ Premium}{popink}{popcream}}\par
  \vspace{1pt}\popsquiggle{poppurple}\par
  \vspace{8pt}
  \begin{tcolorbox}[enhanced,colback=poporange,colframe=popink,boxrule=2.8pt,arc=13pt,
      drop shadow=popink,left=18pt,right=18pt,top=12pt,bottom=13pt,after skip=10pt]
    \poppill{#2 \textbullet\ #3}{popink}{popcream}\hspace{5pt}\poppill{#4}{white}{popink}\par\vspace{8pt}
    {\popdisplay\fontsize{14}{14}\selectfont\color{popink}LEÇON #5}\par\vspace{4pt}
    {\raggedright\hyphenpenalty=10000\exhyphenpenalty=10000\popdisplay\fontsize{25}{27}\selectfont\color{popcream}\MakeUppercase{#1}\par}
    \vspace{8pt}
    {\popxbold\fontsize{9.5}{9.5}\selectfont\color{popink}\MakeUppercase{Durée conseillée : #6}}
  \end{tcolorbox}
  \begin{tcolorbox}[enhanced,colback=white,colframe=popink,boxrule=2.2pt,arc=10pt,
      drop shadow=popink,left=16pt,right=16pt,top=10pt,bottom=10pt,after skip=8pt]
    \poppill{Dans cette leçon}{poppurple}{white}\par\vspace{5pt}
    \popsemi\fontsize{9.3}{12.6}\selectfont\color{popink}#7
  \end{tcolorbox}
  \noindent\begin{minipage}[t]{0.487\linewidth}
    \begin{tcolorbox}[enhanced,colback=popgreen,colframe=popink,boxrule=2.2pt,arc=10pt,
        drop shadow=popink,left=11pt,right=11pt,top=10pt,bottom=10pt,height=3.1cm,valign=center]
      \tcbox[colback=white,colframe=popink,boxrule=1.3pt,arc=5pt,boxsep=0pt,nobeforeafter,
        left=3pt,right=3pt,top=3pt,bottom=3pt,valign=center]{\qrcode[height=1.9cm]{https://play.google.com/store/apps/details?id=com.luvvix.onbuch&hl=fr}}%
      \hspace{8pt}\begin{minipage}[c]{0.5\linewidth}
        {\popdisplay\fontsize{11}{12.5}\selectfont\color{white}\MakeUppercase{Télécharge OnBuch}}\par\vspace{4pt}
        {\raggedright\popsemi\fontsize{8.4}{11}\selectfont\color{white}Gratuit \textbullet\ Google Play\\Tuteur IA, annales, concours, résultats\par}
      \end{minipage}
    \end{tcolorbox}
  \end{minipage}\hfill
  \begin{minipage}[t]{0.487\linewidth}
    \begin{tcolorbox}[enhanced,colback=poppurple,colframe=popink,boxrule=2.2pt,arc=10pt,
        drop shadow=popink,left=11pt,right=11pt,top=10pt,bottom=10pt,height=3.1cm,valign=center]
      \tikz[baseline=-0.7ex]\node[circle,fill=white,draw=popink,line width=1.3pt,minimum size=1.6cm,inner sep=0]{\popdisplay\fontsize{24}{24}\selectfont\color{poppurple}+};%
      \hspace{8pt}\begin{minipage}[c]{0.6\linewidth}
        {\popdisplay\fontsize{11}{12.5}\selectfont\color{white}\MakeUppercase{Passe à OnBuch+}}\par\vspace{4pt}
        {\raggedright\popsemi\fontsize{8.4}{11}\selectfont\color{white}Tous les cours signés, Tuteur IA illimité, fiches PDF — onglet OnBuch+ de l'app\par}
      \end{minipage}
    \end{tcolorbox}
  \end{minipage}
  \vspace{8pt}
  \popcontenu
  \vfill
  \signatureonbuch
  \restoregeometry
  \newpage
}

% Rangée « Ce cours contient » (page de garde)
\newcommand{\poptile}[3]{% #1 couleur #2 gros texte #3 légende
  \begin{tcolorbox}[enhanced,colback=#1,colframe=popink,boxrule=2pt,arc=9pt,shadow={2.5pt}{-2.5pt}{0pt}{popink},
      left=6pt,right=6pt,top=9pt,bottom=9pt,halign=center,valign=center,height=1.75cm,width=\linewidth,nobeforeafter]
    {\popdisplay\fontsize{12.5}{13}\selectfont\color{white}\MakeUppercase{#2}}\par\vspace{3pt}
    {\centering\popsemi\fontsize{7.8}{9.5}\selectfont\color{white}#3\par}
  \end{tcolorbox}}
\newcommand{\popcontenu}{%
  \par\noindent{\popxboldls\fontsize{7.5}{7.5}\selectfont\color{popmuted}CE COURS SIGNÉ CONTIENT}\par\vspace{7pt}
  \noindent\begin{minipage}[t]{0.235\linewidth}\poptile{popblue}{Cours}{complet, illustré, pas à pas}\end{minipage}\hfill
  \begin{minipage}[t]{0.235\linewidth}\poptile{poporange}{Méthodes}{et exercices résolus}\end{minipage}\hfill
  \begin{minipage}[t]{0.235\linewidth}\poptile{poppink}{Exercices}{type examen + corrigés}\end{minipage}\hfill
  \begin{minipage}[t]{0.235\linewidth}\poptile{popgreen}{Fiche bilan}{l'essentiel à retenir}\end{minipage}\par}

% Sommaire
\newtcolorbox{sommaire}{enhanced,colback=white,colframe=popink,boxrule=2.2pt,arc=10pt,
  drop shadow=popink,left=18pt,right=18pt,top=13pt,bottom=13pt,before skip=6pt,after skip=20pt,
  before upper={\poppill{Sommaire}{poppurple}{white}\par\vspace{9pt}\popsemi\fontsize{10.6}{14.5}\selectfont\color{popink}}}

% Signature de fin de cours — poussée tout en bas de la dernière page
\newcommand{\findecours}{%
  \par\vfill\nopagebreak
  \begin{center}\popsquiggle{poporange}\end{center}\vspace{4pt}
  \signatureonbuch
}
\AtBeginDocument{\def\llbracket{\mathopen{[\mkern-3.5mu[}}\def\rrbracket{\mathclose{]\mkern-3.5mu]}}} % intervalles d'entiers (glyphe ⟦ absent de la police)
% Glyphes absents de Fira Math : redéfinis à partir de glyphes disponibles
\AtBeginDocument{%
  \def\setminus{\mathbin{\backslash}}\let\smallsetminus\setminus
  \def\vdots{\mathord{\vbox{\baselineskip3.2pt\lineskiplimit0pt\kern4pt\hbox{.}\hbox{.}\hbox{.}}}}%
  \def\ddots{\mathinner{\mkern1mu\raise6.5pt\hbox{.}\mkern2mu\raise3.6pt\hbox{.}\mkern2mu\raise0.7pt\hbox{.}\mkern1mu}}%
  \def\triangle{\mathord{\scalebox{0.9}{$\symup{\Delta}$}}}%
  \def\nabla{\mathord{\raisebox{\depth}{\scalebox{1}[-1]{$\symup{\Delta}$}}}}%
  \def\checkmark{\popcheck{popdark}}%
  \def\star{\mathbin{\ast}}\def\bigstar{\mathord{\ast}}%
  \def\diamond{\mathbin{\scalebox{0.6}{$\Diamond$}}}%
}

%% FIGFIX-BEGIN v5 — correctifs globaux des figures (docs/onbuch-generation/tools/figfix)
\usepackage{adjustbox}\usepackage{etoolbox}\tcbuselibrary{raster}
% toute figure TikZ/pgfplots plus large que la ligne est réduite à la largeur disponible
\BeforeBeginEnvironment{tikzpicture}{\begin{adjustbox}{max width=\linewidth}}
\AfterEndEnvironment{tikzpicture}{\end{adjustbox}}
% légendes pgfplots lisibles par-dessus les courbes
\pgfplotsset{every axis/.append style={legend style={fill=white,fill opacity=0.92,text opacity=1,draw=black!15,inner sep=3pt}}}
% étiquettes d'arêtes (syntaxe edge["..."]) avec fond blanc
\tikzset{every edge quotes/.append style={fill=white,inner sep=1.2pt}}
%% FIGFIX-END
\begin{document}

\pagegarde{Nombres complexes : approche géométrique}{Mathématiques}{Tle C}{Configurations et transformations élémentaires du plan}{13}{16 h}{Cette leçon donne une vision géométrique des nombres complexes : affixes, modules et arguments deviennent des outils pour décrire des ensembles de points (cercles, disques, médiatrices, arcs de cercles) et pour transformer le plan. Elle fournit les techniques essentielles du Bac (formes trigonométrique et exponentielle, formule de Moivre, racines n-ièmes, lignes de niveau) et leur application à la géométrie plane.
\vspace{4pt}
\begin{multicols}{2}\small
\begin{itemize}
\item À la fin de cette leçon, je sais déterminer l'affixe d'un point, d'un vecteur et interpréter géométriquement module et argument
\item À la fin de cette leçon, je sais déterminer et représenter un ensemble de points défini par une condition sur le module (cercle, disque, médiatrice, cercle d'Apollonius)
\item À la fin de cette leçon, je sais passer de la forme algébrique à la forme trigonométrique ou exponentielle et réciproquement
\item À la fin de cette leçon, je sais utiliser les propriétés des arguments d'un produit, d'un quotient et d'une puissance
\item À la fin de cette leçon, je sais appliquer la formule de Moivre pour linéariser cosⁿx et sinⁿx ou exprimer cos(nx) et sin(nx)
\item À la fin de cette leçon, je sais déterminer et représenter les racines n-ièmes d'un nombre complexe
\end{itemize}\end{multicols}}

\begin{sommaire}
\begin{multicols}{2}
\begin{enumerate}
\item Plan complexe, affixes et interprétation géométrique
\item Ensembles de points définis par une condition sur le module
\item Argument, forme trigonométrique et forme exponentielle
\item Opérations sur les arguments et applications géométriques
\item Formule de Moivre, linéarisation et racines n-ièmes
\item Lignes de niveau et résolution de problèmes de géométrie
\item Activité d'intégration
\item Méthodes et exercices résolus
\item Exercices
\item Corrigés des exercices
\item Fiche bilan
\end{enumerate}
\end{multicols}
\end{sommaire}

\begin{objectifs}
\begin{itemize}
\item À la fin de cette leçon, je sais déterminer l'affixe d'un point, d'un vecteur et interpréter géométriquement module et argument
\item À la fin de cette leçon, je sais déterminer et représenter un ensemble de points défini par une condition sur le module (cercle, disque, médiatrice, cercle d'Apollonius)
\item À la fin de cette leçon, je sais passer de la forme algébrique à la forme trigonométrique ou exponentielle et réciproquement
\item À la fin de cette leçon, je sais utiliser les propriétés des arguments d'un produit, d'un quotient et d'une puissance
\item À la fin de cette leçon, je sais appliquer la formule de Moivre pour linéariser cosⁿx et sinⁿx ou exprimer cos(nx) et sin(nx)
\item À la fin de cette leçon, je sais déterminer et représenter les racines n-ièmes d'un nombre complexe
\item À la fin de cette leçon, je sais déterminer et représenter une ligne de niveau définie par une condition sur un argument
\item À la fin de cette leçon, je sais utiliser les nombres complexes pour résoudre un problème de géométrie plane (alignement, orthogonalité, angles, configurations)
\end{itemize}
\end{objectifs}

\begin{prerequis}
\begin{itemize}
\item Nombres complexes : forme algébrique, opérations, conjugué, module (début de l'année de Terminale)
\item Trigonométrie de Première : cercle trigonométrique, angles orientés, cosinus et sinus remarquables
\item Géométrie vectorielle : coordonnées d'un vecteur, norme, colinéarité, produit scalaire
\item Équations du second degré dans ℂ et discriminant
\item Ensembles de points classiques : cercle, médiatrice, bissectrice (Seconde/Première)
\end{itemize}
\end{prerequis}

\newpage

\coursec{Plan complexe, affixes et interprétation géométrique}

\textbf{Situation d'accroche.} Un technicien de la SONATREL doit localiser sur une carte du réseau électrique de Yaoundé tous les points situés à exactement 5 km d'un transformateur, puis ceux plus proches du poste d'Essos que du poste de Nlongkak dans un rapport donné. Un ingénieur télécom d'Orange Cameroun veut déterminer la zone couverte par une antenne et l'angle sous lequel un village voit deux relais. Un charpentier à Douala doit partager un plateau circulaire en 7 parts rigoureusement égales. Tous ces problèmes de distances, d'angles et de rotations semblent très différents. Pourtant, ils se résolvent tous élégamment en associant à chaque point du plan un unique nombre, appelé \cle{nombre complexe}. C'est toute la puissance de l'approche géométrique des nombres complexes : traduire la géométrie en calcul algébrique, calculer, puis revenir à la géométrie.

\textbf{Problématique.} Comment un nombre complexe peut-il représenter un point ou un vecteur du plan, et comment les opérations sur les nombres complexes traduisent-elles les distances, les milieux, l'alignement et l'orthogonalité ?

\courssub{Rappels : l'ensemble $\mathbb{C}$ et la forme algébrique}

\textbf{Intuition.} Dans $\mathbb{R}$, l'équation $x^2 = -1$ n'a pas de solution : un carré est toujours positif. Les mathématiciens ont alors fait un pari audacieux : \emph{inventer} un nombre dont le carré vaut $-1$, noté $i$, et construire autour de lui un ensemble plus vaste, stable par addition et multiplication. Ce pari, couronné de succès au XIX\textsuperscript{e} siècle, a donné naissance à l'ensemble $\mathbb{C}$, aujourd'hui indispensable en électricité, en télécommunications et en mécanique des fluides.

\begin{definition}[L'ensemble des nombres complexes]
Il existe un ensemble, noté $\mathbb{C}$, qui contient $\mathbb{R}$, muni d'une addition et d'une multiplication prolongeant celles de $\mathbb{R}$, et d'un élément noté $i$ tel que
\[ i^2 = -1. \]
Tout nombre complexe $z$ s'écrit de manière \emph{unique} sous la forme
\[ z = x + iy, \quad \text{avec } (x, y) \in \mathbb{R}^2, \]
appelée \cle{forme algébrique} de $z$. Le réel $x$ est la \cle{partie réelle} de $z$, notée $\operatorname{Re}(z)$ ; le réel $y$ est la \cle{partie imaginaire}, notée $\operatorname{Im}(z)$.
\end{definition}

\begin{attention}[La partie imaginaire est un réel]
Si $z = 3 - 2i$, alors $\operatorname{Im}(z) = -2$ et \textbf{non pas} $-2i$. La partie imaginaire d'un complexe est toujours un nombre \emph{réel}. C'est une erreur classique au Baccalauréat.
\end{attention}

\begin{definition}[Conjugué et module]
Soit $z = x + iy$ un nombre complexe.
\begin{itemize}
\item Le \cle{conjugué} de $z$ est le nombre $\bar{z} = x - iy$.
\item Le \cle{module} de $z$ est le réel positif $|z| = \sqrt{x^2 + y^2}$.
\end{itemize}
On rappelle les propriétés fondamentales : $z\bar{z} = |z|^2$, $\overline{z + z'} = \bar{z} + \bar{z'}$, $\overline{zz'} = \bar{z}\,\bar{z'}$, et $|zz'| = |z|\,|z'|$.
\end{definition}

\begin{exemplebox}[Vérification de l'identité $z\bar{z} = |z|^2$]
Prenons $z = 3 + 2i$. Alors $\bar{z} = 3 - 2i$ et
\begin{align*}
z\bar{z} &= (3 + 2i)(3 - 2i) = 9 - (2i)^2 = 9 - 4i^2 = 9 + 4 = 13, \\
|z|^2 &= 3^2 + 2^2 = 13.
\end{align*}
Les deux coïncident, comme annoncé.
\end{exemplebox}

\courssub{Le plan complexe : image d'un complexe, affixe d'un point}

\textbf{Intuition.} Un nombre complexe $z = x + iy$ est entièrement déterminé par le couple de réels $(x, y)$. Or un couple de réels, c'est exactement ce qui repère un point dans un plan muni d'un repère. L'idée géniale consiste donc à \emph{identifier} le nombre $z$ au point de coordonnées $(x, y)$ : le plan devient une « carte » de l'ensemble $\mathbb{C}$, exactement comme la carte de Yaoundé représente la ville réelle.

\begin{definition}[Plan complexe, image, affixe]
Le plan est muni d'un repère orthonormé direct $(O; \vec{u}, \vec{v})$, appelé alors \cle{plan complexe}.
\begin{itemize}
\item À tout nombre complexe $z = x + iy$, on associe le point $M(x, y)$, appelé \cle{image} de $z$. On note $M(z)$.
\item Réciproquement, à tout point $M(x, y)$, on associe le nombre complexe $z = x + iy$, appelé \cle{affixe} de $M$. On note $z_M$ ou $\operatorname{aff}(M)$.
\item L'axe $(O; \vec{u})$ est l'\cle{axe des réels} (affixes réelles) ; l'axe $(O; \vec{v})$ est l'\cle{axe des imaginaires purs} (affixes de la forme $iy$, $y \in \mathbb{R}$).
\end{itemize}
\end{definition}

\begin{aretenir}[Le dictionnaire points $\leftrightarrow$ complexes]
Dans le plan complexe, on passe librement du point $M$ à son affixe $z_M$ et inversement. Dire « le point $M(3 + 2i)$ » signifie : le point de coordonnées $(3, 2)$. Toute la leçon repose sur ce va-et-vient.
\end{aretenir}

\begin{popfigure}
\begin{center}
\begin{tikzpicture}[scale=1.1, >=Stealth]
\draw[->, popmuted] (-1.2,0) -- (4.5,0) node[below left] {axe des réels};
\draw[->, popmuted] (0,-2.8) -- (0,3.2) node[below left] {axe des imaginaires};
\draw[popline, dashed] (3,0) -- (3,2) -- (0,2);
\draw[popline, dashed] (3,0) -- (3,-2) -- (0,-2);
\draw[very thick, poporange, ->] (0,0) -- (3,2) node[midway, above left, popdark] {$OM = |z|$};
\fill[popblue] (3,2) circle (2.2pt) node[above right, popdark] {$M(z = 3 + 2i)$};
\fill[poppink] (3,-2) circle (2.2pt) node[below right, popdark] {$M'(\bar{z} = 3 - 2i)$};
\fill[popdark] (0,0) circle (2pt) node[below left] {$O$};
\draw[<->, popgreen, thick] (3.35,2) -- (3.35,-2) node[midway, right, popdark] {symétrie / axe des réels};
\foreach \x in {1,2,3} \draw (\x,0.08) -- (\x,-0.08) node[below, popmuted] {\x};
\foreach \y in {1,2} \draw (0.08,\y) -- (-0.08,\y) node[left, popmuted] {\y};
\foreach \y in {-1,-2} \draw (0.08,\y) -- (-0.08,\y) node[left, popmuted] {\y};
\end{tikzpicture}
\end{center}
\legende{Le point $M$ d'affixe $z = 3 + 2i$, son conjugué $M'$ d'affixe $\bar{z} = 3 - 2i$, symétrique de $M$ par rapport à l'axe des réels, et le module $|z| = OM = \sqrt{13}$.}
\end{popfigure}

\courssub{Affixe d'un vecteur, affixe du milieu}

\textbf{Intuition.} Un vecteur $\overrightarrow{AB}$ est caractérisé par son déplacement : « combien d'unités horizontalement, combien verticalement ». Si $A$ a pour coordonnées $(x_A, y_A)$ et $B$ les coordonnées $(x_B, y_B)$, ce déplacement vaut $(x_B - x_A, y_B - y_A)$. En langage complexe, cela s'écrit de façon remarquablement simple : $z_B - z_A$.

\begin{definition}[Affixe d'un vecteur]
Soient $A$ et $B$ deux points d'affixes respectives $z_A$ et $z_B$. L'\cle{affixe} du vecteur $\overrightarrow{AB}$ est
\[ z_{\overrightarrow{AB}} = z_B - z_A. \]
Plus généralement, un vecteur $\vec{w}$ de coordonnées $(a, b)$ a pour affixe $z_{\vec{w}} = a + ib$.
\end{definition}

\begin{propriete}[Distance et affixe du milieu]
Soient $A(z_A)$ et $B(z_B)$ deux points du plan complexe.
\begin{enumerate}
\item La distance $AB$ est le module de la différence des affixes :
\[ AB = |z_B - z_A|. \]
En particulier, pour tout point $M(z)$ : $OM = |z|$.
\item Le milieu $I$ du segment $[AB]$ a pour affixe
\[ z_I = \frac{z_A + z_B}{2}. \]
\end{enumerate}
\end{propriete}

\begin{experience}[Démonstration guidée (exigible)]
On se propose de démontrer ces deux résultats à partir des coordonnées.
\begin{enumerate}
\item Écris $z_A = x_A + iy_A$ et $z_B = x_B + iy_B$. Calcule $z_B - z_A$ sous forme algébrique : tu dois obtenir $(x_B - x_A) + i(y_B - y_A)$.
\item Applique la définition du module : $|z_B - z_A| = \sqrt{(x_B - x_A)^2 + (y_B - y_A)^2}$. Reconnais la formule de la distance entre deux points en repère orthonormé : c'est exactement $AB$. Conclus.
\item Pour le milieu : les coordonnées de $I$ sont $\left(\dfrac{x_A + x_B}{2}, \dfrac{y_A + y_B}{2}\right)$. Son affixe est donc
\[ z_I = \frac{x_A + x_B}{2} + i\,\frac{y_A + y_B}{2} = \frac{(x_A + iy_A) + (x_B + iy_B)}{2} = \frac{z_A + z_B}{2}. \]
\end{enumerate}
Le point clé est que les formules complexes \emph{condensent} en une seule égalité les deux formules sur les coordonnées.
\end{experience}

\begin{attention}[Sens de la différence]
L'affixe du vecteur $\overrightarrow{AB}$ est $z_B - z_A$ (« arrivée moins départ »), et \textbf{pas} $z_A - z_B$. En revanche, la distance ne dépend pas du sens : $AB = |z_B - z_A| = |z_A - z_B|$, car $|-w| = |w|$. Astuce de mémorisation : $\overrightarrow{AB}$ se lit « de $A$ vers $B$ », donc $z_B - z_A$.
\end{attention}

\begin{popfigure}
\begin{center}
\begin{tikzpicture}[scale=1.0, >=Stealth]
\draw[->, popmuted] (-0.8,0) -- (5,0);
\draw[->, popmuted] (0,-0.8) -- (0,4);
\coordinate (O) at (0,0);
\coordinate (A) at (1,1);
\coordinate (B) at (4,2);
\coordinate (C) at (3,1); % Corrigé : z_B - z_A = 3+i
\draw[very thick, popblue, ->] (A) -- (B) node[midway, below right, popdark] {$\overrightarrow{AB}$};
\draw[very thick, poporange, ->] (O) -- (C) node[midway, above left, popdark] {$\overrightarrow{OC}$};
\draw[dashed, popline] (A) -- (O);
\draw[dashed, popline] (B) -- (C);
\fill[popdark] (A) circle (2pt) node[below left] {$A(z_A)$};
\fill[popdark] (B) circle (2pt) node[right] {$B(z_B)$};
\fill[popdark] (C) circle (2pt) node[above right] {$C(z_B - z_A)$};
\fill[popdark] (O) circle (2pt) node[below left] {$O$};
\node[popmuted] at (2.5,0.35) {$OABC$ parallélogramme};
\end{tikzpicture}
\end{center}
\legende{Le vecteur $\overrightarrow{AB}$ est égal au vecteur $\overrightarrow{OC}$ où $C$ a pour affixe $z_B - z_A$ : le quadrilatère $OABC$ est un parallélogramme. L'affixe d'un vecteur se lit donc aussi comme l'affixe du point extrémité quand on part de l'origine.}
\end{popfigure}

\begin{propriete}[Interprétation géométrique du conjugué]
Soit $M$ le point d'affixe $z$ et $M'$ le point d'affixe $\bar{z}$. Alors $M'$ est le symétrique de $M$ par rapport à l'axe des réels. En particulier, $z$ est réel si et seulement si $M$ appartient à l'axe des réels, c'est-à-dire si et seulement si $\bar{z} = z$.
\end{propriete}

En effet, si $z = x + iy$, alors $M(x, y)$ et $M'(x, -y)$ : même abscisse, ordonnées opposées, ce qui caractérise la symétrie orthogonale d'axe $(O; \vec{u})$.

\courssub{Colinéarité, alignement et orthogonalité}

\textbf{Intuition.} Dire que trois points $A$, $B$, $C$ sont alignés revient à dire que le vecteur $\overrightarrow{AC}$ est un « multiple » du vecteur $\overrightarrow{AB}$ par un nombre \emph{réel}. En langage complexe, multiplier par un réel, c'est rester sur la même direction. Le quotient des affixes des deux vecteurs doit donc être un nombre réel. De même, une perpendicularité correspond à un « quart de tour », que le nombre $i$ réalise : le quotient doit être un imaginaire pur.

\begin{propriete}[Alignement et orthogonalité]
Soient $A$, $B$, $C$ trois points d'affixes $z_A$, $z_B$, $z_C$, avec $A \neq B$.
\begin{enumerate}
\item Les points $A$, $B$, $C$ sont alignés si et seulement si
\[ \frac{z_C - z_A}{z_B - z_A} \in \mathbb{R}. \]
\item Les droites $(AB)$ et $(AC)$ sont perpendiculaires si et seulement si
\[ \frac{z_C - z_A}{z_B - z_A} \in i\mathbb{R} \quad \text{(imaginaire pur)}. \]
\end{enumerate}
Le point 2 sera démontré rigoureusement avec les arguments dans la section 4 ; on l'admet pour l'instant.
\end{propriete}

\begin{experience}[Démonstration guidée du critère d'alignement]
\begin{enumerate}
\item Rappelle le critère de colinéarité vu en Première : $A$, $B$, $C$ alignés $\Leftrightarrow$ il existe un réel $k$ tel que $\overrightarrow{AC} = k\,\overrightarrow{AB}$.
\item Traduis cette égalité vectorielle en affixes : $z_C - z_A = k(z_B - z_A)$.
\item Comme $A \neq B$, on a $z_B - z_A \neq 0$ ; divise : $\dfrac{z_C - z_A}{z_B - z_A} = k \in \mathbb{R}$. Réciproquement, si ce quotient est un réel $k$, alors $\overrightarrow{AC} = k\overrightarrow{AB}$, donc les points sont alignés. Conclus.
\end{enumerate}
\end{experience}

\begin{methode}[Tester l'alignement de trois points]
Pour montrer que $A(z_A)$, $B(z_B)$, $C(z_C)$ sont alignés :
\begin{enumerate}
\item Calculer $z_B - z_A$ et $z_C - z_A$.
\item Former le quotient $q = \dfrac{z_C - z_A}{z_B - z_A}$ et l'écrire sous forme algébrique.
\item Conclure : si $\operatorname{Im}(q) = 0$, les points sont alignés ; sinon, ils ne le sont pas. Si de plus $\operatorname{Re}(q) = 0$ et $q \neq 0$, alors $(AB) \perp (AC)$.
\end{enumerate}
\end{methode}

\begin{methode}[Montrer qu'un quadrilatère est un parallélogramme]
Pour montrer que $ABCD$ (sommets pris dans l'ordre) est un parallélogramme :
\begin{enumerate}
\item Calculer les affixes des milieux des diagonales : milieu de $[AC]$ d'affixe $\dfrac{z_A + z_C}{2}$ et milieu de $[BD]$ d'affixe $\dfrac{z_B + z_D}{2}$.
\item Comparer : $ABCD$ est un parallélogramme si et seulement si les diagonales ont le même milieu, c'est-à-dire
\[ z_A + z_C = z_B + z_D. \]
\item Variante équivalente : vérifier que $\overrightarrow{AB} = \overrightarrow{DC}$, c'est-à-dire $z_B - z_A = z_C - z_D$.
\end{enumerate}
\end{methode}

\begin{exoresolu}[Distances, milieu et alignement]
On donne les points $A(1 + i)$, $B(4 - 2i)$ et $C(-1 + 3i)$.
\begin{enumerate}
\item Calculer la distance $AB$.
\item Déterminer l'affixe du milieu $I$ de $[AC]$.
\item Les points $A$, $B$ et $C$ sont-ils alignés ?
\end{enumerate}
\tcblower
\textbf{1. Distance $AB$.} On calcule l'affixe de $\overrightarrow{AB}$ :
\[ z_B - z_A = (4 - 2i) - (1 + i) = 3 - 3i. \]
Donc
\[ AB = |z_B - z_A| = |3 - 3i| = \sqrt{3^2 + (-3)^2} = \sqrt{18} = 3\sqrt{2}. \]

\textbf{2. Milieu de $[AC]$.}
\[ z_I = \frac{z_A + z_C}{2} = \frac{(1 + i) + (-1 + 3i)}{2} = \frac{4i}{2} = 2i. \]
Le milieu $I$ est le point de coordonnées $(0, 2)$, situé sur l'axe des imaginaires purs.

\textbf{3. Alignement.} On a $z_B - z_A = 3 - 3i$ et
\[ z_C - z_A = (-1 + 3i) - (1 + i) = -2 + 2i. \]
Formons le quotient :
\begin{align*}
q &= \frac{z_C - z_A}{z_B - z_A} = \frac{-2 + 2i}{3 - 3i} = \frac{2(-1 + i)}{3(1 - i)}. \\
\text{Or } -1 + i &= -(1 - i), \text{ donc} \\
q &= \frac{-2(1 - i)}{3(1 - i)} = -\frac{2}{3} \in \mathbb{R}.
\end{align*}
Le quotient est réel : les points $A$, $B$ et $C$ sont \textbf{alignés}. On a même montré au passage que $\overrightarrow{AC} = -\dfrac{2}{3}\overrightarrow{AB}$ : le point $C$ est du côté opposé à $B$ par rapport à $A$.
\end{exoresolu}

\begin{exemplebox}[Un parallélogramme par les affixes]
Soient $A(1 + i)$, $B(4 - 2i)$ et $D(2 + 4i)$. Cherchons l'affixe du point $E$ tel que $ABED$ soit un parallélogramme. La condition $z_A + z_E = z_B + z_D$ donne
\[ z_E = z_B + z_D - z_A = (4 - 2i) + (2 + 4i) - (1 + i) = 5 + i. \]
Vérification : $z_B - z_A = 3 - 3i$ et $z_E - z_D = (5 + i) - (2 + 4i) = 3 - 3i$. Les vecteurs $\overrightarrow{AB}$ et $\overrightarrow{DE}$ ont la même affixe, donc sont égaux : $ABED$ est bien un parallélogramme.
\end{exemplebox}

\begin{attention}[Pièges de rédaction au Baccalauréat]
\begin{itemize}
\item Ne jamais écrire $AB = z_B - z_A$ : une distance est un \emph{réel positif}, c'est $AB = |z_B - z_A|$ avec le module.
\item Pour appliquer le critère d'alignement, vérifier d'abord que $A \neq B$ (sinon le quotient n'a pas de sens).
\item Le quotient $\dfrac{z_C - z_A}{z_B - z_A}$ doit être donné sous \emph{forme algébrique} avant de conclure : un quotient non simplifié ne permet pas de décider s'il est réel.
\end{itemize}
\end{attention}

\begin{savaistu}[Des nombres « impossibles » devenus indispensables]
Les nombres complexes sont nés en Italie au XVI\textsuperscript{e} siècle, lorsque Cardan et Bombelli résolvaient les équations du troisième degré : ils manipulaient des racines carrées de nombres négatifs qu'ils jugeaient eux-mêmes « sophistiqués » et « inutiles ». Trois siècles plus tard, Gauss et Argand leur donnent une existence géométrique solide en les représentant dans le plan — exactement ce que tu viens d'apprendre. Aujourd'hui, les ingénieurs de la SONATREL décrivent les courants alternatifs par des nombres complexes (l'« impédance complexe »), et les ingénieurs d'Orange ou de MTN les utilisent pour modéliser les ondes : les nombres « imaginaires » font fonctionner des réseaux bien réels.
\end{savaistu}

\begin{exercice}[À toi de jouer]
Dans le plan complexe, on donne $A(-2 + i)$, $B(3 - i)$ et $C(1 + 5i)$.
\begin{enumerate}
\item Placer ces points dans un repère orthonormé.
\item Calculer $AB$, $AC$ et $BC$.
\item Déterminer l'affixe du point $D$ tel que $ABCD$ soit un parallélogramme.
\item Les droites $(AB)$ et $(AC)$ sont-elles perpendiculaires ? Justifier par un calcul de quotient.
\end{enumerate}
\end{exercice}

\begin{corrige}[Exercice]
\textbf{1.} $A(-2, 1)$, $B(3, -1)$, $C(1, 5)$ : placement direct.

\textbf{2.} $z_B - z_A = 5 - 2i$ donc $AB = \sqrt{25 + 4} = \sqrt{29}$. $z_C - z_A = 3 + 4i$ donc $AC = \sqrt{9 + 16} = 5$. $z_C - z_B = -2 + 6i$ donc $BC = \sqrt{4 + 36} = \sqrt{40} = 2\sqrt{10}$.

\textbf{3.} $ABCD$ parallélogramme $\Leftrightarrow z_A + z_C = z_B + z_D$, donc
\[ z_D = z_A + z_C - z_B = (-2 + i) + (1 + 5i) - (3 - i) = -4 + 7i. \]

\textbf{4.} On calcule
\begin{align*}
q &= \frac{z_C - z_A}{z_B - z_A} = \frac{3 + 4i}{5 - 2i} = \frac{(3 + 4i)(5 + 2i)}{(5 - 2i)(5 + 2i)} \\
  &= \frac{15 + 6i + 20i + 8i^2}{29} = \frac{7 + 26i}{29}.
\end{align*}
Ce quotient n'est ni réel ($\operatorname{Im}(q) = \tfrac{26}{29} \neq 0$) ni imaginaire pur ($\operatorname{Re}(q) = \tfrac{7}{29} \neq 0$) : les droites $(AB)$ et $(AC)$ ne sont \textbf{pas} perpendiculaires.
\end{corrige}

\begin{aretenir}[L'essentiel de la section]
\begin{itemize}
\item Plan complexe $(O; \vec{u}, \vec{v})$ orthonormé direct : le point $M(x, y)$ a pour affixe $z = x + iy$.
\item Affixe d'un vecteur : $z_{\overrightarrow{AB}} = z_B - z_A$ ; distance : $AB = |z_B - z_A|$ ; milieu de $[AB]$ : $z_I = \dfrac{z_A + z_B}{2}$.
\item Le conjugué $\bar{z}$ correspond à la symétrie par rapport à l'axe des réels.
\item $A$, $B$, $C$ alignés $\Leftrightarrow \dfrac{z_C - z_A}{z_B - z_A} \in \mathbb{R}$ ; $(AB) \perp (AC) \Leftrightarrow$ ce quotient est imaginaire pur.
\item $ABCD$ parallélogramme $\Leftrightarrow z_A + z_C = z_B + z_D$.
\end{itemize}
\end{aretenir}

\coursec{Ensembles de points définis par une condition sur le module}

Dans la section précédente, nous avons appris à plonger le plan géométrique dans le monde des nombres complexes : à tout point $M$ du plan correspond son \cle{affixe} $z$, à tout vecteur $\vec{u}$ correspond un nombre complexe, et surtout, la distance entre deux points $A(a)$ et $M(z)$ s'écrit simplement $AM = |z - a|$. Cette dernière identité est un trésor : elle signifie que toute condition géométrique portant sur des distances peut se traduire par une équation ou une inéquation portant sur des modules. Réciproquement, une équation en modules décrit un lieu géométrique. C'est exactement l'objet de cette section : lire une condition sur les modules comme on lit une phrase géométrique, puis représenter l'ensemble de points correspondant. Au Baccalauréat, ces questions sont presque systématiques dans l'exercice de nombres complexes : « déterminer et représenter l'ensemble des points $M$ tels que... ». Maîtrisons donc ce réflexe fondamental.

\courssub{Le cercle : ensemble des points tels que $|z - a| = \alpha$}

Commençons par une situation concrète. Un opérateur téléphonique (MTN ou Orange) installe une antenne relais en un point $\Omega$ et garantit une couverture optimale dans un rayon de $2$ km. Où se trouvent les clients situés exactement à la limite de la zone ? Sur un cercle de centre $\Omega$ et de rayon $2$. Traduisons cela en langage complexe : si l'antenne est placée au point $\Omega$ d'affixe $a$, et si $M$ a pour affixe $z$, alors $M$ est à la limite de la zone si et seulement si $\Omega M = 2$, c'est-à-dire $|z - a| = 2$.

\begin{propriete}[Cercle défini par un module]
Soit $\Omega$ le point d'affixe $a$ et $\alpha$ un réel strictement positif. L'ensemble des points $M$ d'affixe $z$ vérifiant
\[ |z - a| = \alpha \]
est le \cle{cercle} de centre $\Omega$ et de rayon $\alpha$.
\end{propriete}

\begin{experience}[Démonstration (immédiate mais à savoir rédiger)]
Soit $M$ un point d'affixe $z$ et $\Omega$ le point d'affixe $a$. D'après l'interprétation géométrique du module vue à la section précédente, la distance entre $\Omega$ et $M$ vaut :
\[ \Omega M = |z_M - z_\Omega| = |z - a|. \]
Par conséquent :
\[ |z - a| = \alpha \iff \Omega M = \alpha \iff M \text{ appartient au cercle de centre } \Omega \text{ et de rayon } \alpha. \]
C'est une simple chaîne d'équivalences, mais elle doit apparaître explicitement dans ta copie : le jury attend la phrase « $\Omega M = |z-a|$ ».
\end{experience}

\begin{exemplebox}[Lecture directe]
L'ensemble des points $M(z)$ tels que $|z - 3 + 2i| = 5$ est le cercle de centre $\Omega(3 - 2i)$ et de rayon $5$. Attention à la réécriture : $|z - 3 + 2i| = |z - (3 - 2i)|$, donc l'affixe du centre est $a = 3 - 2i$ et non $-3 + 2i$.
\end{exemplebox}

\courssub{Le disque : ensembles définis par $|z - a| \leq \alpha$ et variantes}

Revenons à notre antenne : les clients \emph{dans} la zone de couverture vérifient $\Omega M \leq 2$, c'est-à-dire $|z - a| \leq 2$. Géométriquement, ce sont les points du cercle et de son intérieur : c'est un \cle{disque}.

\begin{propriete}[Disques et extérieurs]
Soit $\Omega$ le point d'affixe $a$ et $\alpha > 0$. Pour $M$ d'affixe $z$ :
\begin{itemize}
\item $|z - a| \leq \alpha$ définit le \cle{disque fermé} de centre $\Omega$ et de rayon $\alpha$ (le cercle est inclus) ;
\item $|z - a| < \alpha$ définit le \cle{disque ouvert} (intérieur strict, cercle exclu) ;
\item $|z - a| > \alpha$ définit l'\cle{extérieur} du cercle (cercle exclu) ;
\item $|z - a| \geq \alpha$ définit l'extérieur du cercle, cercle inclus.
\end{itemize}
\end{propriete}

La démonstration est identique à la précédente : $|z - a| \leq \alpha \iff \Omega M \leq \alpha$, ce qui signifie que $M$ est à une distance de $\Omega$ inférieure ou égale à $\alpha$.

\begin{attention}[Frontière incluse ou exclue : le détail qui coûte des points]
Lors du tracé, la convention est la suivante :
\begin{itemize}
\item inégalité \textbf{large} ($\leq$ ou $\geq$) : le cercle frontière est tracé en \textbf{trait plein} et fait partie de l'ensemble ;
\item inégalité \textbf{stricte} ($<$ ou $>$) : le cercle frontière est tracé en \textbf{pointillés} et n'en fait pas partie.
\end{itemize}
Précise toujours dans ta rédaction si la frontière est incluse ou non.
\end{attention}

\begin{exoresolu}[Déterminer et représenter un disque]
\textbf{Énoncé.} Le plan est muni d'un repère orthonormé direct $(O\,;\vec{u},\vec{v})$. Déterminer et représenter l'ensemble $(E)$ des points $M$ d'affixe $z$ tels que $|z - 1 + 2i| \leq 3$.
\tcblower
\textbf{Solution.} On réécrit la condition sous la forme $|z - a| \leq \alpha$ :
\[ |z - 1 + 2i| \leq 3 \iff |z - (1 - 2i)| \leq 3. \]
Soit $\Omega$ le point d'affixe $a = 1 - 2i$. Alors $|z - a| = \Omega M$, donc :
\[ |z - (1 - 2i)| \leq 3 \iff \Omega M \leq 3. \]
L'ensemble $(E)$ est donc le \textbf{disque fermé} de centre $\Omega(1\,;-2)$ et de rayon $3$, frontière incluse. Pour le représenter : on place $\Omega$ de coordonnées $(1\,;-2)$, on trace le cercle de centre $\Omega$ et de rayon $3$ en trait plein, puis on hachure l'intérieur du disque.
\end{exoresolu}

\begin{popfigure}
\begin{center}
\begin{tikzpicture}[scale=0.75]
% Grille et axes
\draw[popline, very thin] (-2.5,-5.5) grid (4.5,1.5);
\draw[-{Stealth}, popdark] (-2.5,0) -- (4.5,0) node[below left] {$x$};
\draw[-{Stealth}, popdark] (0,-5.5) -- (0,1.5) node[below left] {$y$};
\draw[popdark] (0,0) node[below left] {$O$};
% Disque hachuré : centre (1, -2), rayon 3
\fill[pattern=north east lines, pattern color=popblue, opacity=0.55] (1,-2) circle (3);
\draw[popblue, very thick] (1,-2) circle (3);
% Centre
\fill[popink] (1,-2) circle (0.07) node[above right, popdark] {$\Omega(1-2i)$};
% Rayon
\draw[poporange, thick, -{Stealth}] (1,-2) -- (4,-2) node[midway, above, popdark] {$3$};
\fill[poporange] (4,-2) circle (0.06);
\end{tikzpicture}
\end{center}
\legende{Le cercle $|z-(1-2i)| = 3$ (trait plein) et le disque fermé $|z-(1-2i)| \leq 3$ (hachuré), de centre $\Omega(1-2i)$ et de rayon $3$.}
\end{popfigure}

\courssub{La médiatrice : ensemble des points tels que $|z - a| = |z - b|$}

Changeons de registre. Deux villes $A$ et $B$ veulent construire un marché commun situé \emph{à égale distance} des deux. Où peut-on le placer ? La réponse est connue depuis le collège : sur la médiatrice du segment $[AB]$. En langage complexe, si $A$ a pour affixe $a$, $B$ pour affixe $b$ et $M$ pour affixe $z$, la condition « $M$ équidistant de $A$ et $B$ » s'écrit $MA = MB$, c'est-à-dire $|z - a| = |z - b|$.

\begin{propriete}[Médiatrice définie par une égalité de modules]
Soit $A$ et $B$ deux points distincts d'affixes respectives $a$ et $b$. L'ensemble des points $M$ d'affixe $z$ vérifiant
\[ |z - a| = |z - b| \]
est la \cle{médiatrice} du segment $[AB]$.
\end{propriete}

\begin{experience}[Démonstration (exigible)]
Soit $M$ d'affixe $z$. On a $MA = |z - a|$ et $MB = |z - b|$. Donc :
\[ |z - a| = |z - b| \iff MA = MB \iff M \text{ est équidistant de } A \text{ et } B \iff M \text{ appartient à la médiatrice de } [AB]. \]
La dernière équivalence est la caractérisation classique de la médiatrice vue au collège.
\end{experience}

\begin{exoresolu}[Un quotient de module égal à 1]
\textbf{Énoncé.} Déterminer l'ensemble $(\Delta)$ des points $M$ d'affixe $z$ tels que $\left|\dfrac{z - i}{z + 2}\right| = 1$.
\tcblower
\textbf{Solution.} La condition exige $z \neq -2$ (le dénominateur ne doit pas s'annuler). On utilise la propriété du module d'un quotient : pour $z \neq -2$,
\[ \left|\frac{z - i}{z + 2}\right| = \frac{|z - i|}{|z + 2|}. \]
Soit $A$ le point d'affixe $i$ et $B$ le point d'affixe $-2$. Alors $|z - i| = MA$ et $|z + 2| = |z - (-2)| = MB$. Donc :
\[ \left|\frac{z - i}{z + 2}\right| = 1 \iff \frac{MA}{MB} = 1 \iff MA = MB. \]
L'ensemble $(\Delta)$ est donc la \textbf{médiatrice du segment $[AB]$}, avec $A(0\,;1)$ et $B(-2\,;0)$. Le point $B$ ne satisfaisant pas l'équation (car $BA \neq 0 = BB$), il n'appartient pas à la médiatrice ; la restriction $z \neq -2$ est donc automatiquement vérifiée par tous les points de la médiatrice. On trace la droite perpendiculaire à $[AB]$ passant par le milieu $I(-1\,;\tfrac{1}{2})$ de $[AB]$.
\end{exoresolu}

\begin{aretenir}[Le cas du quotient égal à 1]
Pour $A(a)$, $B(b)$ et $M(z)$ avec $z \neq b$ :
\[ \left|\frac{z - a}{z - b}\right| = 1 \iff MA = MB \iff M \in \text{médiatrice de } [AB]. \]
Pense toujours à écrire $\left|\dfrac{z-a}{z-b}\right| = \dfrac{|z-a|}{|z-b|}$ : c'est la clé qui transforme un quotient de modules en rapport de distances.
\end{aretenir}

\begin{popfigure}
\begin{center}
\begin{tikzpicture}[scale=1.0]
\draw[popline, very thin] (-3.5,-2.5) grid (1.5,3.5);
\draw[-{Stealth}, popdark] (-3.5,0) -- (1.5,0) node[below left] {$x$};
\draw[-{Stealth}, popdark] (0,-2.5) -- (0,3.5) node[below left] {$y$};
\draw[popdark] (0,0) node[below left] {$O$};
% Points A et B : A(0,1) affixe i, B(-2,0) affixe -2
\fill[popink] (0,1) circle (0.08) node[above right, popdark] {$A(i)$};
\fill[popink] (-2,0) circle (0.08) node[below left, popdark] {$B(-2)$};
% Segment [AB] en pointillés
\draw[popmuted, dashed, thick] (0,1) -- (-2,0);
% Milieu I(-1, 0.5)
\fill[poporange] (-1,0.5) circle (0.06) node[right=2pt, popdark, fill=white, inner sep=1pt] {$I$};
% Médiatrice : perpendiculaire à (AB) passant par I(-1, 0.5)
% Vecteur AB = (-2, -1). Vecteur directeur médiatrice = (1, -2) ou (-1, 2).
% Points pour tracé : I + t*(1, -2). t=-1.5 -> (-2.5, 3.5). t=1.5 -> (0.5, -2.5).
\draw[popblue, very thick] (-2.5,3.5) -- (0.5,-2.5);
\draw[popblue] (0.5,-2.5) node[below right, popdark] {$(\Delta)$};
% Codage angle droit en I
\draw[poporange] (-1,0.5) ++(0.15,0.075) -- ++(0.075,-0.15) -- ++(-0.15,-0.075);
\end{tikzpicture}
\end{center}
\legende{La médiatrice $(\Delta)$ de $[AB]$, ensemble des points $M(z)$ tels que $|z - i| = |z + 2|$ : tous les points de $(\Delta)$ sont équidistants de $A(i)$ et de $B(-2)$.}
\end{popfigure}

\courssub{Le cercle d'Apollonius : ensemble des points tels que $\left|\dfrac{z - a}{z - b}\right| = \alpha$, $\alpha \neq 1$}

Que se passe-t-il si le rapport des distances n'est pas égal à $1$ ? Cherchons les points $M$ tels que $MA = 2\,MB$ : $M$ est deux fois plus loin de $A$ que de $B$. Intuitivement, $M$ est « attiré » vers $B$ ; le lieu n'est plus une droite. Le résultat, connu depuis l'Antiquité grecque, est remarquable : c'est un cercle, appelé \cle{cercle d'Apollonius}.

\begin{propriete}[Cercle d'Apollonius]
Soit $A$ et $B$ deux points distincts d'affixes $a$ et $b$, et $\alpha$ un réel strictement positif, $\alpha \neq 1$. L'ensemble des points $M$ d'affixe $z$ vérifiant
\[ \left|\frac{z - a}{z - b}\right| = \alpha \quad (\text{avec } z \neq b) \]
est un \cle{cercle}, appelé cercle d'Apollonius du segment $[AB]$ pour le rapport $\alpha$. Ses intersections avec la droite $(AB)$ sont les deux points $M$ de cette droite vérifiant $MA = \alpha\, MB$ ; elles sont diamétralement opposées sur le cercle.
\end{propriete}

\begin{experience}[Démonstration guidée : la mise au carré]
Plaçons-nous dans un repère orthonormé et posons $z = x + iy$ avec $(x,y) \in \mathbb{R}^2$. Notons $A(a)$, $B(b)$. Comme les deux membres sont positifs, l'équation équivaut à :
\[ |z - a|^2 = \alpha^2 |z - b|^2. \]
Or pour tout complexe, $|z - a|^2 = (x - x_A)^2 + (y - y_A)^2$ et de même $|z - b|^2 = (x - x_B)^2 + (y - y_B)^2$. L'équation devient :
\[ (x - x_A)^2 + (y - y_A)^2 = \alpha^2\left[(x - x_B)^2 + (y - y_B)^2\right]. \]
En développant et en regroupant :
\[ (1 - \alpha^2)x^2 + (1 - \alpha^2)y^2 + (\text{termes du premier degré}) = 0. \]
Comme $\alpha \neq 1$, on a $1 - \alpha^2 \neq 0$ : on peut diviser toute l'équation par $1 - \alpha^2$. On obtient une équation de la forme
\[ x^2 + y^2 + ux + vy + w = 0, \]
qui est l'équation d'un cercle (on le met sous forme canonique $(x - x_0)^2 + (y - y_0)^2 = R^2$ en complétant les carrés). C'est bien ce qui était annoncé. Remarque : si $\alpha = 1$, le facteur $1 - \alpha^2$ s'annule, les termes en $x^2$ et $y^2$ disparaissent et il reste une équation de \emph{droite} : on retrouve la médiatrice !
\end{experience}

\begin{exemplebox}[Le cercle d'Apollonius de $|(z-1)/(z+1)| = 2$]
Cherchons l'ensemble des points $M(z)$ tels que $\left|\dfrac{z-1}{z+1}\right| = 2$, avec $A(1)$ et $B(-1)$. La condition s'écrit $MA = 2\,MB$. Cherchons d'abord les points de la droite $(AB)$ (l'axe réel, $z = x \in \mathbb{R}$) :
\[ |x - 1| = 2|x + 1|. \]
En élevant au carré : $(x-1)^2 = 4(x+1)^2$, soit $x^2 - 2x + 1 = 4x^2 + 8x + 4$, c'est-à-dire $3x^2 + 10x + 3 = 0$. Le discriminant vaut $100 - 36 = 64$, d'où $x = \dfrac{-10 \pm 8}{6}$, soit $x = -\dfrac{1}{3}$ ou $x = -3$. Les points $C\left(-\tfrac{1}{3}\right)$ et $D(-3)$ sont donc sur le cercle, et ils sont diamétralement opposés : le centre est le milieu $\Omega\left(-\tfrac{5}{3}\right)$ et le rayon est $R = \dfrac{1}{2}\left|-\tfrac{1}{3} - (-3)\right| = \dfrac{4}{3}$.
\end{exemplebox}

\begin{popfigure}
\begin{center}
\begin{tikzpicture}[scale=1.1]
\draw[popline, very thin] (-4.5,-2.5) grid (2.5,2.5);
\draw[-{Stealth}, popdark] (-4.5,0) -- (2.5,0) node[below left] {$x$};
\draw[-{Stealth}, popdark] (0,-2.5) -- (0,2.5) node[below left] {$y$};
\draw[popdark] (0,0) node[below left] {$O$};
% Cercle d'Apollonius : centre (-5/3, 0), rayon 4/3
\draw[popblue, very thick] (-1.6667,0) circle (1.3333);
% Points A, B
\fill[popink] (1,0) circle (0.08) node[above right, popdark] {$A(1)$};
\fill[popink] (-1,0) circle (0.08) node[above left, popdark] {$B(-1)$};
% Points C et D (diamètre)
\fill[poporange] (-0.3333,0) circle (0.08) node[above right=1pt, popdark, fill=white, inner sep=1pt] {$C\left(-\frac{1}{3}\right)$};
\fill[poporange] (-3,0) circle (0.08) node[above left, popdark] {$D(-3)$};
% Centre
\fill[popgreen] (-1.6667,0) circle (0.07) node[below, popdark, fill=white, inner sep=1pt] {$\Omega\left(-\frac{5}{3}\right)$};
% Diamètre
\draw[poporange, thick] (-3,0) -- (-0.3333,0);
\end{tikzpicture}
\end{center}
\legende{Le cercle d'Apollonius $\left|\dfrac{z-1}{z+1}\right| = 2$ : ensemble des points $M$ tels que $MA = 2\,MB$. Les points $C$ et $D$ de l'axe réel vérifiant la condition forment un diamètre.}
\end{popfigure}

\courssub{Technique de réduction : factoriser le coefficient de $z$}

Dans les sujets du Baccalauréat, les conditions ne sont presque jamais données sous la forme « propre » $|z - a| = \alpha$. Elles arrivent déguisées : $|2z - 4 + 2i| = 6$, $|3z - 6i| = 9$, etc. Le réflexe salvateur consiste à \cle{factoriser} le coefficient de $z$, puis à utiliser la propriété du module d'un produit : $|k \cdot w| = |k| \cdot |w|$.

\begin{methode}[Identifier un ensemble défini par une condition sur un module]
\begin{enumerate}
\item \textbf{Factoriser} le coefficient de $z$ dans le module : $|kz + c| = |k| \cdot \left|z + \dfrac{c}{k}\right|$.
\item \textbf{Diviser} les deux membres de l'équation (ou de l'inéquation) par $|k|$ pour obtenir la forme réduite $|z - a| = r$ (ou $\leq$, $>$, etc.).
\item \textbf{Réécrire} l'affixe sous la forme $z - a$ en faisant attention aux signes : $z + 2 - i = z - (-2 + i)$.
\item \textbf{Traduire} géométriquement : $|z - a| = r$ est un cercle, $|z - a| \leq r$ un disque fermé, $|z - a| = |z - b|$ une médiatrice, $\left|\dfrac{z-a}{z-b}\right| = \alpha$ ($\alpha \neq 1$) un cercle d'Apollonius.
\item \textbf{Conclure} en précisant la nature de l'ensemble, son centre (ou ses points remarquables), son rayon, et si la frontière est incluse ou exclue.
\end{enumerate}
\end{methode}

\begin{exemplebox}[Réduction en action]
Déterminons l'ensemble des points $M(z)$ tels que $|2z - 4 + 2i| = 6$.
\begin{align*}
|2z - 4 + 2i| = 6 &\iff |2(z - 2 + i)| = 6 \\
&\iff |2| \times |z - 2 + i| = 6 \\
&\iff 2\,|z - (2 - i)| = 6 \\
&\iff |z - (2 - i)| = 3.
\end{align*}
L'ensemble cherché est le cercle de centre $\Omega(2 - i)$ et de rayon $3$.
\end{exemplebox}

\begin{attention}[L'erreur classique : oublier de diviser par le coefficient de $z$]
Face à $|3z - 6i| = 9$, beaucoup d'élèves concluent trop vite : « cercle de rayon $9$ ». C'est \textbf{faux} ! Il faut factoriser :
\[ |3z - 6i| = 9 \iff 3\,|z - 2i| = 9 \iff |z - 2i| = 3. \]
Le rayon est $3$, pas $9$. Retiens cette règle d'or : \textbf{le coefficient de $z$ doit toujours être ramené à $1$} avant toute lecture géométrique. Autre piège du même type : dans $|z - 3 + 2i|$, l'affixe du centre est $3 - 2i$ (on réécrit $z - (3 - 2i)$), et non $-3 + 2i$. Prends le temps d'écrire la forme $z - a$ explicitement.
\end{attention}

\courssub{Synthèse et représentation graphique}

Regroupons dans un tableau les quatre situations fondamentales de cette section.

\begin{center}
\begin{tabularx}{\linewidth}{|l|X|}
\hline
\rowcolor{popblueL} \textbf{Condition} ($M$ d'affixe $z$) & \textbf{Ensemble de points} \\
\hline
$|z - a| = \alpha$ \; ($\alpha > 0$) & Cercle de centre $\Omega(a)$ et de rayon $\alpha$ \\
\hline
$|z - a| \leq \alpha$ & Disque fermé de centre $\Omega(a)$ et de rayon $\alpha$ (frontière incluse) \\
\hline
$|z - a| < \alpha$ & Disque ouvert (frontière exclue, tracée en pointillés) \\
\hline
$|z - a| > \alpha$ & Extérieur du cercle de centre $\Omega(a)$ et de rayon $\alpha$ (frontière exclue) \\
\hline
$|z - a| = |z - b|$ & Médiatrice du segment $[AB]$, avec $A(a)$ et $B(b)$ \\
\hline
$\left|\dfrac{z - a}{z - b}\right| = 1$ & Médiatrice de $[AB]$ (car $MA = MB$) \\
\hline
$\left|\dfrac{z - a}{z - b}\right| = \alpha$, $\alpha > 0$, $\alpha \neq 1$ & Cercle d'Apollonius : $MA = \alpha\, MB$ \\
\hline
\end{tabularx}
\end{center}

Pour le tracé, adopte une routine systématique : place d'abord les points remarquables ($\Omega$, $A$, $B$) à partir de leurs affixes, trace la frontière (cercle au compas, médiatrice à la règle et à l'équerre, en trait plein si elle est incluse, en pointillés sinon), puis hachure la région solution en testant au besoin un point simple — par exemple l'origine $O$ si elle n'est pas sur la frontière : si ses coordonnées vérifient la condition, elle est dans la zone à hachurer.

\begin{exercice}[Pour t'entraîner]
Le plan est rapporté à un repère orthonormé direct $(O\,;\vec{u},\vec{v})$. Déterminer la nature géométrique de chacun des ensembles suivants et les représenter :
\begin{enumerate}
\item $(E_1)$ : ensemble des points $M(z)$ tels que $|z + 3 - i| = 2$ ;
\item $(E_2)$ : ensemble des points $M(z)$ tels que $|4z - 8i| \leq 12$ ;
\item $(E_3)$ : ensemble des points $M(z)$ tels que $|z - 1 - i| = |z + 3|$ ;
\item $(E_4)$ : ensemble des points $M(z)$ tels que $\left|\dfrac{z - 2}{z - i}\right| = 1$.
\end{enumerate}
\end{exercice}

\begin{corrige}[Exercice]
\begin{enumerate}
\item $|z + 3 - i| = |z - (-3 + i)| = 2$ : $(E_1)$ est le cercle de centre $\Omega(-3 + i)$ et de rayon $2$.
\item $|4z - 8i| \leq 12 \iff 4|z - 2i| \leq 12 \iff |z - 2i| \leq 3$ : $(E_2)$ est le disque fermé de centre $\Omega(2i)$ et de rayon $3$, frontière incluse.
\item Avec $A(1 + i)$ et $B(-3)$ : $|z - (1+i)| = |z - (-3)| \iff MA = MB$ : $(E_3)$ est la médiatrice de $[AB]$.
\item Pour $z \neq i$, avec $A(2)$ et $B(i)$ : $\left|\dfrac{z-2}{z-i}\right| = \dfrac{MA}{MB} = 1 \iff MA = MB$ : $(E_4)$ est la médiatrice de $[AB]$. Le point $B$ n'appartient pas à cette médiatrice, donc aucune restriction n'est à faire.
\end{enumerate}
\end{corrige}

\begin{savaistu}[Apollonius de Perge, un génie de la géométrie]
Apollonius de Perge (environ 262--190 av. J.-C.), surnommé « le Grand Géomètre », a étudié les lieux de points dont le rapport des distances à deux points fixes est constant, plus de deux mille ans avant l'invention des nombres complexes ! Ses travaux sur les coniques (ellipse, parabole, hyperbole — c'est lui qui leur a donné leurs noms) sont encore utilisés aujourd'hui pour décrire les orbites des satellites. Les cercles d'Apollonius interviennent aussi en physique moderne, notamment en optique et dans l'étude des champs électriques créés par deux charges ponctuelles.
\end{savaistu}

Nous savons désormais traduire en géométrie toute condition portant sur des \emph{modules}, c'est-à-dire sur des distances. Mais un nombre complexe non nul ne possède pas seulement un module : il possède aussi un \emph{argument}, qui code une information angulaire. C'est cette seconde coordonnée, tout aussi riche géométriquement, que nous allons explorer dans la section suivante avec les formes trigonométrique et exponentielle.

\coursec{Argument, forme trigonométrique et forme exponentielle}

Dans la section précédente, nous avons appris à localiser les points du plan complexe grâce au \cle{module} : la condition $|z - a| = \alpha$ décrit un cercle, $|z - a| \leq \alpha$ un disque. Le module répond à la question : « à quelle \emph{distance} de l'origine se trouve le point $M$ ? ». Mais une distance seule ne suffit pas à repérer un point : sur la route Douala--Yaoundé, dire « je suis à \SI{50}{km} de Douala » ne précise pas dans quelle \emph{direction}. Il nous manque une information angulaire. C'est exactement le rôle de l'\cle{argument} : donner la direction du point $M$ vue depuis l'origine. Avec le module (la distance) et l'argument (la direction), on dispose d'un système de coordonnées complet, dit \emph{polaire}, qui va révolutionner notre façon de calculer avec les complexes.

\courssub{Argument d'un nombre complexe non nul}

\textbf{Intuition.} Soit $M$ un point d'affixe $z$ non nulle, c'est-à-dire distinct de l'origine $O$. La demi-droite $[OM)$ fait un certain angle avec l'axe des réels positifs $(O;\vec{u})$. Cet angle, mesuré en radians et muni d'un sens de rotation (le sens trigonométrique, inverse des aiguilles d'une montre), est ce qu'on appelle un argument de $z$.

\begin{definition}[Argument d'un nombre complexe non nul]
Soit $z$ un nombre complexe non nul et $M$ son image dans le plan complexe $(O;\vec{u},\vec{v})$. On appelle \cle{argument} de $z$ toute mesure, en radians, de l'angle orienté $(\vec{u},\overrightarrow{OM})$. On note :
\[ \arg(z) = \theta \ [2\pi] \]
ce qui signifie que $\theta$ est défini à $2k\pi$ près ($k \in \mathbb{Z}$) : un complexe non nul possède une \textbf{infinité} d'arguments, qui diffèrent tous d'un multiple de $2\pi$.
\end{definition}

\begin{attention}[L'argument de zéro n'existe pas]
Le nombre $z = 0$ n'a \textbf{pas d'argument} : le point $O$ ne définit aucune direction. Toutes les notions de cette section (argument, forme trigonométrique, forme exponentielle) supposent $z \neq 0$. Avant d'écrire $\arg(z)$, vérifie toujours que $z \neq 0$.
\end{attention}

Parmi cette infinité d'arguments, on en privilégie un seul, pour pouvoir parler sans ambiguïté.

\begin{definition}[Argument principal]
Soit $z$ un complexe non nul. L'\cle{argument principal} de $z$ est l'\textbf{unique} argument de $z$ appartenant à l'intervalle $\left]-\pi\,;\,\pi\right]$. On le note parfois $\mathrm{Arg}(z)$.
\end{definition}

Par exemple, les réels $\dfrac{\pi}{3}$, $\dfrac{\pi}{3} + 2\pi = \dfrac{7\pi}{3}$ et $\dfrac{\pi}{3} - 2\pi = -\dfrac{5\pi}{3}$ sont tous des arguments d'un même complexe, mais seul $\dfrac{\pi}{3}$ est l'argument principal.

\begin{propriete}[Arguments des cas particuliers]
Pour un complexe non nul $z$ :
\begin{itemize}
\item $z$ est un \textbf{réel strictement positif} si et seulement si $\arg(z) = 0 \ [2\pi]$ ;
\item $z$ est un \textbf{réel strictement négatif} si et seulement si $\arg(z) = \pi \ [2\pi]$ ;
\item $z$ est un \textbf{imaginaire pur} de partie imaginaire positive si et seulement si $\arg(z) = \dfrac{\pi}{2} \ [2\pi]$ ;
\item $z$ est un \textbf{imaginaire pur} de partie imaginaire négative si et seulement si $\arg(z) = -\dfrac{\pi}{2} \ [2\pi]$.
\end{itemize}
\end{propriete}

Ces résultats sont immédiats géométriquement : un réel positif se trouve sur la demi-droite $[O;\vec{u})$, un réel négatif sur la demi-droite opposée, un imaginaire pur sur l'axe $(O;\vec{v})$.

\begin{propriete}[Lecture géométrique : argument d'une différence]
Soient $A$ et $B$ deux points distincts d'affixes $z_A$ et $z_B$. Alors :
\[ \arg(z_B - z_A) = (\vec{u},\overrightarrow{AB}) \ [2\pi]. \]
Autrement dit, l'argument de la différence des affixes donne la \textbf{direction du vecteur} $\overrightarrow{AB}$ par rapport à l'axe des réels.
\end{propriete}

En effet, le vecteur d'affixe $z_B - z_A$ est précisément $\overrightarrow{AB}$ : on applique la définition de l'argument à ce vecteur.

\courssub{Forme trigonométrique}

\textbf{Intuition.} Plaçons $M(z)$ avec $z = x + iy$ non nul, et notons $r = |z| = OM$ et $\theta$ un argument de $z$. En projetant $M$ sur les axes, la trigonométrie du collège (dans le cercle de rayon $r$) donne immédiatement $x = r\cos\theta$ et $y = r\sin\theta$. Le point $M$ est donc entièrement déterminé par le couple $(r,\theta)$.

\begin{propriete}[Forme trigonométrique — démonstration exigible]
Tout nombre complexe non nul $z$ s'écrit de manière unique (à $2\pi$ près pour l'angle) :
\[ z = r(\cos\theta + i\sin\theta) \quad \text{avec } r = |z| > 0 \text{ et } \theta = \arg(z)\ [2\pi]. \]
Cette écriture est appelée \cle{forme trigonométrique} de $z$.
\end{propriete}

\begin{experience}[Démonstration guidée du théorème]
Soit $z$ un complexe non nul, $r = |z| > 0$.
\begin{enumerate}
\item \textbf{Ramener au cercle unité.} Posons $Z = \dfrac{z}{r}$. Alors $|Z| = \dfrac{|z|}{r} = \dfrac{r}{r} = 1$ : le complexe $Z$ appartient au cercle trigonométrique.
\item \textbf{Paramétrer le cercle unité.} Écrivons $Z = a + ib$ avec $a, b \in \mathbb{R}$. La condition $|Z| = 1$ s'écrit $a^2 + b^2 = 1$. Or tout point du cercle trigonométrique est de la forme $(\cos\theta\,;\,\sin\theta)$ : il existe donc un réel $\theta$ tel que $a = \cos\theta$ et $b = \sin\theta$.
\item \textbf{Conclure.} On a $Z = \cos\theta + i\sin\theta$, donc $z = rZ = r(\cos\theta + i\sin\theta)$. De plus, l'image de $z$ étant sur la demi-droite portée par l'image de $Z$, l'angle $\theta$ est bien une mesure de $(\vec{u},\overrightarrow{OM})$, c'est-à-dire un argument de $z$. $\square$
\end{enumerate}
\end{experience}

\begin{popfigure}
\begin{tikzpicture}[scale=2.4,>=Stealth]
\draw[->,popmuted] (-1.5,0) -- (1.9,0) node[below right] {Réels};
\draw[->,popmuted] (0,-1.3) -- (0,1.6) node[above left] {Imaginaires};
\draw[popline] (0,0) circle (1);
\coordinate (O) at (0,0);
\coordinate (M) at (50:1.35);
\draw[thick,popblue] (O) -- (M) node[above right,popdark] {$M(z)$};
\draw[poporange,thick] (0.55,0) arc (0:50:0.55) node[midway,right] {$\theta$};
\draw[dashed,popgreen] (M) -- (M |- O) node[below,popdark] {$r\cos\theta$};
\draw[dashed,popgreen] (M) -- (O -| M) node[left,popdark] {$r\sin\theta$};
\fill[popink] (M) circle (0.025);
\fill[popdark] (O) circle (0.02) node[below left] {$O$};
\node[popblue] at (25:0.75) {$r = |z|$};
\end{tikzpicture}
\legende{Un point $M(z)$ repéré par son module $r$ et son argument $\theta$ : $z = r(\cos\theta + i\sin\theta)$.}
\end{popfigure}

\begin{attention}[Un piège classique : le module doit être positif]
L'écriture $z = -2\left(\cos\dfrac{\pi}{6} + i\sin\dfrac{\pi}{6}\right)$ n'est \textbf{pas} une forme trigonométrique, car le facteur $-2$ est négatif : or le module est par définition positif. Pour corriger, on utilise $-1 = \cos\pi + i\sin\pi$ et les formules d'addition des angles :
\[ z = 2\left(\cos\left(\pi + \tfrac{\pi}{6}\right) + i\sin\left(\pi + \tfrac{\pi}{6}\right)\right) = 2\left(\cos\tfrac{7\pi}{6} + i\sin\tfrac{7\pi}{6}\right). \]
Retiens : dans une forme trigonométrique $r(\cos\theta + i\sin\theta)$, on exige $r > 0$ et le \textbf{même} angle $\theta$ dans le cosinus et le sinus, reliés par un signe $+$.
\end{attention}

\courssub{Notation exponentielle}

\begin{definition}[Exponentielle complexe et forme exponentielle]
Pour tout réel $\theta$, on pose :
\[ e^{i\theta} = \cos\theta + i\sin\theta. \]
Tout complexe non nul s'écrit alors sous \cle{forme exponentielle} :
\[ z = r\,e^{i\theta} \quad \text{avec } r = |z| > 0 \text{ et } \theta = \arg(z)\ [2\pi]. \]
\end{definition}

Cette notation, due à Euler, n'est pas une simple commodité d'écriture : elle possède les mêmes propriétés algébriques que l'exponentielle réelle ($e^{i\theta} \times e^{i\theta'} = e^{i(\theta+\theta')}$, etc.), ce que nous exploiterons dans la section suivante. Remarque au passage que $|e^{i\theta}| = \cos^2\theta + \sin^2\theta = 1$ : les nombres $e^{i\theta}$ parcourent exactement le cercle trigonométrique.

\begin{popfigure}
\begin{tikzpicture}[scale=2.6,>=Stealth]
\draw[->,popmuted] (-1.4,0) -- (1.5,0) node[below right] {Réels};
\draw[->,popmuted] (0,-1.3) -- (0,1.4) node[above left] {Imaginaires};
\draw[popline] (0,0) circle (1);
\foreach \a/\n/\p in {30/{\frac{\pi}{6} : e^{i\pi/6} = \frac{\sqrt3}{2} + \frac{i}{2}}/right,
                       45/{\frac{\pi}{4} : e^{i\pi/4} = \frac{\sqrt2}{2} + i\frac{\sqrt2}{2}}/above right,
                       60/{\frac{\pi}{3} : e^{i\pi/3} = \frac{1}{2} + i\frac{\sqrt3}{2}}/above right,
                       90/{\frac{\pi}{2} : e^{i\pi/2} = i}/above} {
  \draw[popblue] (0,0) -- (\a:1);
  \fill[popink] (\a:1) circle (0.022);
  \node[popdark,font=\small,\p] at (\a:1.08) {$\n$};
}
\draw[poporange,thick] (0.35,0) arc (0:30:0.35);
\draw[poporange,thick] (0.45,0) arc (0:45:0.45);
\draw[poporange,thick] (0.55,0) arc (0:60:0.55);
\draw[poporange,thick] (0.25,0) arc (0:90:0.25);
\fill[popdark] (0,0) circle (0.02) node[below left] {$O$};
\fill[popdark] (1,0) circle (0.02) node[below right] {$1$};
\fill[popdark] (-1,0) circle (0.02) node[below left] {$-1 = e^{i\pi}$};
\end{tikzpicture}
\legende{Le cercle trigonométrique : les arguments remarquables et les complexes $e^{i\theta}$ associés.}
\end{popfigure}

\begin{savaistu}[La plus belle formule des mathématiques]
En prenant $\theta = \pi$, on obtient $e^{i\pi} = \cos\pi + i\sin\pi = -1$, c'est-à-dire :
\[ e^{i\pi} + 1 = 0. \]
Cette \cle{identité d'Euler} relie en une seule ligne les cinq constantes fondamentales des mathématiques : $0$, $1$, $e$, $\pi$ et $i$. Elle est régulièrement citée comme la plus belle formule des mathématiques. Son auteur, Leonhard Euler (1707--1783), est le mathématicien le plus prolifique de l'histoire : ses œuvres complètes occupent plus de 80 volumes !
\end{savaistu}

\courssub{Passage de la forme algébrique à la forme trigonométrique}

\begin{methode}[Algorithme en 3 étapes : algébrique $\to$ trigonométrique]
Soit $z = x + iy$ un complexe non nul donné sous forme algébrique.
\begin{enumerate}
\item \textbf{Calculer le module} : $r = |z| = \sqrt{x^2 + y^2}$.
\item \textbf{Factoriser par $r$} : $z = r\left(\dfrac{x}{r} + i\,\dfrac{y}{r}\right)$, ce qui fournit $\cos\theta = \dfrac{x}{r}$ et $\sin\theta = \dfrac{y}{r}$.
\item \textbf{Identifier $\theta$} : repérer sur le cercle trigonométrique l'unique angle de $\left]-\pi\,;\,\pi\right]$ ayant ce cosinus et ce sinus (le cosinus seul ne suffit pas : deux angles opposés ont le même cosinus !). Conclure : $z = r(\cos\theta + i\sin\theta) = r e^{i\theta}$.
\end{enumerate}
\end{methode}

\begin{exemplebox}[Deux exemples chiffrés complets]
\textbf{1.} Soit $z = 1 + i\sqrt{3}$.
\begin{align*}
r &= \sqrt{1^2 + (\sqrt{3})^2} = \sqrt{4} = 2,\\
z &= 2\left(\frac{1}{2} + i\,\frac{\sqrt{3}}{2}\right) \quad\Rightarrow\quad \cos\theta = \frac{1}{2},\ \sin\theta = \frac{\sqrt{3}}{2} \quad\Rightarrow\quad \theta = \frac{\pi}{3}\ [2\pi].
\end{align*}
Donc $z = 2\left(\cos\dfrac{\pi}{3} + i\sin\dfrac{\pi}{3}\right) = 2e^{i\pi/3}$.

\textbf{2.} Soit $z = -\sqrt{2} - i\sqrt{2}$.
\begin{align*}
r &= \sqrt{(-\sqrt{2})^2 + (-\sqrt{2})^2} = \sqrt{4} = 2,\\
z &= 2\left(-\frac{\sqrt{2}}{2} - i\,\frac{\sqrt{2}}{2}\right) \quad\Rightarrow\quad \cos\theta = -\frac{\sqrt{2}}{2},\ \sin\theta = -\frac{\sqrt{2}}{2}.
\end{align*}
Le cosinus et le sinus sont tous deux négatifs : $\theta$ est dans le troisième quadrant. L'argument principal est $\theta = -\dfrac{3\pi}{4}$.
Donc $z = 2\left(\cos\left(-\tfrac{3\pi}{4}\right) + i\sin\left(-\tfrac{3\pi}{4}\right)\right) = 2e^{-3i\pi/4}$.
\end{exemplebox}

\courssub{Passage de la forme trigonométrique à la forme algébrique}

Le sens inverse est un simple calcul : on \textbf{développe} $r(\cos\theta + i\sin\theta)$ en remplaçant $\cos\theta$ et $\sin\theta$ par leurs valeurs remarquables.

\begin{exoresolu}[Forme algébrique de $z = 4e^{2i\pi/3}$]
Énoncé. Déterminer la forme algébrique de $z = 4\left(\cos\dfrac{2\pi}{3} + i\sin\dfrac{2\pi}{3}\right)$, puis placer son image dans le plan complexe.
\tcblower
Solution détaillée. On utilise les valeurs remarquables : $\cos\dfrac{2\pi}{3} = -\dfrac{1}{2}$ et $\sin\dfrac{2\pi}{3} = \dfrac{\sqrt{3}}{2}$. Donc :
\[ z = 4\left(-\frac{1}{2} + i\,\frac{\sqrt{3}}{2}\right) = -2 + 2i\sqrt{3}. \]
Vérification : $|z| = \sqrt{(-2)^2 + (2\sqrt{3})^2} = \sqrt{4 + 12} = 4$, conforme au module donné. L'image de $z$ est le point de coordonnées $(-2\,;\,2\sqrt{3})$, situé dans le deuxième quadrant, ce qui est cohérent avec un argument de $\dfrac{2\pi}{3}$.
\end{exoresolu}

\begin{attention}[Deux réflexes de vérification]
Après tout passage d'une forme à l'autre, vérifie que :
\begin{enumerate}
\item le \textbf{module} est conservé : recalcule $\sqrt{x^2+y^2}$ et compare à $r$ ;
\item le \textbf{quadrant} est cohérent : les signes de $x$ et $y$ doivent correspondre à la position de $\theta$ sur le cercle trigonométrique. Un argument de $\dfrac{2\pi}{3}$ avec une forme algébrique à partie réelle positive est forcément faux.
\end{enumerate}
\end{attention}

\begin{aretenir}[L'essentiel de la section]
\begin{itemize}
\item Pour $z \neq 0$ d'image $M$ : $\arg(z) = (\vec{u},\overrightarrow{OM})\ [2\pi]$ ; l'argument principal est l'unique argument dans $\left]-\pi\,;\,\pi\right]$.
\item Formes trigonométrique et exponentielle : $z = r(\cos\theta + i\sin\theta) = r e^{i\theta}$, avec $r = |z| > 0$.
\item Cas particuliers : réel positif $\Rightarrow \theta = 0$ ; réel négatif $\Rightarrow \theta = \pi$ ; imaginaire pur $\Rightarrow \theta = \pm\dfrac{\pi}{2}$.
\item Lecture géométrique : $\arg(z_B - z_A) = (\vec{u},\overrightarrow{AB})\ [2\pi]$.
\item Algébrique $\to$ trigonométrique : $r = \sqrt{x^2+y^2}$, puis $\cos\theta = \dfrac{x}{r}$, $\sin\theta = \dfrac{y}{r}$, et on identifie $\theta$ avec les \textbf{deux} conditions.
\end{itemize}
\end{aretenir}

\begin{exercice}[Pour s'entraîner]
\begin{enumerate}
\item Déterminer la forme trigonométrique et la forme exponentielle de : $z_1 = -1 + i$ ; $z_2 = \sqrt{3} - i$ ; $z_3 = -5i$.
\item Déterminer la forme algébrique de : $z_4 = 3e^{-i\pi/6}$ ; $z_5 = 2\left(\cos\dfrac{5\pi}{6} + i\sin\dfrac{5\pi}{6}\right)$.
\item Soient $A(1)$ et $B(1 + 2i)$. Déterminer une mesure de l'angle $(\vec{u},\overrightarrow{AB})$.
\end{enumerate}
\end{exercice}

\begin{corrige}[Exercice]
\begin{enumerate}
\item $z_1 = -1 + i$ : $r = \sqrt{2}$, $\cos\theta = -\dfrac{\sqrt{2}}{2}$, $\sin\theta = \dfrac{\sqrt{2}}{2}$, donc $\theta = \dfrac{3\pi}{4}$ et $z_1 = \sqrt{2}\,e^{3i\pi/4}$.
$z_2 = \sqrt{3} - i$ : $r = 2$, $\cos\theta = \dfrac{\sqrt{3}}{2}$, $\sin\theta = -\dfrac{1}{2}$, donc $\theta = -\dfrac{\pi}{6}$ et $z_2 = 2e^{-i\pi/6}$.
$z_3 = -5i$ est un imaginaire pur négatif : $z_3 = 5e^{-i\pi/2}$.
\item $z_4 = 3\left(\cos\dfrac{\pi}{6} - i\sin\dfrac{\pi}{6}\right) = 3\left(\dfrac{\sqrt{3}}{2} - \dfrac{i}{2}\right) = \dfrac{3\sqrt{3}}{2} - \dfrac{3}{2}i$.
$z_5 = 2\left(-\dfrac{\sqrt{3}}{2} + \dfrac{i}{2}\right) = -\sqrt{3} + i$.
\item $\overrightarrow{AB}$ a pour affixe $z_B - z_A = 2i$, imaginaire pur de partie imaginaire positive : $(\vec{u},\overrightarrow{AB}) = \dfrac{\pi}{2}\ [2\pi]$.
\end{enumerate}
\end{corrige}

Nous savons désormais écrire un complexe sous trois formes : algébrique, trigonométrique et exponentielle. Mais toute la puissance de la forme exponentielle apparaîtra dans la section suivante, où nous verrons que \emph{multiplier des complexes, c'est multiplier les modules et additionner les arguments} — une propriété qui transforme les calculs et éclaire la géométrie.

\coursec{Opérations sur les arguments et applications géométriques}

Dans la section précédente, nous avons appris à écrire un nombre complexe non nul sous forme trigonométrique $z = r(\cos\theta + i\sin\theta)$ ou exponentielle $z = re^{i\theta}$, où $\theta$ est un argument de $z$. Mais toute la puissance de cette écriture apparaît lorsqu'on fait des \emph{opérations} : multiplier deux complexes revient à multiplier leurs modules et \emph{additionner leurs arguments}. C'est cette propriété remarquable qui va transformer le plan complexe en une véritable « machine à géométrie » : angles, alignements, perpendicularités deviendront de simples calculs d'arguments. C'est l'objet de cette section, l'une des plus rentables de tout le programme de Terminale pour le Baccalauréat.

\courssub{Arguments d'un produit, d'un quotient, d'une puissance}

\textbf{Intuition : multiplier, c'est « tourner et dilater ».}\par
Prenons deux complexes en forme exponentielle : $z = re^{i\theta}$ et $z' = r'e^{i\theta'}$. Si l'on applique naïvement la règle des exposants (que nous justifierons rigoureusement), on obtient :
\[ zz' = rr'\,e^{i(\theta+\theta')}. \]
Autrement dit, le point image de $zz'$ s'obtient à partir de l'image de $z$ en \textbf{dilatant} d'un facteur $r'$ et en \textbf{tournant} d'un angle $\theta'$. Multiplier par $i = e^{i\pi/2}$, c'est tourner d'un quart de tour ; multiplier par $-1 = e^{i\pi}$, c'est faire un demi-tour. Toute la géométrie des arguments découle de cette idée.

\begin{propriete}[Propriétés algébriques des arguments]
Pour tous nombres complexes \textbf{non nuls} $z$ et $z'$, et pour tout entier $n \in \mathbb{Z}$ :
\begin{align*}
|zz'| &= |z|\,|z'| & \arg(zz') &= \arg(z) + \arg(z') \ [2\pi] \\
\left|\dfrac{1}{z}\right| &= \dfrac{1}{|z|} & \arg\left(\dfrac{1}{z}\right) &= -\arg(z) \ [2\pi] \\
\left|\dfrac{z}{z'}\right| &= \dfrac{|z|}{|z'|} & \arg\left(\dfrac{z}{z'}\right) &= \arg(z) - \arg(z') \ [2\pi] \\
|z^n| &= |z|^n & \arg(z^n) &= n\arg(z) \ [2\pi]
\end{align*}
On a de plus, pour $z \neq 0$ :
\[ \arg(\bar{z}) = -\arg(z) \ [2\pi] \qquad \text{et} \qquad \arg(-z) = \arg(z) + \pi \ [2\pi]. \]
\emph{Démonstration exigible au Baccalauréat pour le produit ; les autres formules s'en déduisent.}
\end{propriete}

\begin{experience}[Démonstration guidée : argument d'un produit]
Écrivons $z$ et $z'$ sous forme trigonométrique : $z = r(\cos\theta + i\sin\theta)$ et $z' = r'(\cos\theta' + i\sin\theta')$.

\textbf{Étape 1 : développer le produit.}
\begin{align*}
zz' &= rr'\big(\cos\theta + i\sin\theta\big)\big(\cos\theta' + i\sin\theta'\big) \\
&= rr'\Big[\big(\cos\theta\cos\theta' - \sin\theta\sin\theta'\big) + i\big(\sin\theta\cos\theta' + \cos\theta\sin\theta'\big)\Big].
\end{align*}

\textbf{Étape 2 : reconnaître les formules d'addition.} Rappelons, pour tous réels $a$ et $b$ :
\[ \cos(a+b) = \cos a\cos b - \sin a\sin b \quad ; \quad \sin(a+b) = \sin a\cos b + \cos a\sin b. \]
Le produit s'écrit donc :
\[ zz' = rr'\big(\cos(\theta+\theta') + i\sin(\theta+\theta')\big). \]

\textbf{Étape 3 : conclure.} Comme $rr' > 0$, cette écriture est bien une forme trigonométrique de $zz'$. On en déduit $|zz'| = rr' = |z||z'|$ et $\arg(zz') = \theta + \theta' = \arg(z) + \arg(z') \ [2\pi]$. $\blacksquare$

\textbf{Conséquences (à savoir retrouver).}
\begin{itemize}
\item Pour l'inverse : $z \times \dfrac{1}{z} = 1$, donc $\arg(z) + \arg\left(\dfrac{1}{z}\right) = \arg(1) = 0 \ [2\pi]$, d'où $\arg\left(\dfrac{1}{z}\right) = -\arg(z) \ [2\pi]$.
\item Pour le quotient : $\dfrac{z}{z'} = z \times \dfrac{1}{z'}$, donc $\arg\left(\dfrac{z}{z'}\right) = \arg(z) - \arg(z') \ [2\pi]$.
\item Pour la puissance : par récurrence, $\arg(z^n) = n\arg(z) \ [2\pi]$ pour $n \in \mathbb{N}$, puis pour $n \in \mathbb{Z}$ grâce à $z^{-n} = 1/z^n$.
\item Pour le conjugué : si $z = re^{i\theta}$, alors $\bar{z} = re^{-i\theta}$, d'où $\arg(\bar{z}) = -\arg(z) \ [2\pi]$.
\item Pour l'opposé : $-z = z \times (-1)$ et $\arg(-1) = \pi \ [2\pi]$, d'où $\arg(-z) = \arg(z) + \pi \ [2\pi]$.
\end{itemize}
\end{experience}

\begin{exemplebox}[Calculs d'arguments]
Soit $z = 1 + i$ et $z' = -\sqrt{3} + i$. On a $|z| = \sqrt{2}$ et $\arg(z) = \dfrac{\pi}{4} \ [2\pi]$ ; $|z'| = 2$ et $\arg(z') = \dfrac{5\pi}{6} \ [2\pi]$. Alors :
\begin{align*}
\arg(zz') &= \dfrac{\pi}{4} + \dfrac{5\pi}{6} = \dfrac{13\pi}{12} \ [2\pi] ;\\
\arg\left(\dfrac{z}{z'}\right) &= \dfrac{\pi}{4} - \dfrac{5\pi}{6} = -\dfrac{7\pi}{12} \ [2\pi] ;\\
\arg(z^4) &= 4 \times \dfrac{\pi}{4} = \pi \ [2\pi], \quad \text{donc } z^4 \text{ est un réel négatif.}
\end{align*}
Vérification : $z^2 = (1+i)^2 = 2i$, donc $z^4 = (2i)^2 = -4$, qui a bien pour argument $\pi$.
\end{exemplebox}

\begin{attention}[Conditions d'application et réduction modulo $2\pi$]
\begin{itemize}
\item Toutes ces formules exigent des complexes \textbf{non nuls} : $0$ n'a pas d'argument.
\item Les égalités sont toujours \textbf{modulo $2\pi$}. Si l'on demande l'argument \emph{principal}, il faut ramener le résultat dans $]{-\pi}\,;\,\pi]$. Par exemple $\arg(zz') = \dfrac{13\pi}{12} \ [2\pi]$ donne l'argument principal $\dfrac{13\pi}{12} - 2\pi = -\dfrac{11\pi}{12}$.
\item $\arg(z^n) = n\arg(z)$ est valable pour $n \in \mathbb{Z}$ tout entier, mais attention : l'argument principal de $z^n$ n'est pas forcément $n$ fois l'argument principal de $z$.
\end{itemize}
\end{attention}

\courssub{Interprétation géométrique : l'argument d'un quotient est un angle}

Voici le résultat central de cette section, celui qui fait le pont définitif entre calcul algébrique et géométrie plane.

\begin{propriete}[Propriété fondamentale : angle de deux vecteurs]
Soient $A$, $B$, $C$ trois points deux à deux distincts, d'affixes respectives $z_A$, $z_B$, $z_C$. Alors :
\[ \left(\overrightarrow{AB}\,,\,\overrightarrow{AC}\right) = \arg\left(\dfrac{z_C - z_A}{z_B - z_A}\right) \ [2\pi]. \]
Plus généralement, pour quatre points $A$, $B$, $C$, $D$ avec $A \neq B$ et $C \neq D$ :
\[ \left(\overrightarrow{AB}\,,\,\overrightarrow{CD}\right) = \arg\left(\dfrac{z_D - z_C}{z_B - z_A}\right) \ [2\pi]. \]
\end{propriete}

\textbf{Justification.} Le vecteur $\overrightarrow{AB}$ a pour affixe $z_B - z_A$ et le vecteur $\overrightarrow{AC}$ a pour affixe $z_C - z_A$. Or, pour deux vecteurs non nuls d'affixes $w_1$ et $w_2$, l'angle $(\vec{w}_1, \vec{w}_2)$ vaut $\arg(w_2) - \arg(w_1) \ [2\pi]$ : c'est l'angle dont il faut tourner $\vec{w}_1$ pour le superposer à $\vec{w}_2$. En appliquant la formule de l'argument d'un quotient :
\[ \left(\overrightarrow{AB}\,,\,\overrightarrow{AC}\right) = \arg(z_C - z_A) - \arg(z_B - z_A) = \arg\left(\dfrac{z_C - z_A}{z_B - z_A}\right) \ [2\pi]. \]

\begin{popfigure}
\begin{center}
\begin{tikzpicture}[scale=1.6,>=Stealth]
\coordinate (O) at (0,0);
\coordinate (A) at (0.6,0.4);
\coordinate (B) at (2.9,0.8);
\coordinate (C) at (1.7,2.1);
% axes
\draw[->,popmuted] (-0.4,0) -- (3.6,0) node[below left] {$\mathbb{R}$};
\draw[->,popmuted] (0,-0.4) -- (0,2.7) node[below left] {$i\mathbb{R}$};
% vecteurs
\draw[->,very thick,popblue] (A) -- (B) node[midway,below] {$\overrightarrow{AB}$};
\draw[->,very thick,poporange] (A) -- (C) node[midway,above left] {$\overrightarrow{AC}$};
% points
\fill[popdark] (A) circle (1.4pt) node[below left] {$A\,(z_A)$};
\fill[popdark] (B) circle (1.4pt) node[right] {$B\,(z_B)$};
\fill[popdark] (C) circle (1.4pt) node[above] {$C\,(z_C)$};
% angle
\draw[->,popink,thick] ($(A)+(8:0.9)$) arc (8:62:0.9) node[midway,right,xshift=4pt] {$\theta$};
\node[popink] at (3.35,2.3) {$\theta = \arg\!\left(\dfrac{z_C - z_A}{z_B - z_A}\right)\ [2\pi]$};
\end{tikzpicture}
\end{center}
\legende{L'angle orienté $(\overrightarrow{AB},\overrightarrow{AC})$ est un argument du quotient $\dfrac{z_C - z_A}{z_B - z_A}$.}
\end{popfigure}

\begin{methode}[Calculer un angle géométrique avec les complexes]
Pour déterminer une mesure de l'angle $(\overrightarrow{AB}, \overrightarrow{AC})$ :
\begin{enumerate}
\item Former les affixes des vecteurs : $z_{\overrightarrow{AB}} = z_B - z_A$ et $z_{\overrightarrow{AC}} = z_C - z_A$.
\item Calculer le quotient $Z = \dfrac{z_C - z_A}{z_B - z_A}$ (numérateur = affixe du \emph{second} vecteur, dénominateur = affixe du \emph{premier} vecteur).
\item Mettre $Z$ sous forme algébrique simplifiée (multiplier par le conjugué du dénominateur).
\item Déterminer un argument de $Z$ (forme trigonométrique ou lecture directe si $Z$ est réel ou imaginaire pur).
\item Conclure : $(\overrightarrow{AB}, \overrightarrow{AC}) = \arg(Z) \ [2\pi]$.
\end{enumerate}
\end{methode}

\begin{attention}[Le piège du quotient inversé]
C'est \textbf{l'erreur la plus fréquente} au Baccalauréat. Retiens bien :
\[ \arg\left(\dfrac{z_C - z_A}{z_B - z_A}\right) = \left(\overrightarrow{AB}\,,\,\overrightarrow{AC}\right) \quad \text{et} \quad \arg\left(\dfrac{z_B - z_A}{z_C - z_A}\right) = \left(\overrightarrow{AC}\,,\,\overrightarrow{AB}\right). \]
Le \emph{dénominateur} correspond au \emph{premier} vecteur de l'angle, le \emph{numérateur} au \emph{second}. Inverser le quotient change l'angle en son opposé : un angle de $+\dfrac{\pi}{3}$ devient $-\dfrac{\pi}{3}$, et un triangle « rectangle direct » devient « rectangle indirect ». Vérifie toujours l'ordre des points dans la question posée avant d'écrire le quotient.
\end{attention}

\courssub{Caractérisations de l'alignement et de la perpendicularité}

De la propriété fondamentale découlent deux critères immédiats et extrêmement utiles.

\begin{propriete}[Alignement et perpendicularité]
Soient $A$, $B$, $C$ trois points deux à deux distincts.
\begin{itemize}
\item Les points $A$, $B$, $C$ sont \cle{alignés} si et seulement si
\[ \arg\left(\dfrac{z_C - z_A}{z_B - z_A}\right) = 0 \ [\pi] \quad \Longleftrightarrow \quad \dfrac{z_C - z_A}{z_B - z_A} \in \mathbb{R}. \]
\item Les droites $(AB)$ et $(AC)$ sont \cle{perpendiculaires} si et seulement si
\[ \arg\left(\dfrac{z_C - z_A}{z_B - z_A}\right) = \dfrac{\pi}{2} \ [\pi] \quad \Longleftrightarrow \quad \dfrac{z_C - z_A}{z_B - z_A} \in i\mathbb{R} \ (\text{imaginaire pur}). \]
\end{itemize}
\end{propriete}

\textbf{Justification.} Dire que $A$, $B$, $C$ sont alignés, c'est dire que les vecteurs $\overrightarrow{AB}$ et $\overrightarrow{AC}$ sont colinéaires, c'est-à-dire que leur angle vaut $0$ ou $\pi$ modulo $2\pi$, autrement dit $0 \ [\pi]$. Un complexe non nul d'argument $0 \ [\pi]$ est un réel. De même, la perpendicularité équivaut à un angle de $\pm\dfrac{\pi}{2}$, c'est-à-dire $\dfrac{\pi}{2} \ [\pi]$, ce qui caractérise les imaginaires purs non nuls.

\begin{aretenir}[Le réflexe à acquérir]
Pour étudier la nature d'une configuration avec les complexes :
\begin{itemize}
\item quotient \textbf{réel} $\Rightarrow$ vecteurs colinéaires, points alignés ;
\item quotient \textbf{imaginaire pur} $\Rightarrow$ vecteurs orthogonaux, droites perpendiculaires ;
\item quotient de \textbf{module 1} $\Rightarrow$ vecteurs de même norme, côtés de même longueur.
\end{itemize}
Et n'oublie pas le critère des modules : $AB = AC \iff |z_B - z_A| = |z_C - z_A|$.
\end{aretenir}

\courssub{Application : étude de la nature d'un triangle}

\begin{exoresolu}[Un triangle rectangle... mais est-il isocèle ?]
Le plan complexe est rapporté à un repère orthonormé direct $(O\,;\vec{u},\vec{v})$. On considère les points $A$, $B$, $C$ d'affixes respectives $z_A = 1$, $z_B = 2 + i$, $z_C = 3i$.

Déterminer une mesure de l'angle $\left(\overrightarrow{BA}\,,\,\overrightarrow{BC}\right)$, puis préciser la nature exacte du triangle $ABC$.
\tcblower
\textbf{Étape 1 : former le quotient adapté à l'angle en $B$.} Le sommet visé est $B$, donc les vecteurs en jeu sont $\overrightarrow{BA}$ et $\overrightarrow{BC}$, d'affixes :
\[ z_A - z_B = 1 - (2 + i) = -1 - i \quad ; \quad z_C - z_B = 3i - (2 + i) = -2 + 2i. \]
On forme le quotient (affixe du second vecteur sur affixe du premier) :
\[ Z = \dfrac{z_C - z_B}{z_A - z_B} = \dfrac{-2 + 2i}{-1 - i}. \]

\textbf{Étape 2 : simplifier.} On multiplie numérateur et dénominateur par le conjugué du dénominateur :
\[ Z = \dfrac{(-2 + 2i)(-1 + i)}{(-1 - i)(-1 + i)} = \dfrac{2 - 2i - 2i + 2i^2}{1 + 1} = \dfrac{2 - 4i - 2}{2} = \dfrac{-4i}{2} = -2i. \]

\textbf{Étape 3 : lire le module et l'argument.} $Z = -2i$ est un imaginaire pur, $|Z| = 2$ et $\arg(Z) = -\dfrac{\pi}{2} \ [2\pi]$. Donc :
\[ \left(\overrightarrow{BA}\,,\,\overrightarrow{BC}\right) = -\dfrac{\pi}{2} \ [2\pi]. \]
L'angle en $B$ est droit : le triangle $ABC$ est \textbf{rectangle en $B$} (rectangle \emph{indirect}, car l'angle est négatif).

\textbf{Étape 4 : comparer les longueurs.} Le module du quotient donne directement le rapport des longueurs :
\[ \dfrac{BC}{BA} = |Z| = 2 \quad \Longrightarrow \quad BC = 2\,BA. \]
Vérifions par le calcul : $BA = |z_A - z_B| = |-1 - i| = \sqrt{2}$ et $BC = |z_C - z_B| = |-2 + 2i| = \sqrt{4 + 4} = 2\sqrt{2}$. Donc $BA \neq BC$ : le triangle n'est \textbf{pas isocèle en $B$}. Serait-il isocèle ailleurs ?
\[ AC = |z_C - z_A| = |3i - 1| = \sqrt{1 + 9} = \sqrt{10}. \]
Les trois longueurs $\sqrt{2}$, $2\sqrt{2}$, $\sqrt{10}$ sont deux à deux distinctes : le triangle n'est isocèle en aucun sommet. Vérification de cohérence par Pythagore : $AB^2 + BC^2 = 2 + 8 = 10 = AC^2$. 

\textbf{Conclusion :} $ABC$ est rectangle en $B$, non isocèle.

\emph{Remarque : cet exemple montre qu'il faut toujours vérifier séparément les deux conditions (angle droit ET égalité de longueurs) avant de conclure « rectangle isocèle ». Un quotient imaginaire pur ne suffit pas : il faut en plus qu'il soit de module $1$.}
\end{exoresolu}

\begin{exoresolu}[Cette fois, un vrai triangle rectangle isocèle]
On considère les points $A(1)$, $B(2 + i)$ et $C'(2 - i)$. Montrer que le triangle $ABC'$ est rectangle isocèle en $A$.
\tcblower
\textbf{Quotient pour l'angle en $A$ :}
\[ Z = \dfrac{z_{C'} - z_A}{z_B - z_A} = \dfrac{(2 - i) - 1}{(2 + i) - 1} = \dfrac{1 - i}{1 + i} = \dfrac{(1 - i)^2}{(1 + i)(1 - i)} = \dfrac{1 - 2i + i^2}{2} = \dfrac{-2i}{2} = -i. \]
$Z = -i$ est imaginaire pur, donc $\left(\overrightarrow{AB}\,,\,\overrightarrow{AC'}\right) = \arg(-i) = -\dfrac{\pi}{2} \ [2\pi]$ : le triangle est rectangle en $A$.

\textbf{Longueurs :} $AB = |1 + i| = \sqrt{2}$ et $AC' = |1 - i| = \sqrt{2}$. Donc $AB = AC'$ : le triangle est bien \textbf{isocèle en $A$}.

\textbf{Conclusion :} $ABC'$ est rectangle isocèle en $A$, indirect (angle $-\dfrac{\pi}{2}$). On pouvait aussi remarquer que $|Z| = |-i| = 1$, ce qui donne directement $\dfrac{AC'}{AB} = 1$, donc $AB = AC'$ : le module du quotient fournit le rapport des longueurs, son argument fournit l'angle.
\end{exoresolu}

\begin{popfigure}
\begin{center}
\begin{tikzpicture}[scale=1.5,>=Stealth]
\coordinate (O) at (0,0);
\coordinate (A) at (1,0);
\coordinate (B) at (2,1);
\coordinate (C) at (2,-1);
% axes
\draw[->,popmuted] (-0.4,0) -- (3.2,0) node[below left] {$\mathbb{R}$};
\draw[->,popmuted] (0,-1.8) -- (0,1.8) node[below left] {$i\mathbb{R}$};
% graduations
\foreach \x in {1,2,3} \draw[popmuted] (\x,0.05) -- (\x,-0.05) node[below] {\scriptsize $\x$};
\foreach \y in {-1,1} \draw[popmuted] (0.05,\y) -- (-0.05,\y) node[left] {\scriptsize $\y$};
% triangle
\draw[thick,popblue,fill=popblueL] (A) -- (B) -- (C) -- cycle;
% angle droit en A
\draw[popink] ($(A)+(45:0.28)$) -- ($(A)+(45:0.28)+(-45:0.28)$) -- ($(A)+(-45:0.28)$);
% arc d'angle
\draw[->,popink,thick] ($(A)+(45:0.55)$) arc (45:-45:0.55) node[midway,right,xshift=3pt] {$-\dfrac{\pi}{2}$};
% points
\fill[popdark] (A) circle (1.3pt) node[above left] {$A\,(1)$};
\fill[popdark] (B) circle (1.3pt) node[above right] {$B\,(2+i)$};
\fill[popdark] (C) circle (1.3pt) node[below right] {$C'\,(2-i)$};
% longueurs
\node[popblue] at (1.2,0.75) {\scriptsize $\sqrt{2}$};
\node[popblue] at (1.2,-0.75) {\scriptsize $\sqrt{2}$};
\end{tikzpicture}
\end{center}
\legende{Le triangle $ABC'$ : $\dfrac{z_{C'}-z_A}{z_B-z_A} = -i$, d'où un angle droit en $A$ et $AB = AC'$.}
\end{popfigure}

\begin{methode}[Montrer qu'un triangle est rectangle isocèle en un sommet $A$]
\begin{enumerate}
\item Calculer le quotient $Z = \dfrac{z_C - z_A}{z_B - z_A}$.
\item Si $Z$ est imaginaire pur : angle droit en $A$ (rectangle).
\item Si de plus $|Z| = 1$ : $AC = AB$ (isocèle).
\item Le signe de l'argument ($+\dfrac{\pi}{2}$ ou $-\dfrac{\pi}{2}$) précise si le triangle est direct ou indirect.
\item Cas idéal : $Z = i$ ou $Z = -i$ donne les deux informations d'un coup.
\end{enumerate}
\end{methode}

\begin{exercice}[Alignement, perpendicularité et angle]
Le plan est muni d'un repère orthonormé direct. On donne les points $A(2 + i)$, $B(5 + 3i)$, $C(8 + 5i)$ et $D(4i)$.
\begin{enumerate}
\item Montrer que les points $A$, $B$, $C$ sont alignés et préciser la relation vectorielle entre $\overrightarrow{AB}$ et $\overrightarrow{AC}$.
\item Montrer que la droite $(AD)$ est perpendiculaire à la droite $(AB)$, et que $AD = AB$.
\item En déduire la nature du triangle $ADB$, puis déterminer une mesure de l'angle $\left(\overrightarrow{DA}\,,\,\overrightarrow{DB}\right)$ par le calcul.
\end{enumerate}
\end{exercice}

\begin{corrige}[Exercice : Alignement, perpendicularité et angle]
\textbf{1. Alignement de $A$, $B$, $C$.} On forme le quotient :
\[ \dfrac{z_C - z_A}{z_B - z_A} = \dfrac{(8 + 5i) - (2 + i)}{(5 + 3i) - (2 + i)} = \dfrac{6 + 4i}{3 + 2i} = \dfrac{2(3 + 2i)}{3 + 2i} = 2. \]
Ce quotient est un \textbf{réel} (égal à $2$), donc les points $A$, $B$, $C$ sont alignés. Mieux : $z_C - z_A = 2(z_B - z_A)$ signifie $\overrightarrow{AC} = 2\,\overrightarrow{AB}$ : le point $B$ est le milieu du segment $[AC]$.

\textbf{2. Perpendicularité de $(AD)$ et $(AB)$.} On calcule :
\[ z_D - z_A = 4i - (2 + i) = -2 + 3i \quad ; \quad z_B - z_A = 3 + 2i. \]
Or $i(3 + 2i) = 3i + 2i^2 = -2 + 3i = z_D - z_A$. Donc :
\[ \dfrac{z_D - z_A}{z_B - z_A} = i, \]
qui est un \textbf{imaginaire pur} : les droites $(AD)$ et $(AB)$ sont perpendiculaires en $A$. De plus $|i| = 1$, donc $AD = AB$ (vérification : $AD = |-2 + 3i| = \sqrt{13}$ et $AB = |3 + 2i| = \sqrt{13}$).

\textbf{3. Nature du triangle et angle.} D'après la question 2, le triangle $ADB$ est \textbf{rectangle isocèle en $A$}, direct (car $\arg(i) = +\dfrac{\pi}{2}$). Calculons l'angle demandé :
\[ \left(\overrightarrow{DA}\,,\,\overrightarrow{DB}\right) = \arg\left(\dfrac{z_B - z_D}{z_A - z_D}\right) = \arg\left(\dfrac{5 + 3i - 4i}{2 + i - 4i}\right) = \arg\left(\dfrac{5 - i}{2 - 3i}\right). \]
On simplifie en multipliant par le conjugué du dénominateur :
\[ \dfrac{5 - i}{2 - 3i} = \dfrac{(5 - i)(2 + 3i)}{(2 - 3i)(2 + 3i)} = \dfrac{10 + 15i - 2i - 3i^2}{4 + 9} = \dfrac{13 + 13i}{13} = 1 + i. \]
Or $\arg(1 + i) = \dfrac{\pi}{4} \ [2\pi]$, donc :
\[ \left(\overrightarrow{DA}\,,\,\overrightarrow{DB}\right) = \dfrac{\pi}{4} \ [2\pi]. \]
Ce résultat est cohérent : dans un triangle rectangle isocèle, les angles aigus valent $\dfrac{\pi}{4}$.

\emph{Leçon de rédaction : au Baccalauréat, les affixes sont choisies pour que les quotients « tombent » sur un réel ou un imaginaire pur simple. Si ton calcul aboutit à un quotient quelconque, reprends d'abord le calcul des affixes des vecteurs avant de conclure.}
\end{corrige}

\begin{savaistu}[Pourquoi cette machine est-elle si puissante ?]
L'idée de représenter les nombres complexes comme des points du plan revient au mathématicien norvégien-danois Caspar \textsc{Wessel} (1799), puis à Jean-Robert \textsc{Argand} (1806) — c'est lui qui a laissé son nom au mot « argument ». Cette représentation a transformé des calculs trigonométriques pénibles en simples multiplications de complexes. Aujourd'hui encore, les ingénieurs en télécommunications — y compris ceux qui dimensionnent les réseaux MTN et Orange au Cameroun — manipulent quotidiennement modules et arguments : un signal sinusoïdal est représenté par un complexe dont l'argument est la \emph{phase} et le module l'\emph{amplitude}, et filtrer un signal revient à le multiplier par un complexe bien choisi.
\end{savaistu}

\begin{aretenir}[L'essentiel de la section]
\begin{itemize}
\item $\arg(zz') = \arg(z) + \arg(z') \ [2\pi]$ ; $\arg\left(\dfrac{z}{z'}\right) = \arg(z) - \arg(z') \ [2\pi]$ ; $\arg(z^n) = n\arg(z) \ [2\pi]$ ; $\arg(\bar z) = -\arg(z) \ [2\pi]$ ; $\arg(-z) = \arg(z) + \pi \ [2\pi]$.
\item Formule d'or : $\left(\overrightarrow{AB},\overrightarrow{AC}\right) = \arg\left(\dfrac{z_C - z_A}{z_B - z_A}\right) \ [2\pi]$ — le \emph{dénominateur} est le \emph{premier} vecteur.
\item Quotient réel $\iff$ alignement ; quotient imaginaire pur $\iff$ perpendicularité ; quotient de module $1$ $\iff$ longueurs égales.
\item Le module du quotient donne le rapport des distances, son argument donne l'angle : un seul calcul, deux informations.
\end{itemize}
\end{aretenir}

\coursec{Formule de Moivre, linéarisation et racines n-ièmes}

Dans la section précédente, nous avons établi les règles de calcul sur les arguments : $\arg(zz') = \arg z + \arg z'$ $[2\pi]$, $\arg\left(\dfrac{z}{z'}\right) = \arg z - \arg z'$ $[2\pi]$, et $\arg(z^n) = n\arg z$ $[2\pi]$. Ces formules cachent un trésor : appliquées à un complexe de module $1$, elles livrent une identité remarquable reliant trigonométrie et puissances entières, la \cle{formule de Moivre}. De cette formule découlent deux techniques fondamentales au programme du Baccalauréat : la \cle{linéarisation} des puissances de cosinus et sinus (indispensable pour calculer des intégrales), et la résolution complète des équations $z^n = a$ (les fameuses \cle{racines n-ièmes}), dont les images dessinent de magnifiques polygones réguliers. C'est le voyage que nous entreprenons maintenant.

\courssub{La formule de Moivre}

\textbf{Intuition.} Multiplier par $e^{i\theta}$, c'est tourner d'un angle $\theta$. Multiplier $n$ fois de suite par $e^{i\theta}$, c'est donc tourner $n$ fois de l'angle $\theta$, c'est-à-dire tourner de l'angle $n\theta$. Autrement dit :
\[
\left(e^{i\theta}\right)^n = e^{in\theta}.
\]
En écrivant les deux membres sous forme trigonométrique, on obtient l'une des plus belles formules des mathématiques.

\begin{propriete}[Formule de Moivre]
Pour tout réel $\theta$ et tout entier $n \in \mathbb{Z}$ :
\[
(\cos\theta + i\sin\theta)^n = \cos(n\theta) + i\sin(n\theta).
\]
\end{propriete}

\begin{experience}[Démonstration de la formule de Moivre (exigible)]
\textbf{Étape 1 : récurrence sur $\mathbb{N}$.} Posons, pour $n \in \mathbb{N}$, la propriété $\mathcal{P}_n$ : $(\cos\theta + i\sin\theta)^n = \cos(n\theta) + i\sin(n\theta)$.

\emph{Initialisation.} Pour $n = 0$ : $(\cos\theta + i\sin\theta)^0 = 1$ et $\cos 0 + i\sin 0 = 1$. La propriété est vraie.

\emph{Hérédité.} Supposons $\mathcal{P}_n$ vraie pour un entier $n \geq 0$. Alors :
\begin{align*}
(\cos\theta + i\sin\theta)^{n+1} &= (\cos\theta + i\sin\theta)^n (\cos\theta + i\sin\theta) \\
&= \big(\cos(n\theta) + i\sin(n\theta)\big)(\cos\theta + i\sin\theta) \quad \text{(hypothèse de récurrence)}\\
&= \cos(n\theta)\cos\theta - \sin(n\theta)\sin\theta + i\big(\sin(n\theta)\cos\theta + \cos(n\theta)\sin\theta\big)\\
&= \cos\big((n+1)\theta\big) + i\sin\big((n+1)\theta\big),
\end{align*}
en reconnaissant les formules d'addition $\cos(a+b)$ et $\sin(a+b)$. Donc $\mathcal{P}_{n+1}$ est vraie. Par récurrence, $\mathcal{P}_n$ est vraie pour tout $n \in \mathbb{N}$.

\textbf{Étape 2 : extension à $\mathbb{Z}$.} Soit $n$ un entier strictement négatif ; posons $m = -n \in \mathbb{N}^*$. Comme $\cos\theta + i\sin\theta$ est de module $1$, son inverse est son conjugué :
\[
\frac{1}{\cos\theta + i\sin\theta} = \cos\theta - i\sin\theta = \cos(-\theta) + i\sin(-\theta).
\]
Donc, en appliquant le cas positif à l'angle $-\theta$ :
\[
(\cos\theta + i\sin\theta)^n = \big(\cos(-\theta) + i\sin(-\theta)\big)^m = \cos(-m\theta) + i\sin(-m\theta) = \cos(n\theta) + i\sin(n\theta).
\]
La formule est donc valable pour tout $n \in \mathbb{Z}$.
\end{experience}

\courssub{Application 1 : exprimer $\cos(n\theta)$ et $\sin(n\theta)$ en fonction de $\cos\theta$ et $\sin\theta$}

L'idée est simple : on développe $(\cos\theta + i\sin\theta)^n$ par la \cle{formule du binôme de Newton}, puis on identifie parties réelle et imaginaire avec $\cos(n\theta) + i\sin(n\theta)$.

\begin{exemplebox}[Expression de $\cos(3\theta)$ et $\sin(3\theta)$]
Développons $(\cos\theta + i\sin\theta)^3$ :
\[
(\cos\theta + i\sin\theta)^3 = \cos^3\theta + 3i\cos^2\theta\sin\theta - 3\cos\theta\sin^2\theta - i\sin^3\theta.
\]
En regroupant parties réelle et imaginaire, puis en identifiant avec $\cos(3\theta) + i\sin(3\theta)$ :
\[
\cos(3\theta) = \cos^3\theta - 3\cos\theta\sin^2\theta = 4\cos^3\theta - 3\cos\theta,
\]
\[
\sin(3\theta) = 3\cos^2\theta\sin\theta - \sin^3\theta = 3\sin\theta - 4\sin^3\theta,
\]
où l'on a utilisé $\sin^2\theta = 1 - \cos^2\theta$ et $\cos^2\theta = 1 - \sin^2\theta$.
\end{exemplebox}

\courssub{Application 2 : la linéarisation}

\textbf{Pourquoi linéariser ?} Pour calculer une intégrale comme $\displaystyle\int_0^{\pi} \cos^4 x \, dx$, on ne connaît pas de primitive directe de $\cos^4 x$. En revanche, on sait intégrer $\cos(kx)$. \emph{Linéariser}, c'est transformer une puissance $\cos^n x$ ou $\sin^n x$ en une \textbf{somme} de termes de la forme $\cos(kx)$ et $\sin(kx)$. L'outil magique : les \cle{formules d'Euler} :
\[
\cos x = \frac{e^{ix} + e^{-ix}}{2}, \qquad \sin x = \frac{e^{ix} - e^{-ix}}{2i}.
\]

\begin{methode}[Linéariser $\cos^n x$ ou $\sin^n x$ en 4 étapes]
\begin{enumerate}
\item \textbf{Remplacer} $\cos x$ par $\dfrac{e^{ix}+e^{-ix}}{2}$ (ou $\sin x$ par $\dfrac{e^{ix}-e^{-ix}}{2i}$).
\item \textbf{Développer} la puissance $n$-ième par la formule du binôme de Newton.
\item \textbf{Regrouper} les termes deux à deux : le terme en $e^{ikx}$ avec le terme en $e^{-ikx}$ (ce sont des conjugués).
\item \textbf{Reconnaître} les formules d'Euler : $e^{ikx} + e^{-ikx} = 2\cos(kx)$ et $e^{ikx} - e^{-ikx} = 2i\sin(kx)$.
\end{enumerate}
\end{methode}

\begin{exemplebox}[Linéarisation de $\cos^4 x$]
Appliquons la méthode.
\begin{align*}
\cos^4 x &= \left(\frac{e^{ix}+e^{-ix}}{2}\right)^4 = \frac{1}{16}\left(e^{ix}+e^{-ix}\right)^4 \\
&= \frac{1}{16}\left(e^{4ix} + 4e^{2ix} + 6 + 4e^{-2ix} + e^{-4ix}\right) \\
&= \frac{1}{16}\Big(\big(e^{4ix}+e^{-4ix}\big) + 4\big(e^{2ix}+e^{-2ix}\big) + 6\Big) \\
&= \frac{1}{16}\big(2\cos(4x) + 8\cos(2x) + 6\big).
\end{align*}
D'où le résultat :
\[
\cos^4 x = \frac{1}{8}\cos(4x) + \frac{1}{2}\cos(2x) + \frac{3}{8}.
\]
On peut alors intégrer sans peine : une primitive est $\dfrac{\sin(4x)}{32} + \dfrac{\sin(2x)}{4} + \dfrac{3x}{8}$.
\end{exemplebox}

\begin{attention}[Pièges de la linéarisation]
\begin{itemize}
\item N'oublie pas le dénominateur : pour $\sin^n x$, le facteur est $(2i)^n$, dont le signe dépend de $n$ (par exemple $(2i)^3 = -8i$, $(2i)^4 = 16$).
\item Le regroupement se fait entre termes d'exposants \emph{opposés} : $e^{2ix}$ avec $e^{-2ix}$, pas avec $e^{ix}$.
\item Vérifie ton résultat en une valeur simple : pour $x = 0$, $\cos^4 0 = 1$ et $\frac{1}{8} + \frac{1}{2} + \frac{3}{8} = 1$. C'est un excellent réflexe de contrôle.
\end{itemize}
\end{attention}

\courssub{Racines n-ièmes de l'unité}

\textbf{Question.} Quels sont les nombres complexes dont la puissance $n$-ième vaut $1$ ? Pour $n = 2$, ce sont $1$ et $-1$. Pour $n = 3$, on connaît $1$, mais aussi $j = e^{2i\pi/3}$ et $j^2$. Généralisons.

\begin{definition}[Racines n-ièmes de l'unité]
Soit $n \in \mathbb{N}^*$. On appelle \cle{racine n-ième de l'unité} toute solution complexe de l'équation $z^n = 1$.
\end{definition}

\begin{propriete}[Description des racines n-ièmes de l'unité]
L'équation $z^n = 1$ admet exactement $n$ solutions distinctes :
\[
\omega_k = e^{\frac{2ik\pi}{n}}, \qquad k \in \{0, 1, \ldots, n-1\}.
\]
De plus, leur somme est nulle : $\displaystyle\sum_{k=0}^{n-1} \omega_k = 0$ (pour $n \geq 2$).
\end{propriete}

\begin{experience}[Démonstration guidée : pourquoi exactement $n$ racines ?]
Soit $z$ une solution de $z^n = 1$, écrite sous forme exponentielle $z = re^{i\theta}$ avec $r > 0$.

\textbf{Étape 1 : identification du module.} $z^n = r^n e^{in\theta} = 1 = 1 \cdot e^{i \cdot 0}$. Deux complexes égaux ont même module : $r^n = 1$, donc $r = 1$ (car $r > 0$).

\textbf{Étape 2 : identification des arguments.} Les arguments sont égaux modulo $2\pi$ : $n\theta = 0 \ [2\pi]$, c'est-à-dire $n\theta = 2k\pi$ pour un $k \in \mathbb{Z}$, soit $\theta = \dfrac{2k\pi}{n}$.

\textbf{Étape 3 : comptage.} Deux entiers $k$ et $k'$ donnent la même racine si et seulement si $\dfrac{2k\pi}{n} \equiv \dfrac{2k'\pi}{n} \ [2\pi]$, c'est-à-dire $k \equiv k' \ [n]$. Il y a donc exactement $n$ classes : celles de $k = 0, 1, \ldots, n-1$.

\textbf{Étape 4 : somme nulle.} Pour $n \geq 2$, $\omega_1 = e^{2i\pi/n} \neq 1$, et les racines sont les puissances $\omega_1^k$. C'est une suite géométrique :
\[
\sum_{k=0}^{n-1} \omega_1^k = \frac{1 - \omega_1^n}{1 - \omega_1} = \frac{1-1}{1-\omega_1} = 0.
\]
\end{experience}

\begin{propriete}[Interprétation géométrique]
Les images des $n$ racines n-ièmes de l'unité sont les sommets d'un \cle{polygone régulier} à $n$ côtés, inscrit dans le cercle unité, dont l'un des sommets est le point d'affixe $1$.
\end{propriete}

En effet, toutes les racines sont de module $1$ (donc sur le cercle trigonométrique) et leurs arguments forment une progression arithmétique de raison $\dfrac{2\pi}{n}$ : les points sont régulièrement espacés.

\begin{popfigure}
\begin{center}
\begin{tikzpicture}[scale=2.6]
\draw[popline, thick] (0,0) circle (1);
\draw[->, popmuted] (-1.4,0) -- (1.4,0) node[below right] {$\mathbb{R}$};
\draw[->, popmuted] (0,-1.4) -- (0,1.4) node[above left] {$i\mathbb{R}$};
\foreach \k/\nom in {0/{\omega_0=1}, 1/{\omega_1}, 2/{\omega_2}, 3/{\omega_3}, 4/{\omega_4}}{
  \coordinate (P\k) at ({72*\k}:1);
  \fill[popblue] (P\k) circle (0.035);
}
\draw[poporange, very thick] (P0) -- (P1) -- (P2) -- (P3) -- (P4) -- cycle;
\node[popdark, right] at (P0) {$\omega_0 = 1$};
\node[popdark, above right] at (P1) {$\omega_1 = e^{2i\pi/5}$};
\node[popdark, above left] at (P2) {$\omega_2$};
\node[popdark, below left] at (P3) {$\omega_3$};
\node[popdark, below right] at (P4) {$\omega_4$};
\draw[popgreen, thick, ->] (0.35,0) arc (0:72:0.35);
\node[popgreen] at (36:0.55) {$\frac{2\pi}{5}$};
\fill[popdark] (0,0) circle (0.03) node[below left] {$O$};
\end{tikzpicture}
\end{center}
\legende{Les cinq racines 5-ièmes de l'unité : sommets d'un pentagone régulier inscrit dans le cercle unité.}
\end{popfigure}

\courssub{Racines n-ièmes d'un nombre complexe non nul}

\begin{propriete}[Racines n-ièmes d'un complexe non nul]
Soit $Z = R e^{i\varphi}$ un complexe non nul ($R > 0$) et $n \in \mathbb{N}^*$. L'équation $z^n = Z$ admet exactement $n$ solutions distinctes :
\[
z_k = R^{1/n}\, e^{\,i\frac{\varphi + 2k\pi}{n}}, \qquad k \in \{0, 1, \ldots, n-1\}.
\]
Leurs images forment un polygone régulier à $n$ côtés inscrit dans le cercle de centre $O$ et de rayon $R^{1/n}$.
\end{propriete}

La démonstration suit exactement le même schéma que pour l'unité : si $z = re^{i\theta}$, alors $z^n = r^n e^{in\theta} = Re^{i\varphi}$ impose $r^n = R$ (donc $r = R^{1/n}$) et $n\theta = \varphi \ [2\pi]$, soit $\theta = \dfrac{\varphi + 2k\pi}{n}$, avec $n$ valeurs distinctes pour $k = 0, \ldots, n-1$.

\begin{aretenir}[Lien avec les racines de l'unité et factorisation]
Si $z_0$ est \textbf{une} racine n-ième de $Z$, alors toutes les autres s'obtiennent en multipliant $z_0$ par les racines n-ièmes de l'unité : $z_k = z_0 \, \omega_k$. On en déduit la factorisation :
\[
z^n - Z = \prod_{k=0}^{n-1} (z - z_k).
\]
En particulier : $z^n - 1 = (z-1)(z - \omega_1)\cdots(z-\omega_{n-1})$, et pour $n = 3$ : $z^3 - 1 = (z-1)(z^2 + z + 1)$, ce qui redonne $j$ et $j^2$ comme racines de $z^2+z+1$.
\end{aretenir}

\begin{exoresolu}[Résoudre $z^3 = -8$ et représenter les solutions]
\textbf{Énoncé.} Résoudre dans $\mathbb{C}$ l'équation $z^3 = -8$, puis représenter les images des solutions dans le plan complexe.
\tcblower
\textbf{Solution.} Écrivons $-8$ sous forme exponentielle : $-8 = 8\,e^{i\pi}$. On a $R = 8$, $\varphi = \pi$, $n = 3$, donc $R^{1/3} = 2$ et :
\[
z_k = 2\, e^{\,i\frac{\pi + 2k\pi}{3}}, \qquad k \in \{0, 1, 2\}.
\]
\begin{itemize}
\item $k = 0$ : $z_0 = 2e^{i\pi/3} = 2\left(\dfrac{1}{2} + i\dfrac{\sqrt{3}}{2}\right) = 1 + i\sqrt{3}$ ;
\item $k = 1$ : $z_1 = 2e^{i\pi} = -2$ (la racine réelle évidente : $(-2)^3 = -8$) ;
\item $k = 2$ : $z_2 = 2e^{5i\pi/3} = 2\left(\dfrac{1}{2} - i\dfrac{\sqrt{3}}{2}\right) = 1 - i\sqrt{3}$.
\end{itemize}
Les trois images sont les sommets d'un triangle équilatéral inscrit dans le cercle de centre $O$ et de rayon $2$. Vérification : $(1+i\sqrt{3})^2 = -2 + 2i\sqrt{3}$, puis $(1+i\sqrt{3})^3 = (1+i\sqrt{3})(-2+2i\sqrt{3}) = -2 - 6 = -8$. Correct.
\end{exoresolu}

\begin{exoresolu}[Racines cubiques de $8i$]
\textbf{Énoncé.} Déterminer les racines cubiques de $Z = 8i$ et préciser la nature de la figure formée par leurs images.
\tcblower
\textbf{Solution.} On écrit $8i = 8\,e^{i\pi/2}$. Donc $R^{1/3} = 2$ et :
\[
z_k = 2\, e^{\,i\frac{\pi/2 + 2k\pi}{3}} = 2\, e^{\,i\left(\frac{\pi}{6} + \frac{2k\pi}{3}\right)}, \qquad k \in \{0,1,2\}.
\]
\begin{itemize}
\item $z_0 = 2e^{i\pi/6} = \sqrt{3} + i$ ;
\item $z_1 = 2e^{i(\pi/6 + 2\pi/3)} = 2e^{5i\pi/6} = -\sqrt{3} + i$ ;
\item $z_2 = 2e^{i(\pi/6 + 4\pi/3)} = 2e^{3i\pi/2} = -2i$.
\end{itemize}
Les trois points sont sur le cercle de rayon $2$, séparés angulairement de $\dfrac{2\pi}{3}$ : ils forment un \textbf{triangle équilatéral}.
\end{exoresolu}

\begin{popfigure}
\begin{center}
\begin{tikzpicture}[scale=1.15]
\draw[popline, thick] (0,0) circle (2);
\draw[->, popmuted] (-2.7,0) -- (2.7,0) node[below right] {$\mathbb{R}$};
\draw[->, popmuted] (0,-2.7) -- (0,2.7) node[above left] {$i\mathbb{R}$};
\coordinate (A) at (30:2);
\coordinate (B) at (150:2);
\coordinate (C) at (270:2);
\draw[poporange, very thick] (A) -- (B) -- (C) -- cycle;
\fill[popblue] (A) circle (0.06) node[right, popdark] {$z_0 = \sqrt{3}+i$};
\fill[popblue] (B) circle (0.06) node[left, popdark] {$z_1 = -\sqrt{3}+i$};
\fill[popblue] (C) circle (0.06) node[below, popdark] {$z_2 = -2i$};
\fill[popgold] (0,2) circle (0.05) node[above right, popdark] {$2i$};
\draw[popgreen, thick, dashed] (0,0) -- (0,2);
\draw[popgreen, thick, ->] (0.6,0) arc (0:30:0.6);
\node[popgreen] at (15:0.95) {$\frac{\pi}{6}$};
\fill[popdark] (0,0) circle (0.05) node[below left] {$O$};
\end{tikzpicture}
\end{center}
\legende{Les racines cubiques de $8i$ : triangle équilatéral inscrit dans le cercle de rayon $2$.}
\end{popfigure}

\begin{attention}[Erreurs fréquentes sur les racines n-ièmes]
\begin{itemize}
\item \textbf{Les valeurs de $k$.} On prend $k \in \{0, 1, \ldots, n-1\}$, soit $n$ valeurs à partir de $0$ — et non $k \in \{1, \ldots, n\}$. Prendre $k$ de $1$ à $n$ donne bien les mêmes racines (car $k=n$ redonne $k=0$), mais c'est une source de confusion et de racines écrites en double. Toute famille de $n$ entiers consécutifs convient, mais $0$ à $n-1$ est la convention sûre.
\item \textbf{Racines de $1$ contre racines de $a$.} Les racines n-ièmes de l'\emph{unité} sont sur le cercle de rayon $1$ ; celles de $a = Re^{i\varphi}$ sont sur le cercle de rayon $R^{1/n}$ et le premier angle est $\varphi/n$, pas $0$. Ne mélange pas les deux situations.
\item \textbf{Le module.} C'est $R^{1/n} = \sqrt[n]{R}$, et non $R/n$ ni $R^n$.
\end{itemize}
\end{attention}

\begin{savaistu}[Des racines de l'unité à ton téléphone]
Les racines n-ièmes de l'unité ne sont pas qu'un joli polygone : elles sont le moteur de la \cle{transformée de Fourier rapide} (FFT), l'algorithme qui décompose un signal en ses fréquences. Quand tu écoutes de la musique en MP3, quand tu appliques un filtre sur une photo, quand ton téléphone capte le réseau MTN ou Orange, des millions de calculs mettant en jeu les nombres $e^{2ik\pi/n}$ sont effectués chaque seconde. La symétrie du polygone régulier des racines de l'unité est précisément ce qui rend l'algorithme si rapide. Les mathématiques de cette leçon tournent littéralement dans ta poche !
\end{savaistu}

\begin{exercice}[Pour t'entraîner]
\begin{enumerate}
\item Linéariser $\sin^3 x$ à l'aide des formules d'Euler.
\item Résoudre dans $\mathbb{C}$ l'équation $z^4 = -16$ et représenter les solutions.
\item Montrer que si $\omega = e^{2i\pi/5}$, alors $1 + \omega + \omega^2 + \omega^3 + \omega^4 = 0$, et en déduire la valeur exacte de $\cos\dfrac{2\pi}{5} + \cos\dfrac{4\pi}{5}$.
\end{enumerate}
\end{exercice}

\begin{corrige}[Exercice : pour t'entraîner]
\textbf{1.} $\sin^3 x = \dfrac{1}{(2i)^3}\left(e^{ix}-e^{-ix}\right)^3 = \dfrac{1}{-8i}\left(e^{3ix} - 3e^{ix} + 3e^{-ix} - e^{-3ix}\right) = \dfrac{1}{-8i}\big(2i\sin 3x - 6i\sin x\big)$, d'où $\sin^3 x = \dfrac{3}{4}\sin x - \dfrac{1}{4}\sin 3x$.

\textbf{2.} $-16 = 16e^{i\pi}$, donc $z_k = 2e^{i(\pi + 2k\pi)/4} = 2e^{i(\pi/4 + k\pi/2)}$ pour $k = 0,1,2,3$ : $z_0 = \sqrt{2}+i\sqrt{2}$, $z_1 = -\sqrt{2}+i\sqrt{2}$, $z_2 = -\sqrt{2}-i\sqrt{2}$, $z_3 = \sqrt{2}-i\sqrt{2}$. Les images forment un carré inscrit dans le cercle de rayon $2$.

\textbf{3.} $\omega$ est racine 5-ième de l'unité distincte de $1$, donc la somme des cinq racines est nulle : $1+\omega+\omega^2+\omega^3+\omega^4 = 0$. En regroupant les termes conjugués ($\omega^4 = \overline{\omega}$, $\omega^3 = \overline{\omega^2}$) : $1 + 2\cos\dfrac{2\pi}{5} + 2\cos\dfrac{4\pi}{5} = 0$, donc $\cos\dfrac{2\pi}{5} + \cos\dfrac{4\pi}{5} = -\dfrac{1}{2}$.
\end{corrige}

Tu maîtrises désormais la formule de Moivre et ses deux grandes applications : la linéarisation, qui te servira en calcul intégral, et les racines n-ièmes, qui transforment une équation algébrique en une superbe figure géométrique. Dans la dernière section, nous utiliserons toute la puissance des complexes — modules, arguments, transformations — pour étudier les lignes de niveau et résoudre des problèmes de géométrie plane dignes du Baccalauréat.

\coursec{Lignes de niveau et résolution de problèmes de géométrie}

Dans la section précédente, la formule de Moivre et les racines $n$-ièmes nous ont montré toute la puissance des arguments pour décrire des rotations et des polygones réguliers. Nous allons maintenant renverser la perspective : au lieu de \emph{construire} des points à partir d'angles, nous partons d'une \emph{condition angulaire} et nous cherchons \emph{tous} les points qui la vérifient. C'est l'idée des \cle{lignes de niveau}, qui donne aux nombres complexes leur visage le plus géométrique — et l'un des outils les plus redoutables au Baccalauréat.

\courssub{Lignes de niveau $\arg\left(\dfrac{z-a}{z-b}\right) = \alpha$}

\textbf{Intuition.} Imagine que tu es sur une pirogue sur le fleuve Wouri et que tu observes deux balises lumineuses $A$ et $B$ sur la rive. Tu mesures l'angle sous lequel tu vois ces deux balises : c'est l'angle $(\overrightarrow{MB}, \overrightarrow{MA})$. Les marins savent depuis des siècles que l'ensemble des positions d'où l'on voit deux points sous un \emph{même angle} est un arc de cercle : c'est le fameux \cle{arc capable}. Les nombres complexes traduisent cela en une ligne.

Soient $A$ et $B$ deux points distincts d'affixes $a$ et $b$, et $M$ un point d'affixe $z$, distinct de $A$ et de $B$. On a vu dans la section 4 que :
\[
\arg\left(\frac{z-a}{z-b}\right) = \arg(z-a) - \arg(z-b) = \left(\overrightarrow{MB},\, \overrightarrow{MA}\right) \quad [2\pi].
\]
La condition $\arg\left(\dfrac{z-a}{z-b}\right) = \alpha$ signifie donc simplement : \emph{« l'angle orienté $(\overrightarrow{MB}, \overrightarrow{MA})$ vaut $\alpha$ »}.

\begin{propriete}[Ligne de niveau modulo $2\pi$ : l'arc capable]
Soient $A(a)$ et $B(b)$ deux points distincts et $\alpha$ un réel \textbf{non congru à $0$ modulo $\pi$}. L'ensemble des points $M(z)$ tels que
\[
\arg\left(\frac{z-a}{z-b}\right) = \alpha \quad [2\pi]
\]
est un \cle{arc de cercle} d'extrémités $A$ et $B$ (l'\emph{arc capable}), \textbf{privé des points $A$ et $B$}.
\end{propriete}

\begin{attention}[Les points $A$ et $B$ sont toujours exclus]
C'est l'erreur la plus fréquente au Bac ! En $M = A$, on a $z - a = 0$, et $0$ n'a \textbf{pas d'argument} ; en $M = B$, le quotient $\dfrac{z-a}{z-b}$ n'est \textbf{pas défini}. Donc $A$ et $B$ n'appartiennent \emph{jamais} à la ligne de niveau, même si l'arc ou le cercle passe par eux. Sur ta figure, marque-les avec des points « creux ».
\end{attention}

\textbf{Cas particuliers fondamentaux (modulo $\pi$).} Deux situations reviennent sans cesse dans les sujets ; elles se formulent naturellement modulo $\pi$ car on y raisonne sur des droites (alignement, orthogonalité) :

\begin{itemize}
\item $\arg\left(\dfrac{z-a}{z-b}\right) = 0 \;[\pi]$ : les vecteurs $\overrightarrow{MA}$ et $\overrightarrow{MB}$ sont colinéaires, donc $M$ appartient à la \textbf{droite $(AB)$ privée de $A$ et $B$}. C'est le critère complexe d'\cle{alignement}.
\item $\arg\left(\dfrac{z-a}{z-b}\right) = \dfrac{\pi}{2} \;[\pi]$ : les vecteurs $\overrightarrow{MA}$ et $\overrightarrow{MB}$ sont orthogonaux, donc $M$ appartient au \textbf{cercle de diamètre $[AB]$ privé de $A$ et $B$}. C'est le critère complexe d'\cle{orthogonalité}.
\end{itemize}

\begin{propriete}[{Cercle de diamètre $[AB]$ — démonstration exigible}]
Soient $A(a)$ et $B(b)$ distincts. Alors :
\[
\arg\left(\frac{z-a}{z-b}\right) = \frac{\pi}{2} \quad [\pi] \iff M \text{ appartient au cercle de diamètre } [AB],\ M \neq A,\ M \neq B.
\]
\end{propriete}

\begin{experience}[Démonstration guidée]
Soit $M(z)$ distinct de $A$ et de $B$.
\begin{enumerate}
\item Traduis la condition sur l'argument en orthogonalité : $\arg\left(\dfrac{z-a}{z-b}\right) = \dfrac{\pi}{2}\;[\pi] \iff (\overrightarrow{MB}, \overrightarrow{MA}) = \dfrac{\pi}{2}\;[\pi] \iff \overrightarrow{MA} \perp \overrightarrow{MB}$.
\item Écris l'orthogonalité avec le produit scalaire : $\overrightarrow{MA} \cdot \overrightarrow{MB} = 0$.
\item Introduis le milieu $I$ de $[AB]$ (identité du parallélogramme scalaire) : $\overrightarrow{MA} \cdot \overrightarrow{MB} = MI^2 - \dfrac{AB^2}{4}$.
\item Conclus : $\overrightarrow{MA} \cdot \overrightarrow{MB} = 0 \iff MI = \dfrac{AB}{2}$, c'est-à-dire $M$ sur le cercle de centre $I$ et de rayon $\dfrac{AB}{2}$ : le cercle de diamètre $[AB]$.
\end{enumerate}
\end{experience}

\begin{propriete}[Ligne de niveau modulo $\pi$ : le cercle complet — admis]
Soient $A(a)$ et $B(b)$ distincts et $\alpha$ un réel non multiple de $\pi$. L'ensemble des points $M(z)$ tels que
\[
\arg\left(\frac{z-a}{z-b}\right) = \alpha \quad [\pi]
\]
est un \cle{cercle passant par $A$ et $B$}, privé de $A$ et $B$.
\end{propriete}

Ce résultat est admis ; il se justifie par le \emph{théorème de l'angle inscrit} : travailler modulo $\pi$ revient à raisonner sur les angles de \emph{droites} $(MB, MA)$, et deux points situés de part et d'autre de la corde $[AB]$ donnent des angles orientés de vecteurs opposés (modulo $2\pi$) mais égaux modulo $\pi$. Les deux arcs capables se réunissent alors en un cercle complet.

\begin{methode}[{Construire une ligne de niveau $\arg\left(\frac{z-a}{z-b}\right) = \alpha \;[2\pi]$}]
\begin{enumerate}
\item Identifier $A(a)$ et $B(b)$, les placer, et vérifier la convention : l'angle mesuré est $(\overrightarrow{MB}, \overrightarrow{MA})$.
\item Choisir un point-test simple (par exemple le milieu $I$ de $[AB]$, ou un point sur la médiatrice) pour savoir de quel côté de $(AB)$ se situe l'arc : le signe de $\alpha$ donne le sens (positif $\Rightarrow$ sens trigonométrique de $\overrightarrow{MB}$ vers $\overrightarrow{MA}$).
\item Utiliser le théorème de l'angle inscrit : l'angle au centre vaut $2\alpha$ ; le centre $\Omega$ est sur la médiatrice de $[AB]$, avec $(\overrightarrow{\Omega B}, \overrightarrow{\Omega A}) = 2\alpha\;[2\pi]$.
\item Tracer l'arc d'extrémités $A$ et $B$ passant par le point-test, et \textbf{exclure $A$ et $B$} (points creux).
\end{enumerate}
\end{methode}

\begin{popfigure}
\begin{center}
\begin{tikzpicture}[scale=1.15, >=Stealth]
\coordinate (A) at (-1.6,0);
\coordinate (B) at (1.6,0);
\coordinate (O) at (0,1.6);
\coordinate (M) at (0,3.86); % Point sur l'arc majeur (supérieur)
% Arc capable (arc majeur, lieu pour alpha=pi/4) : angle au centre intercepte = 90 deg (arc mineur inferieur).
% Arc majeur va de B (315 deg) a A (225 deg) en passant par le haut (90 deg). Sens croissant : 315 -> 225+360=585. Delta = 270.
\draw[popblue, very thick] (B) arc[start angle=-45, end angle=225, radius=2.26];
% Arc mineur (intercepte, pointillés) : de A (225) a B (315) en passant par le bas (270). Delta = 90.
\draw[popmuted, dashed, thin] (A) arc[start angle=225, end angle=315, radius=2.26];
% Points
\draw[popdark] (A) circle (1.6pt) node[below left] {$A(a)$};
\draw[popdark] (B) circle (1.6pt) node[below right] {$B(b)$};
\fill[popink] (M) circle (1.8pt) node[above] {$M(z)$};
\fill[poporange] (O) circle (1.5pt) node[left] {$\Omega$};
\draw[popmuted, thin] (O) -- (A) (O) -- (B);
\draw[popgreen, thin] (M) -- (A) (M) -- (B);
% Angle au point M (angle inscrit = 45 deg)
\draw[popgreen, thick, ->] ($(M)+(-0.6,-0.4)$) arc[start angle=210, end angle=330, radius=0.7] node[midway, below, popgreen] {$\alpha=\frac{\pi}{4}$};
% Angle au centre (angle au centre intercepte = 90 deg = 2alpha)
\draw[poporange, thick, ->] ($(O)+(-0.45,-0.45)$) arc[start angle=225, end angle=315, radius=0.64] node[midway, below, poporange] {$2\alpha$};
\draw[popmuted, dashed] (A) -- (B);
\end{tikzpicture}
\end{center}
\legende{Arc capable : ligne de niveau $\arg\left(\frac{z-a}{z-b}\right) = \frac{\pi}{4}\;[2\pi]$. L'angle inscrit $(\overrightarrow{MB},\overrightarrow{MA})$ vaut $\alpha$ (arc majeur bleu), l'angle au centre intercepté vaut $2\alpha$ (arc mineur pointillé). Les points $A$ et $B$ sont exclus.}
\end{popfigure}

\begin{popfigure}
\begin{center}
\begin{tikzpicture}[scale=1.05, >=Stealth]
\coordinate (A) at (-2,0);
\coordinate (B) at (2,0);
\coordinate (I) at (0,0);
\coordinate (M) at (1.2,1.6);
\draw[popblue, very thick] (I) circle (2);
\draw[popdark] (A) circle (1.6pt) node[below left] {$A(a)$};
\draw[popdark] (B) circle (1.6pt) node[below right] {$B(b)$};
\fill[popmuted] (I) circle (1.4pt) node[below] {$I$};
\fill[popink] (M) circle (1.8pt) node[above right] {$M(z)$};
\draw[popgreen, thick] (M) -- (A) (M) -- (B);
% Angle droit
\draw[popgreen] ($(M)+(-0.28,-0.21)$) -- ($(M)+(-0.49,-0.43)$) -- ($(M)+(-0.37,-0.58)$);
\draw[popmuted, thin] (A) -- (B);
\end{tikzpicture}
\end{center}
\legende{Ligne de niveau $\arg\left(\frac{z-a}{z-b}\right) = \frac{\pi}{2}\;[\pi]$ : le cercle de diamètre $[AB]$, privé de $A$ et $B$. En tout point $M$, l'angle $\widehat{AMB}$ est droit.}
\end{popfigure}

\begin{exemplebox}[Exemple chiffré guidé]
Déterminons l'ensemble $(\mathcal{E})$ des points $M(z)$ tels que $\arg\left(\dfrac{z-2}{z+2i}\right) = \dfrac{\pi}{2}\;[\pi]$.
\begin{itemize}
\item \textbf{Lecture des points} : $z - 2 = z - a$ avec $a = 2$, donc $A(2\,;\,0)$ ; $z + 2i = z - b$ avec $b = -2i$, donc $B(0\,;\,-2)$.
\item \textbf{Nature du lieu} : condition $= \frac{\pi}{2}\;[\pi]$, donc $(\mathcal{E})$ est le cercle de diamètre $[AB]$ privé de $A$ et $B$.
\item \textbf{Éléments du cercle} : centre $I$, milieu de $[AB]$ : $I(1\,;\,-1)$ ; rayon $r = \dfrac{AB}{2} = \dfrac{|{-2i} - 2|}{2} = \dfrac{\sqrt{4+4}}{2} = \sqrt{2}$.
\item \textbf{Conclusion} : $(\mathcal{E})$ est le cercle de centre $I(1\,;\,-1)$ et de rayon $\sqrt{2}$, privé des points $A(2\,;\,0)$ et $B(0\,;\,-2)$. Équation cartésienne possible : $(x-1)^2 + (y+1)^2 = 2$.
\end{itemize}
\end{exemplebox}

\begin{savaistu}[L'arc capable en navigation]
Le principe de l'arc capable est utilisé depuis le XVIII\textsuperscript{e} siècle en navigation maritime et aérienne : un navigateur qui mesure, au sextant puis au \emph{goniomètre} radio, l'angle sous lequel il voit deux balises connues (phares, émetteurs) sait qu'il se trouve sur un arc de cercle précis. En croisant les mesures prises sur deux paires de balises, il obtient sa position à l'intersection de deux arcs. C'est l'ancêtre du positionnement par satellites — et une application directe de la ligne de niveau $\arg\left(\frac{z-a}{z-b}\right) = \alpha$ !
\end{savaistu}

\courssub{Résoudre un problème de géométrie plane avec les nombres complexes}

\textbf{Idée directrice.} Les complexes transforment la géométrie en calcul : les \emph{distances} deviennent des \cle{modules}, les \emph{angles} deviennent des \cle{arguments}, les \emph{rotations} deviennent des \emph{multiplications par $e^{i\theta}$}. Un problème de géométrie devient alors une série d'équations que l'on résout méthodiquement.

\begin{methode}[Plan d'action en 4 étapes]
\begin{enumerate}
\item \textbf{Choisir un repère adapté} : placer l'origine sur un point remarquable de la figure (centre, sommet), l'axe des réels sur une droite importante. Souvent, l'énoncé impose le repère : lis-le attentivement.
\item \textbf{Traduire les hypothèses} : égalité de distances $\rightarrow$ égalité de modules ; orthogonalité $\rightarrow$ quotient imaginaire pur ($\arg = \frac{\pi}{2}\;[\pi]$) ; alignement $\rightarrow$ quotient réel ($\arg = 0\;[\pi]$) ; angle donné $\rightarrow$ rotation.
\item \textbf{Calculer} : manipuler les affixes (factoriser, conjuguer, passer en forme exponentielle) jusqu'à la relation cherchée.
\item \textbf{Conclure géométriquement} : revenir à une phrase de géométrie (« le triangle est équilatéral », « les points sont cocycliques »). Un résultat algébrique sans interprétation ne rapporte pas tous les points !
\end{enumerate}
\end{methode}

\begin{propriete}[Complément utile au Bac : écritures complexes des transformations]
Soit $\Omega$ un point d'affixe $\omega$.
\begin{itemize}
\item La \cle{rotation} de centre $\Omega$ et d'angle $\theta$ est l'application qui à $M(z)$ associe $M'(z')$ avec :
\[ z' - \omega = e^{i\theta}(z - \omega). \]
\item L'\cle{homothétie} de centre $\Omega$ et de rapport $k \in \mathbb{R}^*$ s'écrit :
\[ z' - \omega = k\,(z - \omega). \]
\end{itemize}
Ces écritures sont un outil extrêmement efficace pour traiter les configurations du plan (triangles équilatéraux, carrés, similitudes).
\end{propriete}

\textbf{Applications types.} Trois configurations classiques :

\begin{itemize}
\item \textbf{Triangle équilatéral} : $ABC$ est équilatéral direct si et seulement si $C$ est l'image de $B$ par la rotation de centre $A$ et d'angle $\frac{\pi}{3}$, c'est-à-dire $c - a = e^{i\pi/3}(b - a)$. Condition équivalente très utile : $a^2 + b^2 + c^2 = ab + bc + ca$, ou encore $a + bj + cj^2 = 0$ avec $j = e^{2i\pi/3}$ (triangle équilatéral, sens direct selon l'ordre).
\item \textbf{Carré} : $ABCD$ est un carré direct si et seulement si $D$ est l'image de $A$ par la rotation de centre $B$ d'angle $\frac{\pi}{2}$ et $C$ l'image de $B$ par la rotation de centre $D$ d'angle $\frac{\pi}{2}$ — plus simplement : $d - a = i(c - b)$ et $b - a = i(c - d)$, ou diagonales de même milieu avec $c - a = i(d - b)$... En pratique : mêmes milieux pour $[AC]$ et $[BD]$, $|c - a| = |d - b|$, et $\arg\left(\frac{d-b}{c-a}\right) = \frac{\pi}{2}\;[\pi]$.
\item \textbf{Cocyclicité} : quatre points distincts $A, B, C, D$ sont cocycliques ou alignés si et seulement si le birapport $\dfrac{(c-a)(d-b)}{(c-b)(d-a)}$ est réel, c'est-à-dire si $\arg\left(\dfrac{c-a}{c-b}\right) = \arg\left(\dfrac{d-a}{d-b}\right)\;[\pi]$ (égalité d'angles inscrits).
\end{itemize}

\begin{exoresolu}[Triangle équilatéral par rotation]
Le plan est muni d'un repère orthonormé direct. On donne $A(1)$ et $B(2 + i\sqrt{3})$. Déterminer l'affixe du point $C$ tel que $ABC$ soit équilatéral direct, puis vérifier que $AC = AB$.
\tcblower
$ABC$ équilatéral direct signifie que $C$ est l'image de $B$ par la rotation de centre $A$ et d'angle $\dfrac{\pi}{3}$ :
\[
c - a = e^{i\pi/3}(b - a).
\]
Or $e^{i\pi/3} = \dfrac{1}{2} + i\dfrac{\sqrt{3}}{2}$ et $b - a = (2 + i\sqrt{3}) - 1 = 1 + i\sqrt{3}$. Donc :
\begin{align*}
c - 1 &= \left(\frac{1}{2} + i\frac{\sqrt{3}}{2}\right)\left(1 + i\sqrt{3}\right) \\
&= \frac{1}{2} + i\frac{\sqrt{3}}{2} + i\frac{\sqrt{3}}{2} + i^2\frac{3}{2} \\
&= \frac{1}{2} - \frac{3}{2} + i\sqrt{3} = -1 + i\sqrt{3}.
\end{align*}
D'où $c = i\sqrt{3}$, c'est-à-dire $C(0\,;\,\sqrt{3})$.
\medskip

Vérification : $AB = |b - a| = |1 + i\sqrt{3}| = \sqrt{1 + 3} = 2$ et $AC = |c - a| = |{-1} + i\sqrt{3}| = \sqrt{1 + 3} = 2$. De plus $BC = |c - b| = |{-2}| = 2$. Les trois côtés mesurent $2$ : le triangle $ABC$ est bien équilatéral.
\end{exoresolu}

\begin{popfigure}
\begin{center}
\begin{tikzpicture}[scale=1.5, >=Stealth]
\coordinate (A) at (1,0);
\coordinate (B) at (2,1.732);
\coordinate (C) at (0,1.732);
% axes
\draw[popmuted, thin, ->] (-0.6,0) -- (2.7,0) node[below left] {$x$};
\draw[popmuted, thin, ->] (0,-0.4) -- (0,2.3) node[below left] {$y$};
% triangle
\draw[popblue, very thick] (A) -- (B) -- (C) -- cycle;
\fill[popink] (A) circle (1.4pt) node[below right] {$A(1)$};
\fill[popink] (B) circle (1.4pt) node[above right] {$B(2+i\sqrt{3})$};
\fill[popink] (C) circle (1.4pt) node[above left] {$C(i\sqrt{3})$};
% arc de rotation
\draw[poporange, thick, ->] ($(A)+(0.55,0)$) arc[start angle=0, end angle=60, radius=0.55] node[midway, right, poporange] {$\frac{\pi}{3}$};
\draw[popmuted, thin] (A) -- ($(A)+(0.9,0)$);
\end{tikzpicture}
\end{center}
\legende{Construction du triangle équilatéral direct $ABC$ : $C$ est l'image de $B$ par la rotation de centre $A$ et d'angle $\frac{\pi}{3}$, soit $c - a = e^{i\pi/3}(b-a)$.}
\end{popfigure}

\begin{attention}[Pièges de rédaction dans les problèmes de géométrie]
\begin{itemize}
\item Ne confonds pas $\arg\left(\dfrac{z-a}{z-b}\right)$ (angle $(\overrightarrow{MB}, \overrightarrow{MA})$) avec $\arg\left(\dfrac{z-b}{z-a}\right)$ (angle $(\overrightarrow{MA}, \overrightarrow{MB})$) : l'ordre des affixes change le signe de l'angle !
\item Une égalité de modules ne suffit pas à prouver un triangle équilatéral : il faut \emph{trois} côtés égaux, ou un côté égal \emph{plus} un angle de $\frac{\pi}{3}$.
\item Dans une rotation $z' - \omega = e^{i\theta}(z - \omega)$, n'oublie pas de \emph{recentrer} : on soustrait $\omega$ des deux côtés. Écrire $z' = e^{i\theta} z$ n'est correct que si le centre est l'origine.
\item Conclus toujours en langage géométrique et précise les points exclus le cas échéant.
\end{itemize}
\end{attention}

\begin{exercice}[Pour t'entraîner]
Le plan complexe est rapporté à un repère orthonormé direct $(O\,;\,\vec{u}, \vec{v})$. On donne $A(-2)$ et $B(2)$.
\begin{enumerate}
\item Déterminer et représenter l'ensemble $(\mathcal{C})$ des points $M(z)$ tels que $\arg\left(\dfrac{z+2}{z-2}\right) = \dfrac{\pi}{2}\;[\pi]$.
\item Déterminer l'ensemble $(\mathcal{D})$ des points $M(z)$ tels que $\arg\left(\dfrac{z+2}{z-2}\right) = 0\;[\pi]$.
\item Soit $C$ le point d'affixe $2i$. Montrer que $C \in (\mathcal{C})$ et préciser la nature du triangle $ABC$.
\end{enumerate}
\end{exercice}

\begin{corrige}[Exercice : pour t'entraîner]
\begin{enumerate}
\item On lit $a = -2$ (point $A$) et $b = 2$ (point $B$). La condition $\arg\left(\dfrac{z-a}{z-b}\right) = \dfrac{\pi}{2}\;[\pi]$ caractérise le cercle de diamètre $[AB]$ privé de $A$ et $B$. Ici $[AB]$ a pour milieu $O$ et $AB = 4$, donc $(\mathcal{C})$ est le cercle de centre $O$ et de rayon $2$, privé de $A$ et $B$ : c'est le cercle d'équation $|z| = 2$, soit $x^2 + y^2 = 4$.
\item La condition $\arg = 0\;[\pi]$ traduit la colinéarité de $\overrightarrow{MA}$ et $\overrightarrow{MB}$ : $(\mathcal{D})$ est la droite $(AB)$, c'est-à-dire l'axe des réels, privée de $A$ et $B$.
\item $C$ a affixe $2i$ : $|2i| = 2$, donc $C$ est sur le cercle de centre $O$ de rayon $2$, et $C \neq A$, $C \neq B$, donc $C \in (\mathcal{C})$. Par conséquent l'angle $\widehat{ACB}$ est droit : le triangle $ABC$ est \textbf{rectangle en $C$}. (De plus $AC = BC = 2\sqrt{2}$, donc il est même rectangle isocèle.)
\end{enumerate}
\end{corrige}

\begin{aretenir}[L'essentiel de la section]
\begin{itemize}
\item $\arg\left(\dfrac{z-a}{z-b}\right) = \alpha\;[2\pi]$ ($\alpha \not\equiv 0[\pi]$) : arc capable d'extrémités $A$ et $B$, privé de $A$ et $B$.
\item $\arg\left(\dfrac{z-a}{z-b}\right) = \alpha\;[\pi]$ ($\alpha \not\equiv 0[\pi]$) : cercle passant par $A$ et $B$, privé de $A$ et $B$.
\item Cas clés modulo $\pi$ : $\alpha = 0\;[\pi] \rightarrow$ droite $(AB)$ ; $\alpha = \frac{\pi}{2}\;[\pi] \rightarrow$ cercle de diamètre $[AB]$.
\item Rotation : $z' - \omega = e^{i\theta}(z-\omega)$ ; homothétie : $z' - \omega = k(z - \omega)$.
\item Méthode en 4 temps : repère adapté, traduction modules/arguments, calcul, conclusion géométrique.
\end{itemize}
\end{aretenir}

\newpage

\coursec{Activité d'intégration}

\courssub{Objectifs de l'activité}

À l'issue de cette activité d'intégration, tu seras capable de :
\begin{itemize}
\item modéliser une situation concrète de couverture téléphonique par des ensembles de points du plan complexe ;
\item déterminer et représenter des ensembles définis par une condition sur le module : disque $\abs{z-a}\leq\alpha$, médiatrice $\abs{z-a}=\abs{z-b}$, cercle d'Apollonius $\abs{\dfrac{z-a}{z-b}}=\alpha$ ;
\item déterminer et représenter une ligne de niveau définie par une condition sur l'argument : $\arg\!\left(\dfrac{z-a}{z-b}\right)=\alpha\ [\pi]$ ;
\item calculer et placer les racines $n$-ièmes d'un nombre complexe pour résoudre un problème de positionnement régulier ;
\item produire, en groupe, une « carte de couverture » complète rassemblant tous les lieux géométriques obtenus.
\end{itemize}

\courssub{Situation}

La ville de Bafoussam, chef-lieu de la région de l'Ouest, connaît une croissance rapide de son réseau mobile. Un opérateur de télécommunications souhaite optimiser la couverture de deux quartiers et fait appel à ta classe de Terminale pour modéliser le problème.

On assimile la ville au \cle{plan complexe} muni d'un repère orthonormé direct $(O\,;\vec{u},\vec{v})$, l'unité étant le kilomètre. Le centre-ville (marché B) est représenté par l'origine $O$. Deux antennes relais sont déjà installées :
\begin{itemize}
\item l'antenne $A$, située au quartier Tougang, d'affixe $a=-2$ ;
\item l'antenne $B$, située sur les hauteurs de Banengo, d'affixe $b=2i$.
\end{itemize}

L'opérateur pose cinq questions techniques :
\begin{enumerate}
\item L'antenne $A$ émet avec une portée de \SI{3}{km}. Quelle zone de la ville est couverte par cette antenne ?
\item Pour équilibrer le trafic, on cherche la ligne le long de laquelle un téléphone est à égale distance des deux antennes. Où passe cette ligne ?
\item Un téléphone situé en un point $M$ reçoit le signal de $B$ avec une puissance double de celle reçue de $A$, ce qui se traduit par $MB = 2\,MA$. Quel est le lieu des positions possibles de ce téléphone ?
\item Un technicien placé en $M$ effectue des mesures angulaires : il voit les deux antennes sous un angle droit, c'est-à-dire que l'angle $\widehat{AMB}$ est droit. Sur quelle ligne se trouve-t-il ?
\item Pour densifier le réseau, l'opérateur veut installer \textbf{six micro-relais} équidistants les uns des autres, répartis régulièrement sur un cercle de centre $O$ et de rayon \SI{2}{km}, l'un d'eux étant placé sur l'axe des réels positifs. Quelles sont les affixes de ces six relais ?
\end{enumerate}

Chaque groupe devra présenter une \textbf{carte finale} de la ville avec tous les lieux tracés en couleurs différentes.

\courssub{Consignes}

\textbf{Travail en groupes de quatre élèves. Durée : 2 heures. Matériel : papier millimétré, règle, compas, rapporteur, calculatrice.}

\begin{enumerate}
\item \textbf{Zone de couverture de A.} Traduire la phrase « $M$ est couvert par $A$ » par une condition sur $\abs{z-a}$, où $z$ est l'affixe de $M$. Déterminer la nature de cet ensemble et le représenter (hachurer la zone couverte).
\item \textbf{Ligne d'équilibre.} Écrire la condition $MA = MB$ à l'aide des modules. En utilisant la forme algébrique $z=x+iy$, déterminer une équation cartésienne de cet ensemble, reconnaître sa nature et le tracer.
\item \textbf{Puissance double.} Traduire $MB = 2\,MA$ par $\abs{\dfrac{z-a}{z-b}} = \dfrac{1}{2}$. En élevant au carré et en posant $z=x+iy$, montrer que l'ensemble cherché est un cercle dont on déterminera le centre $\Omega$ (affixe $\omega$) et le rayon $R$. Le tracer.
\item \textbf{Angle droit.} Justifier que $\widehat{AMB}$ est droit si et seulement si $\arg\!\left(\dfrac{z-b}{z-a}\right) = \dfrac{\pi}{2}\ [\pi]$. En déduire, à l'aide du théorème de l'angle inscrit (ou par le calcul avec $z=x+iy$), que $M$ appartient au cercle de diamètre $[AB]$, privé de $A$ et $B$. Déterminer le centre et le rayon de ce cercle et le tracer.
\item \textbf{Les six micro-relais.}
\begin{enumerate}
\item Justifier que les affixes cherchées sont les solutions de l'équation $z^6 = 64$.
\item Résoudre cette équation en écrivant $64$ sous forme exponentielle.
\item Donner les six affixes sous forme exponentielle, puis sous forme algébrique, et placer les six relais sur la carte.
\end{enumerate}
\item \textbf{Synthèse.} Regrouper tous les tracés sur une même figure soignée, avec une légende en couleurs, puis rédiger une courte note technique (5 lignes) à destination de l'opérateur : où conseillez-vous d'implanter un commerce qui a besoin d'être couvert par les deux antennes ?
\end{enumerate}

\courssub{Ressources à mobiliser}

\begin{aretenir}[Boîte à outils de l'activité]
\begin{itemize}
\item \textbf{Distance :} si $M(z)$ et $A(a)$, alors $AM = \abs{z-a}$.
\item \textbf{Disque :} $\abs{z-a}\leq\alpha$ est le disque fermé de centre $A(a)$ et de rayon $\alpha$ (avec $\alpha > 0$).
\item \textbf{Médiatrice :} $\abs{z-a}=\abs{z-b}$ est la médiatrice du segment $[AB]$.
\item \textbf{Cercle d'Apollonius :} pour $\alpha>0$, $\alpha\neq 1$, l'ensemble $\abs{\dfrac{z-a}{z-b}}=\alpha$ est un cercle, obtenu en développant $\abs{z-a}^2=\alpha^2\abs{z-b}^2$ ; pour $\alpha = 1$, on retrouve la médiatrice de $[AB]$.
\item \textbf{Arguments et angles :} $\arg\!\left(\dfrac{z-b}{z-a}\right) = \left(\overrightarrow{MA},\overrightarrow{MB}\right)\ [2\pi]$.
\item \textbf{Ligne de niveau :} $\arg\!\left(\dfrac{z-b}{z-a}\right)=\dfrac{\pi}{2}\ [\pi]$ est le cercle de diamètre $[AB]$ privé de $A$ et $B$.
\item \textbf{Racines $n$-ièmes :} les solutions de $z^n = R\,e^{i\theta}$ (avec $R>0$) sont les $z_k = \sqrt[n]{R}\,e^{i\left(\frac{\theta}{n}+\frac{2k\pi}{n}\right)}$, $k\in\{0,1,\dots,n-1\}$ ; ce sont les sommets d'un polygone régulier à $n$ côtés inscrit dans le cercle de centre $O$ et de rayon $\sqrt[n]{R}$.
\end{itemize}
\end{aretenir}

\begin{attention}[Avant de te lancer]
\begin{itemize}
\item Dans $\arg\!\left(\dfrac{z-b}{z-a}\right)$, le \textbf{numérateur} correspond au point $B$ et le \textbf{dénominateur} au point $A$ : l'angle est $\left(\overrightarrow{MA},\overrightarrow{MB}\right)$, dans cet ordre. Inverser numérateur et dénominateur change l'angle en son opposé.
\item Le quotient $\dfrac{z-b}{z-a}$ n'a de sens que pour $z \neq a$ ; et son argument n'est défini que s'il est non nul, donc aussi pour $z \neq b$. Toute ligne de niveau en argument est donc \emph{privée} de $A$ et $B$.
\item Pour élever une égalité de modules au carré, on utilise $\abs{u}^2 = u\,\overline{u}$, ou plus simplement, si $u = x'+iy'$, $\abs{u}^2 = x'^2 + y'^2$.
\end{itemize}
\end{attention}

\courssub{Démarche et corrigé commenté}

\begin{exoresolu}[Le réseau d'antennes de Bafoussam — corrigé complet]
On pose $z = x+iy$ l'affixe du point $M$, avec $a=-2$ et $b=2i$.

\textbf{1) Zone couverte par l'antenne A.}

$M$ est couvert par $A$ si et seulement si $AM \leq 3$, c'est-à-dire $\abs{z-(-2)}\leq 3$, soit
\[ \abs{z+2}\leq 3. \]
C'est le \textbf{disque fermé} de centre $A(-2)$ et de rayon $3$. On trace le cercle de centre $A$ et de rayon $3$, puis on hachure l'intérieur : c'est la zone de couverture. (Le cercle lui-même correspond à l'égalité $\abs{z+2}=3$ : la limite exacte de portée.)

\textbf{2) Ligne d'équilibre (médiatrice).}

$MA = MB \iff \abs{z+2} = \abs{z-2i}$. En élevant au carré avec $z=x+iy$ :
\begin{align*}
\abs{z+2}^2 &= \abs{z-2i}^2\\
(x+2)^2 + y^2 &= x^2 + (y-2)^2\\
x^2 + 4x + 4 + y^2 &= x^2 + y^2 - 4y + 4\\
4x + 4y &= 0 \quad\Longleftrightarrow\quad y = -x.
\end{align*}
L'ensemble cherché est la \textbf{droite d'équation $y=-x$} : c'est bien la médiatrice de $[AB]$. Vérification : le milieu de $[AB]$ a pour affixe $\dfrac{a+b}{2} = \dfrac{-2+2i}{2} = -1+i$, dont les coordonnées $(-1\,;1)$ vérifient $y=-x$. De plus $\overrightarrow{AB}$ a pour affixe $b-a = 2+2i$, de direction la droite $y=x$, orthogonale à $y=-x$. Tout est cohérent.

\textbf{3) Cercle d'Apollonius ($MB = 2\,MA$).}

$MB = 2\,MA \iff \abs{z-2i} = 2\abs{z+2} \iff \abs{\dfrac{z+2}{z-2i}} = \dfrac{1}{2}$. En élevant au carré :
\begin{align*}
\abs{z-2i}^2 &= 4\abs{z+2}^2\\
x^2 + (y-2)^2 &= 4\left[(x+2)^2 + y^2\right]\\
x^2 + y^2 - 4y + 4 &= 4x^2 + 16x + 16 + 4y^2\\
0 &= 3x^2 + 3y^2 + 16x + 4y + 12.
\end{align*}
En divisant par $3$ : $x^2 + y^2 + \dfrac{16}{3}x + \dfrac{4}{3}y + 4 = 0$. Complétons les carrés :
\[ \left(x+\dfrac{8}{3}\right)^2 - \dfrac{64}{9} + \left(y+\dfrac{2}{3}\right)^2 - \dfrac{4}{9} + 4 = 0
\iff \left(x+\dfrac{8}{3}\right)^2 + \left(y+\dfrac{2}{3}\right)^2 = \dfrac{32}{9}. \]
C'est le \textbf{cercle de centre $\Omega$ d'affixe $\omega = -\dfrac{8}{3}-\dfrac{2}{3}i$ et de rayon $R = \sqrt{\dfrac{32}{9}} = \dfrac{4\sqrt{2}}{3} \approx 1{,}89$ km}. On remarque que $\Omega$ est aligné avec $A$ et $B$ : en effet $\omega = \dfrac{4a - b}{3}$, donc $\overrightarrow{B\Omega} = \dfrac{4}{3}\overrightarrow{BA}$ ; c'est une propriété classique du cercle d'Apollonius, dont le centre est toujours sur la droite $(AB)$.

\textbf{4) Le technicien et l'angle droit.}

L'angle $\left(\overrightarrow{MA},\overrightarrow{MB}\right)$ a pour mesure $\arg\!\left(\dfrac{z-b}{z-a}\right)\ [2\pi]$. Dire que $\widehat{AMB}$ est droit signifie :
\[ \arg\!\left(\dfrac{z-b}{z-a}\right) = \dfrac{\pi}{2}\ [\pi]. \]
D'après le cours (ligne de niveau $\alpha = \dfrac{\pi}{2}\ [\pi]$), l'ensemble des points $M$ est le \textbf{cercle de diamètre $[AB]$, privé des points $A$ et $B$} (en $A$, le quotient n'est pas défini ; en $B$, il est nul et son argument n'existe pas).

Son centre est le milieu de $[AB]$, d'affixe $-1+i$, et son rayon est $\dfrac{AB}{2} = \dfrac{\abs{b-a}}{2} = \dfrac{\abs{2+2i}}{2} = \dfrac{2\sqrt{2}}{2} = \sqrt{2}$ km.

\emph{Vérification par le calcul :} le quotient $\dfrac{z-b}{z-a}$ est imaginaire pur non nul si et seulement si $(z-b)\overline{(z-a)}$ est imaginaire pur non nul. Or
\[ (z-b)\overline{(z-a)} = (x+i(y-2))(x+2-iy), \]
dont la partie réelle est $x(x+2)+(y-2)y = x^2+y^2+2x-2y$. Elle est nulle si et seulement si
\[ x^2+y^2+2x-2y=0 \iff (x+1)^2+(y-1)^2=2 : \]
on retrouve bien le cercle de centre $(-1\,;1)$ et de rayon $\sqrt{2}$, privé de $A$ et $B$.

\textbf{5) Les six micro-relais.}

\textbf{a)} Les relais sont sur le cercle de centre $O$ et de rayon $2$, donc $\abs{z}=2$, et régulièrement espacés, donc leurs arguments diffèrent de $\dfrac{2\pi}{6}=\dfrac{\pi}{3}$. En élevant à la puissance $6$ : $z^6$ a pour module $2^6=64$ et pour argument $6\arg z$, multiple de $2\pi$. Les affixes vérifient donc $z^6 = 64$. Réciproquement, toute solution de $z^6=64$ a pour module $\sqrt[6]{64}=2$ et les six solutions sont régulièrement espacées : l'équation donne exactement les positions cherchées.

\textbf{b)} On écrit $64 = 64\,e^{i\cdot 0}$. Les racines sixièmes sont :
\[ z_k = \sqrt[6]{64}\,e^{i\frac{2k\pi}{6}} = 2\,e^{i\frac{k\pi}{3}}, \qquad k\in\{0,1,2,3,4,5\}. \]

\textbf{c)} Formes algébriques :
\begin{center}
\begin{tabularx}{\linewidth}{|c|l|X|}
\hline
\rowcolor{popblueL} $k$ & Forme exponentielle & Forme algébrique \\
\hline
$0$ & $2e^{i\cdot 0}$ & $z_0 = 2$ \\
\hline
$1$ & $2e^{i\pi/3}$ & $z_1 = 1 + i\sqrt{3}$ \\
\hline
$2$ & $2e^{2i\pi/3}$ & $z_2 = -1 + i\sqrt{3}$ \\
\hline
$3$ & $2e^{i\pi}$ & $z_3 = -2$ \\
\hline
$4$ & $2e^{4i\pi/3}$ & $z_4 = -1 - i\sqrt{3}$ \\
\hline
$5$ & $2e^{5i\pi/3}$ & $z_5 = 1 - i\sqrt{3}$ \\
\hline
\end{tabularx}
\end{center}
Les six relais sont les sommets d'un \textbf{hexagone régulier} inscrit dans le cercle de centre $O$ et de rayon $2$. On remarque que le relais $z_3 = -2$ coïncide exactement avec l'antenne $A$ : l'opérateur pourra donc mutualiser ce site !

\textbf{6) Carte finale et note technique.}

Sur la carte : le disque de couverture (hachuré), la médiatrice $y=-x$, le cercle d'Apollonius de centre $\Omega$, le cercle de diamètre $[AB]$, et l'hexagone des six relais. \emph{Note à l'opérateur :} pour être couvert par les deux antennes à la fois, un commerce doit se trouver dans le disque de centre $A$ et de rayon $3$ \textbf{et} du bon côté de la médiatrice ; la zone proche du segment $[AB]$ (intérieure au cercle de diamètre $[AB]$) garantit en outre une excellente qualité de signal des deux antennes.
\tcblower
\textbf{Commentaires méthodologiques.}
\begin{itemize}
\item La clé de chaque question est la \textbf{traduction} : distance $\to$ module, angle $\to$ argument, positionnement régulier $\to$ racines $n$-ièmes.
\item Pour les questions 2, 3 et 4, le passage par la forme algébrique $z=x+iy$ et l'élévation au carré $\abs{z-a}^2 = (z-a)(\overline{z}-\overline{a})$ est la méthode la plus sûre ; elle aboutit toujours à une équation cartésienne que l'on reconnaît (droite ou cercle).
\item \textbf{Piège fréquent (question 4) :} ne pas oublier d'\emph{exclure} les points $A$ et $B$ du cercle, car le quotient $\dfrac{z-b}{z-a}$ n'y est pas défini ou n'y a pas d'argument.
\item \textbf{Piège fréquent (question 5) :} $\sqrt[6]{64} = 64^{1/6} = 2$ (et non $64/6$ !) ; les six arguments sont $\dfrac{2k\pi}{6}$ pour $k=0,\dots,5$, et non $\dfrac{2k\pi}{64}$.
\item \textbf{Rédaction au Bac :} pour un cercle obtenu par équation cartésienne, toujours conclure en donnant explicitement le \emph{centre} (affixe ou coordonnées) et le \emph{rayon}, puis la nature de l'ensemble.
\end{itemize}
\end{exoresolu}

\begin{popfigure}
\begin{center}
\begin{tikzpicture}[scale=0.85,>=Stealth]
\draw[popline] (-4.5,0) -- (3.5,0) node[below left] {\footnotesize axe réel};
\draw[popline] (0,-3.5) -- (0,4) node[below left] {\footnotesize axe imaginaire};
% disque de couverture
\fill[popblue!18] (-2,0) circle (3);
\draw[popblue,thick] (-2,0) circle (3);
% médiatrice y=-x
\draw[poporange,thick,dashed] (-3.2,3.2) -- (2.6,-2.6) node[pos=0.92,above,sloped] {\footnotesize médiatrice};
% cercle diamètre AB
\draw[popgreen,thick] (-1,1) circle (1.414);
% cercle Apollonius
\draw[poppurple,thick] (-2.6667,-0.6667) circle (1.8856);
% hexagone des relais
\draw[popgold,thick] (2,0) -- (1,1.732) -- (-1,1.732) -- (-2,0) -- (-1,-1.732) -- (1,-1.732) -- cycle;
\foreach \p in {(2,0),(1,1.732),(-1,1.732),(-2,0),(-1,-1.732),(1,-1.732)}
  \fill[popgold] \p circle (0.07);
% points
\fill[popink] (-2,0) circle (0.08) node[below left] {\footnotesize $A(-2)$};
\fill[popink] (0,2) circle (0.08) node[above right] {\footnotesize $B(2i)$};
\fill[popink] (0,0) circle (0.08) node[below right] {\footnotesize $O$};
\fill[poppurple] (-2.6667,-0.6667) circle (0.07) node[below] {\footnotesize $\Omega$};
\node[popblue] at (-3.6,2.4) {\footnotesize couverture $A$};
\node[popgreen] at (0.9,2.6) {\footnotesize cercle diamètre $[AB]$};
\node[popgold] at (1.9,-1.9) {\footnotesize relais};
\end{tikzpicture}
\end{center}
\legende{Carte finale de la modélisation : disque de couverture (bleu), médiatrice (orange), cercle d'Apollonius (violet), cercle de diamètre $[AB]$ (vert) et hexagone des six micro-relais (doré).}
\end{popfigure}

\begin{savaistu}[Apollonius de Perge]
Le cercle obtenu à la question 3 porte le nom d'\emph{Apollonius de Perge} (III\textsuperscript{e} siècle av. J.-C.), mathématicien grec célèbre pour son traité des coniques. Il a montré que le lieu des points dont le rapport des distances à deux points fixes est constant (et différent de $1$) est un cercle. Aujourd'hui, ce même calcul est utilisé en \textbf{radiolocalisation} : en comparant les puissances reçues de deux antennes, un téléphone peut estimer sa position sur un cercle d'Apollonius ; avec trois antennes, l'intersection de deux tels cercles le localise presque exactement. C'est l'un des principes de la géolocalisation par les réseaux mobiles, utilisée chaque jour à Bafoussam comme ailleurs !
\end{savaistu}

\courssub{Bilan de l'activité}

Cette activité a permis de mettre en évidence que les nombres complexes constituent un \cle{langage unifié} pour la géométrie plane :
\begin{itemize}
\item une \textbf{condition de distance} se traduit par une \textbf{condition sur un module} : $\abs{z-a}\leq\alpha$ donne un disque, $\abs{z-a}=\abs{z-b}$ une médiatrice, et $\abs{\dfrac{z-a}{z-b}}=\alpha$ un cercle d'Apollonius ;
\item une \textbf{condition d'angle} se traduit par une \textbf{condition sur un argument} : la ligne de niveau $\arg\!\left(\dfrac{z-b}{z-a}\right)=\dfrac{\pi}{2}\ [\pi]$ redonne le théorème classique « angle droit $\iff$ cercle de diamètre $[AB]$ » ;
\item un \textbf{positionnement régulier} se traduit par la recherche de \textbf{racines $n$-ièmes} : les six micro-relais forment un hexagone régulier, image géométrique des solutions de $z^6=64$.
\end{itemize}
Tu as également constaté la puissance du va-et-vient entre forme algébrique (pour obtenir des équations cartésiennes) et forme exponentielle (pour les angles, rotations et racines). C'est exactement cette double lecture qui sera attendue au Baccalauréat : savoir choisir la forme adaptée au problème, justifier chaque traduction, et conclure par une interprétation géométrique claire. Enfin, la démarche de modélisation (situation réelle $\to$ traduction mathématique $\to$ résolution $\to$ retour à la réalité, avec la note technique à l'opérateur) illustre l'approche par les compétences : les mathématiques ne servent pas seulement à calculer, mais à \emph{décider}.

\newpage

\coursec{Méthodes et exercices résolus}

\coursec{Nombres complexes : approche géométrique}

\courssub{Le plan complexe, affixes et interprétation géométrique}

\begin{definition}[Plan complexe et affixes]
On munit le plan d'un repère orthonormé direct $(O\,;\vec{\imath},\vec{\jmath})$. À tout point $M$ du plan, on associe un unique nombre complexe $z$, appelé \cle{affixe} de $M$, tel que $\overrightarrow{OM}=x\vec{\imath}+y\vec{\jmath}$ avec $z=x+iy$ ($x,y\in\mathbb{R}$). Réciproquement, à tout $z\in\mathbb{C}$ correspond un unique point $M$. Le plan muni de cette correspondance est le \cle{plan complexe}.
\end{definition}

\begin{propriete}[Affixe d'un vecteur]
Si $A(a)$ et $B(b)$ sont deux points, le vecteur $\overrightarrow{AB}$ a pour affixe $b-a$. En particulier, $\overrightarrow{OM}$ a pour affixe $z$ si $M(z)$.
\end{propriete}

\begin{definition}[Module et argument]
Pour $z=x+iy \neq 0$ :
\begin{itemize}
\item Le \cle{module} de $z$ est $|z|=\sqrt{x^2+y^2}$. C'est la distance $OM$.
\item Un \cle{argument} de $z$ est une mesure (en radians) de l'angle orienté $(\vec{\imath}\,;\overrightarrow{OM})$. L'ensemble des arguments est $\arg(z)=\{\theta_0+2k\pi \mid k\in\mathbb{Z}\}$.
\item L'\cle{argument principal} est l'unique argument $\theta \in ]-\pi\,;\pi]$. On le note $\mathrm{Arg}(z)$ ou $\arg(z)$ par abus.
\end{itemize}
Pour $z=0$, le module est $0$ et l'argument n'est pas défini.
\end{definition}

\begin{propriete}[Formes trigonométrique et exponentielle]
Tout $z\neq 0$ s'écrit de manière unique :
\[ z = |z|(\cos\theta + i\sin\theta) = |z|\,e^{i\theta} \qquad (\theta = \arg(z)). \]
\emph{Formules d'Euler :} $\cos\theta = \frac{e^{i\theta}+e^{-i\theta}}{2}$, $\sin\theta = \frac{e^{i\theta}-e^{-i\theta}}{2i}$.
\end{propriete}

\begin{propriete}[Propriétés du module et de l'argument]
Pour $z,z'\in\mathbb{C}^*$ :
\begin{align*}
|zz'| &= |z|\,|z'|, & \arg(zz') &= \arg z + \arg z' \ [2\pi],\\
\left|\frac{z}{z'}\right| &= \frac{|z|}{|z'|}, & \arg\left(\frac{z}{z'}\right) &= \arg z - \arg z' \ [2\pi],\\
|z^n| &= |z|^n, & \arg(z^n) &= n\arg z \ [2\pi],\\
|\overline{z}| &= |z|, & \arg(\overline{z}) &= -\arg z \ [2\pi],\\
|-z| &= |z|, & \arg(-z) &= \arg z + \pi \ [2\pi].
\end{align*}
\end{propriete}

\begin{aretenir}[Interprétation géométrique des opérations]
\begin{itemize}
\item \textbf{Addition :} $z+z'$ est l'affixe du quatrième sommet du parallélogramme construit sur $\overrightarrow{OM}$ et $\overrightarrow{OM'}$ (règle du parallélogramme).
\item \textbf{Multiplication par un réel :} homothétie de centre $O$.
\item \textbf{Multiplication par $e^{i\alpha}$ :} rotation de centre $O$ et d'angle $\alpha$.
\item \textbf{Multiplication générale :} $zz'$ correspond à la composée d'une homothétie de rapport $|z'|$ et d'une rotation d'angle $\arg(z')$ (ou l'inverse).
\item \textbf{Conjugaison :} symétrie axiale par rapport à l'axe réel.
\end{itemize}
\end{aretenir}

\begin{exemplebox}[Passage algébrique $\to$ trigonométrique]
Soit $z = -1 + i\sqrt{3}$.
$|z| = \sqrt{1+3}=2$. $z = 2\left(-\frac12 + i\frac{\sqrt3}{2}\right)$.
$\cos\theta=-\frac12$, $\sin\theta=\frac{\sqrt3}{2} \Rightarrow \theta = \frac{2\pi}{3}$ (2e quadrant).
$z = 2e^{i\frac{2\pi}{3}}$.
\end{exemplebox}

\courssub{Savoir-faire 1 : Déterminer et représenter un ensemble de points défini par une condition sur le module}

\begin{methode}[Ensembles de points et conditions sur le module]
Pour déterminer l'ensemble des points $M(z)$ vérifiant une condition sur les modules, on procède ainsi :
\begin{enumerate}
\item \textbf{Traduire géométriquement.} Repérer les points fixes $A(a)$ et $B(b)$ et écrire $|z-a|=AM$, $|z-b|=BM$.
\item \textbf{Identifier la configuration classique :}
\begin{itemize}
\item $|z-a|=\alpha$ ($\alpha>0$) : \cle{cercle} de centre $A$ et de rayon $\alpha$ ;
\item $|z-a|\leq\alpha$ : \cle{disque fermé} de centre $A$ et de rayon $\alpha$ ;
\item $|z-a|=|z-b|$ : \cle{médiatrice} du segment $[AB]$ ;
\item $\left|\dfrac{z-a}{z-b}\right|=\alpha$ avec $\alpha>0$, $\alpha\neq 1$ : \cle{cercle d'Apollonius} (on peut retrouver son centre et son rayon en écrivant $MA=\alpha\, MB$, en élevant au carré et en utilisant $MA^2=(z-a)(\overline{z}-\overline{a})$).
\end{itemize}
\item \textbf{Réduire si nécessaire.} Si la condition porte sur une expression composée (par exemple $|2z-3+i|$), factoriser : $|2z-3+i|=2\left|z-\dfrac{3-i}{2}\right|$.
\item \textbf{Conclure par une phrase} précisant la nature de l'ensemble, son centre/rayon ou sa position, puis \textbf{représenter} dans le plan complexe.
\end{enumerate}
\textbf{Réflexe :} toujours écrire la condition sous la forme $|z - (\text{affixe})|$ avant de conclure ; le signe dans la parenthèse est un moins.
\end{methode}

\begin{exoresolu}[Ensemble de points et cercle d'Apollonius]
Le plan complexe est rapporté au repère orthonormé direct $(O\,;\vec{u},\vec{v})$. On considère les points $A$ et $B$ d'affixes respectives $a=1$ et $b=-2$.
\begin{enumerate}
\item Déterminer et représenter l'ensemble $(\mathcal{E})$ des points $M$ d'affixe $z$ tels que $|z-1|=|z+2|$.
\item Déterminer la nature de l'ensemble $(\mathcal{F})$ des points $M$ d'affixe $z$ tels que $\left|\dfrac{z-1}{z+2}\right|=2$.
\end{enumerate}
\tcblower
\textbf{1.} La condition $|z-1|=|z+2|$ s'écrit $|z-a|=|z-b|$, c'est-à-dire $MA=MB$. L'ensemble des points équidistants de $A$ et $B$ est la \textbf{médiatrice du segment $[AB]$}.

Cherchons son équation. Posons $z=x+iy$ avec $(x,y)\in\mathbb{R}^2$ :
\begin{align*}
|z-1|^2 &= |z+2|^2\\
(x-1)^2+y^2 &= (x+2)^2+y^2\\
x^2-2x+1 &= x^2+4x+4\\
-6x &= 3 \quad\Longleftrightarrow\quad x=-\tfrac{1}{2}.
\end{align*}
$(\mathcal{E})$ est donc la droite d'équation $x=-\dfrac{1}{2}$ : c'est bien la médiatrice de $[AB]$, qui passe par le milieu $I\left(-\dfrac{1}{2}\right)$ de $[AB]$ et qui est verticale (car $[AB]$ est horizontal, $A$ et $B$ étant sur l'axe réel).

\textbf{2.} La condition s'écrit $\dfrac{MA}{MB}=2$, soit $MA=2\,MB$ avec $M\neq B$. Comme le rapport est différent de $1$, on reconnaît un cercle d'Apollonius. Élevons au carré avec $z=x+iy$ :
\begin{align*}
|z-1|^2 &= 4|z+2|^2\\
(x-1)^2+y^2 &= 4\left[(x+2)^2+y^2\right]\\
x^2-2x+1+y^2 &= 4x^2+16x+16+4y^2\\
0 &= 3x^2+18x+15+3y^2\\
x^2+6x+5+y^2 &= 0\\
(x+3)^2-9+5+y^2 &= 0\\
(x+3)^2+y^2 &= 4.
\end{align*}
$(\mathcal{F})$ est donc le \textbf{cercle de centre $\Omega(-3)$ et de rayon $2$}.

\textbf{Conclusion :} $(\mathcal{E})$ est la médiatrice de $[AB]$ (droite d'équation $x=-\tfrac12$) et $(\mathcal{F})$ est le cercle de centre $\Omega$ d'affixe $-3$ et de rayon $2$. On vérifie au passage que $B(-2)$ n'appartient pas à $(\mathcal{F})$ : $(-2+3)^2+0=1\neq 4$, ce qui est cohérent avec l'exclusion de $B$.
\end{exoresolu}

\begin{popfigure}
\begin{tikzpicture}[scale=1.1, >=Stealth]
\draw[popline] (-5.5,0) -- (2.5,0) node[right] {$\mathbb{R}$};
\draw[popline] (0,-3) -- (0,3) node[above] {$i\mathbb{R}$};
% Points A, B
\fill[popblue] (1,0) circle (2pt) node[below right] {$A(1)$};
\fill[popblue] (-2,0) circle (2pt) node[below left] {$B(-2)$};
% Mediatrice x = -0.5
\draw[popgreen, thick, dashed] (-0.5,-3) -- (-0.5,3) node[above] {$(\mathcal{E})$};
\fill[popgreen] (-0.5,0) circle (2pt) node[above right] {$I(-\frac12)$};
% Apollonius Circle center -3, radius 2
\fill[poppurple] (-3,0) circle (2pt) node[below] {$\Omega(-3)$};
\draw[poppurple, thick] (-3,0) circle (2);
\node[poppurple] at (-3, 2.3) {$(\mathcal{F})$};
% Graduations
\foreach \x in {-5,-4,-3,-2,-1,1,2} \draw (\x,0.1) -- (\x,-0.1) node[below] {\small $\x$};
\foreach \y in {-2,-1,1,2} \draw (0.1,\y) -- (-0.1,\y) node[left] {\small $\y i$};
\end{tikzpicture}
\legende{Représentation de la médiatrice $(\mathcal{E})$ et du cercle d'Apollonius $(\mathcal{F})$.}
\end{popfigure}

\courssub{Savoir-faire 2 : Passer de la forme algébrique à la forme trigonométrique ou exponentielle}

\begin{methode}[Forme algébrique $\leftrightarrow$ forme trigonométrique]
\textbf{Sens algébrique $\to$ trigonométrique} (pour $z=a+ib\neq 0$) :
\begin{enumerate}
\item Calculer le \cle{module} : $r=|z|=\sqrt{a^2+b^2}$.
\item Factoriser : $z=r\left(\dfrac{a}{r}+i\dfrac{b}{r}\right)$ et chercher $\theta$ tel que $\cos\theta=\dfrac{a}{r}$ et $\sin\theta=\dfrac{b}{r}$.
\item Identifier $\theta$ à l'aide du \textbf{cercle trigonométrique} (valeurs remarquables : $\dfrac{1}{2}$, $\dfrac{\sqrt2}{2}$, $\dfrac{\sqrt3}{2}$) et des \textbf{signes} de $a$ et $b$ pour choisir le bon quadrant.
\item Conclure : $z=r(\cos\theta+i\sin\theta)=r\,e^{i\theta}$. L'argument principal est celui qui appartient à $]-\pi\,;\pi]$.
\end{enumerate}
\textbf{Sens trigonométrique $\to$ algébrique :} développer $r(\cos\theta+i\sin\theta)$ en remplaçant $\cos\theta$ et $\sin\theta$ par leurs valeurs exactes.

\textbf{Réflexe de vérification :} le module calculé doit être positif, et on peut contrôler que $\left(\dfrac{a}{r}\right)^2+\left(\dfrac{b}{r}\right)^2=1$.
\end{methode}

\begin{exoresolu}[Mise sous forme trigonométrique]
On considère le nombre complexe $z=-\sqrt{3}+i$.
\begin{enumerate}
\item Calculer le module de $z$.
\item Déterminer l'argument principal de $z$, puis écrire $z$ sous forme trigonométrique et sous forme exponentielle.
\item En déduire la forme algébrique de $z^{6}$.
\end{enumerate}
\tcblower
\textbf{1.} $|z|=\sqrt{(-\sqrt3)^2+1^2}=\sqrt{3+1}=\sqrt4=2$.

\textbf{2.} On factorise par le module :
\[ z=2\left(-\frac{\sqrt3}{2}+i\,\frac{1}{2}\right). \]
On cherche $\theta\in\,]-\pi\,;\pi]$ tel que $\cos\theta=-\dfrac{\sqrt3}{2}$ et $\sin\theta=\dfrac{1}{2}$. Sur le cercle trigonométrique, l'angle dont le cosinus vaut $-\dfrac{\sqrt3}{2}$ est $\pm\dfrac{5\pi}{6}$ ; le sinus positif impose $\theta=\dfrac{5\pi}{6}$ (deuxième quadrant, cohérent avec $a<0$ et $b>0$).

Donc :
\[ z=2\left(\cos\frac{5\pi}{6}+i\sin\frac{5\pi}{6}\right)=2\,e^{i\frac{5\pi}{6}}, \qquad \arg(z)=\frac{5\pi}{6}\ [2\pi]. \]

\textbf{3.} Par la formule de Moivre :
\[ z^{6}=2^{6}\left(\cos\frac{6\times5\pi}{6}+i\sin\frac{6\times5\pi}{6}\right)=64\left(\cos 5\pi+i\sin 5\pi\right). \]
Or $\cos 5\pi=\cos\pi=-1$ et $\sin 5\pi=0$, donc :
\[ z^{6}=-64. \]
\textbf{Conclusion :} $z=2e^{i\frac{5\pi}{6}}$ et $z^6=-64$, qui est un réel négatif : le passage par la forme trigonométrique a rendu le calcul immédiat, alors qu'un développement direct de $(-\sqrt3+i)^6$ aurait été pénible.
\end{exoresolu}

\begin{popfigure}
\begin{tikzpicture}[scale=1.5, >=Stealth]
\draw[popline] (-2.5,0) -- (2.5,0) node[right] {$\mathbb{R}$};
\draw[popline] (0,-2.5) -- (0,2.5) node[above] {$i\mathbb{R}$};
\draw[popblueL] (0,0) circle (2);
\draw[popblue, thick, ->] (0,0) -- (-1.732,1) node[midway, above left] {$2$};
\fill[popblue] (-1.732,1) circle (2pt) node[above left] {$z=-\sqrt3+i$};
\draw[poporange, dashed] (-1.732,0) -- (-1.732,1);
\draw[poporange, dashed] (0,0) -- (-1.732,0);
\draw[poporange, ->] (0.5,0) arc (0:150:0.5) node[midway, right] {$\frac{5\pi}{6}$};
\foreach \x in {-2,-1,1,2} \draw (\x,0.05) -- (\x,-0.05) node[below] {\small $\x$};
\foreach \y in {-2,-1,1,2} \draw (0.05,\y) -- (-0.05,\y) node[left] {\small $\y i$};
\end{tikzpicture}
\legende{Représentation de $z=-\sqrt3+i$ : module $2$, argument $\frac{5\pi}{6}$.}
\end{popfigure}

\courssub{Savoir-faire 3 : Utiliser les opérations sur les arguments en géométrie}

\begin{methode}[Arguments et configurations géométriques]
Les formules à connaître parfaitement (pour $z,z'$ non nuls) :
\[ \arg(zz')=\arg z+\arg z'\ [2\pi], \qquad \arg\left(\frac{z}{z'}\right)=\arg z-\arg z'\ [2\pi], \]
\[ \arg(z^n)=n\arg z\ [2\pi], \qquad \arg(\overline{z})=-\arg z\ [2\pi], \qquad \arg(-z)=\arg z+\pi\ [2\pi]. \]
Pour prouver une configuration géométrique avec trois points $A(a)$, $B(b)$, $C(c)$ :
\begin{enumerate}
\item Former le quotient $Z=\dfrac{c-a}{b-a}$ (rapport des affixes des vecteurs $\overrightarrow{AC}$ et $\overrightarrow{AB}$).
\item Calculer module et argument de $Z$ :
\begin{itemize}
\item $|Z|=\dfrac{AC}{AB}$ donne le rapport des longueurs ;
\item $\arg(Z)=\left(\overrightarrow{AB}\,;\overrightarrow{AC}\right)\ [2\pi]$ donne l'angle orienté.
\end{itemize}
\item Conclure selon les cas remarquables :
\begin{itemize}
\item $Z$ réel $\Longleftrightarrow$ $A,B,C$ alignés ;
\item $Z$ imaginaire pur $\Longleftrightarrow$ $(AB)\perp(AC)$ ;
\item $|Z|=1$ et $\arg Z=\pm\dfrac{\pi}{3}$ $\Longleftrightarrow$ $ABC$ équilatéral ;
\item $|Z|=1$ et $\arg Z=\pm\dfrac{\pi}{2}$ $\Longleftrightarrow$ $ABC$ rectangle isocèle en $A$.
\end{itemize}
\end{enumerate}
\textbf{Réflexe :} l'angle $\left(\overrightarrow{AB}\,;\overrightarrow{AC}\right)$ se lit sur $\dfrac{c-a}{b-a}$ : \emph{arrivée moins origine} au numérateur et au dénominateur, avec $A$ comme sommet de l'angle.
\end{methode}

\begin{exoresolu}[Nature d'un triangle et construction d'un carré]
Dans le plan complexe, on considère les points $A$, $B$, $C$ d'affixes respectives :
\[ a=1+i, \qquad b=3+2i, \qquad c=3i. \]
\begin{enumerate}
\item Démontrer que le triangle $ABC$ est rectangle isocèle en $A$.
\item Déterminer l'affixe du point $D$ tel que $ABDC$ soit un carré.
\end{enumerate}
\tcblower
\textbf{1.} Formons le quotient $Z=\dfrac{c-a}{b-a}$ :
\begin{align*}
c-a &= (3i)-(1+i)=-1+2i,\\
b-a &= (3+2i)-(1+i)=2+i,\\
Z &= \frac{-1+2i}{2+i}=\frac{(-1+2i)(2-i)}{(2+i)(2-i)}=\frac{-2+i+4i-2i^2}{5}=\frac{5i}{5}=i.
\end{align*}
Calculons le module : $|Z|=|i|=1$. Donc $\dfrac{AC}{AB}=1$, c'est-à-dire $AC=AB$ : le triangle est \textbf{isocèle en $A$}.

L'argument principal de $Z$ est $\arg(i)=\dfrac{\pi}{2}$. Donc $\left(\overrightarrow{AB}\,;\overrightarrow{AC}\right)=\dfrac{\pi}{2}\ [2\pi]$ : les droites $(AB)$ et $(AC)$ sont \textbf{perpendiculaires}.

Le triangle $ABC$ est donc \textbf{rectangle isocèle en $A$}.

\textbf{2.} Pour que $ABDC$ soit un carré, il suffit que $ABDC$ soit un parallélogramme (ce qui est vrai si $\overrightarrow{AB}=\overrightarrow{CD}$) car nous venons de prouver que $AB=AC$ et $(AB)\perp(AC)$.
La condition $\overrightarrow{AB}=\overrightarrow{CD}$ s'écrit $b-a = d-c$, d'où :
\[ d = b - a + c = (3+2i) - (1+i) + 3i = 2 + 4i. \]

\textbf{Vérification :} $A(1,1)$, $B(3,2)$, $D(2,4)$, $C(0,3)$. $\overrightarrow{AB}=(2,1)$, $\overrightarrow{BD}=(-1,2)$, $\overrightarrow{DC}=(-2,-1)$, $\overrightarrow{CA}=(1,-2)$. Les quatre côtés ont la longueur $\sqrt{5}$ et les produits scalaires consécutifs sont nuls : c'est bien un carré.

\textbf{Conclusion :} Le triangle $ABC$ est rectangle isocèle en $A$ et l'affixe de $D$ est $d=2+4i$.
\end{exoresolu}

\begin{popfigure}
\begin{tikzpicture}[scale=1.2, >=Stealth]
\draw[popline] (-1,0) -- (4,0) node[right] {$\mathbb{R}$};
\draw[popline] (0,-1) -- (0,5) node[above] {$i\mathbb{R}$};
\coordinate (A) at (1,1);
\coordinate (B) at (3,2);
\coordinate (C) at (0,3);
\coordinate (D) at (2,4);
\draw[popblue, thick] (A) -- (B) -- (D) -- (C) -- cycle;
\draw[poporange, dashed] (A) -- (D);
\draw[poporange, dashed] (B) -- (C);
\foreach \p/\label/\pos in {A/$A(1+i)$/below left, B/$B(3+2i)$/right, C/$C(3i)$/left, D/$D(2+4i)$/above right}
  \fill[popblue] (\p) circle (2pt) node[\pos] {\label};
% Right angle mark
\draw[popgreen] (1.2,1) -- (1.2,1.2) -- (1,1.2);
\draw[popgreen] (3.2,2) -- (3,2.2) -- (2.8,2); % approx at B? No, angle at A.
% Angle at A is between AB and AC. AB=(2,1), AC=(-1,2).
% Mark at A
\draw[popgreen] ($(A)!0.3!(B)$) -- ($($(A)!0.3!(B)$)!0.3!-90:($(A)!0.3!(C)$)$) -- ($(A)!0.3!(C)$);
\foreach \x in {1,2,3} \draw (\x,0.05) -- (\x,-0.05) node[below] {\small $\x$};
\foreach \y in {1,2,3,4} \draw (0.05,\y) -- (-0.05,\y) node[left] {\small $\y i$};
\end{tikzpicture}
\legende{Le carré $ABDC$ construit sur le triangle rectangle isocèle $ABC$.}
\end{popfigure}

\courssub{Savoir-faire 4 : Formule de Moivre et linéarisation}

\begin{propriete}[Formule de Moivre]
Pour tout entier $n$ et tout réel $\theta$ :
\[ (\cos\theta + i\sin\theta)^n = \cos(n\theta) + i\sin(n\theta) \quad \text{ou} \quad (e^{i\theta})^n = e^{in\theta}. \]
Cette formule est \textbf{exigible au Bac} (démonstration par récurrence pour $n\in\mathbb{N}$, extension aux négatifs par inverse).
\end{propriete}

\begin{methode}[Linéariser $\cos^n x$ et $\sin^n x$]
Pour \cle{linéariser} (exprimer $\cos^n x$ ou $\sin^n x$ comme combinaison linéaire de $\cos(kx)$ et $\sin(kx)$) :
\begin{enumerate}
\item Partir des \textbf{formules d'Euler} : $\cos x=\dfrac{e^{ix}+e^{-ix}}{2}$ et $\sin x=\dfrac{e^{ix}-e^{-ix}}{2i}$.
\item Élever à la puissance $n$ et développer avec le \textbf{binôme de Newton}.
\item Regrouper les termes deux à deux : $e^{ikx}+e^{-ikx}=2\cos(kx)$ et $e^{ikx}-e^{-ikx}=2i\sin(kx)$.
\item Simplifier et conclure.
\end{enumerate}
\textbf{Réflexes :} le résultat de la linéarisation de $\cos^n x$ ne contient que des cosinus ; celui de $\sin^n x$ ne contient que des cosinus si $n$ est pair, que des sinus si $n$ est impair. La linéarisation sert surtout à calculer des \textbf{intégrales} et des primitives.
\end{methode}

\begin{exoresolu}[Linéarisation de $\cos^3 x$]
\begin{enumerate}
\item Linéariser $\cos^3 x$.
\item En déduire la valeur exacte de $\displaystyle I=\int_0^{\frac{\pi}{6}}\cos^3 x\,\mathrm{d}x$.
\end{enumerate}
\tcblower
\textbf{1.} D'après les formules d'Euler, $\cos x=\dfrac{e^{ix}+e^{-ix}}{2}$. Donc :
\begin{align*}
\cos^3 x &= \left(\frac{e^{ix}+e^{-ix}}{2}\right)^3=\frac{1}{8}\left(e^{ix}+e^{-ix}\right)^3\\
&=\frac{1}{8}\left(e^{3ix}+3e^{2ix}e^{-ix}+3e^{ix}e^{-2ix}+e^{-3ix}\right) \quad \text{(Binôme de Newton)}\\
&=\frac{1}{8}\left(e^{3ix}+3e^{ix}+3e^{-ix}+e^{-3ix}\right)\\
&=\frac{1}{8}\left[\left(e^{3ix}+e^{-3ix}\right)+3\left(e^{ix}+e^{-ix}\right)\right]\\
&=\frac{1}{8}\left[2\cos(3x)+6\cos x\right].
\end{align*}
D'où la formule de linéarisation :
\[ \cos^3 x=\frac{1}{4}\cos(3x)+\frac{3}{4}\cos x. \]

\textbf{2.} On en déduit :
\begin{align*}
I &= \int_0^{\frac{\pi}{6}}\left(\frac14\cos(3x)+\frac34\cos x\right)\mathrm{d}x\\
&=\left[\frac{1}{12}\sin(3x)+\frac34\sin x\right]_0^{\frac{\pi}{6}}\\
&=\frac{1}{12}\sin\frac{\pi}{2}+\frac34\sin\frac{\pi}{6}-0\\
&=\frac{1}{12}\times 1+\frac34\times\frac12=\frac{1}{12}+\frac{3}{8}=\frac{2}{24}+\frac{9}{24}=\frac{11}{24}.
\end{align*}
\textbf{Conclusion :} $\cos^3 x=\dfrac{1}{4}\cos(3x)+\dfrac{3}{4}\cos x$ et $I=\dfrac{11}{24}$.
\end{exoresolu}

\courssub{Savoir-faire 5 : Déterminer les racines n-ièmes d'un nombre complexe}

\begin{methode}[Racines n-ièmes]
Pour résoudre $z^n=a$ avec $a\neq 0$ :
\begin{enumerate}
\item Écrire $a$ sous \textbf{forme exponentielle} : $a=R\,e^{i\varphi}$ (avec $\varphi = \arg(a)$).
\item Poser $z=r\,e^{i\theta}$ : l'équation devient $r^n e^{in\theta}=R e^{i\varphi}$, d'où $r=\sqrt[n]{R}$ (racine $n$-ième réelle positive) et $n\theta=\varphi\ [2\pi]$.
\item Les $n$ solutions sont :
\[ z_k=\sqrt[n]{R}\; e^{i\left(\frac{\varphi}{n}+\frac{2k\pi}{n}\right)}, \qquad k\in\{0,1,\dots,n-1\}. \]
\item \textbf{Interprétation géométrique :} les images des $z_k$ sont les sommets d'un \cle{polygone régulier} à $n$ côtés, inscrit dans le cercle de centre $O$ et de rayon $\sqrt[n]{R}$.
\end{enumerate}
\textbf{Cas particulier fondamental :} les racines $n$-ièmes de l'unité sont $u_k=e^{\frac{2ik\pi}{n}}$, $k=0,\dots,n-1$. Si l'on connaît une solution particulière $z_0$ de $z^n=a$, toutes les solutions sont $z_0 u_k$.
\end{methode}

\begin{exoresolu}[Racines quatrièmes]
\begin{enumerate}
\item Résoudre dans $\mathbb{C}$ l'équation $z^4=-8+8\sqrt{3}\,i$.
\item Représenter les solutions dans le plan complexe et préciser la nature du polygone formé par leurs images.
\end{enumerate}
\tcblower
\textbf{1.} Écrivons $a=-8+8\sqrt3\,i$ sous forme exponentielle.
\[ |a|=\sqrt{64+192}=\sqrt{256}=16, \qquad a=16\left(-\frac12+i\frac{\sqrt3}{2}\right)=16\,e^{i\frac{2\pi}{3}}. \]
Posons $z=re^{i\theta}$ : $z^4=r^4 e^{4i\theta}=16e^{i\frac{2\pi}{3}}$ donne $r^4=16$, soit $r=2$ (car $r>0$), et $4\theta=\dfrac{2\pi}{3}\ [2\pi]$, c'est-à-dire $\theta=\dfrac{\pi}{6}+\dfrac{k\pi}{2}$, $k\in\{0,1,2,3\}$.

Les quatre solutions sont :
\[ z_0=2e^{i\frac{\pi}{6}}=\sqrt3+i, \qquad z_1=2e^{i\frac{2\pi}{3}}=-1+\sqrt3\,i, \]
\[ z_2=2e^{i\frac{7\pi}{6}}=-\sqrt3-i, \qquad z_3=2e^{i\frac{5\pi}{3}}=1-\sqrt3\,i. \]

\textbf{Vérification partielle :} $z_0^2=(\sqrt3+i)^2=3+2\sqrt3\,i-1=2+2\sqrt3\,i$, puis $z_0^4=(2+2\sqrt3\,i)^2=4+8\sqrt3\,i-12=-8+8\sqrt3\,i$. \checkmark

\textbf{2.} Les quatre points $M_0,M_1,M_2,M_3$ sont situés sur le cercle de centre $O$ et de rayon $2$, et leurs arguments diffèrent de $\dfrac{\pi}{2}$ : ce sont les sommets d'un \textbf{carré} inscrit dans ce cercle, dont une diagonale est portée par la droite $(OM_0M_2)$.

\textbf{Conclusion :} $S=\left\{\sqrt3+i\,;\ -1+\sqrt3\,i\,;\ -\sqrt3-i\,;\ 1-\sqrt3\,i\right\}$, et les images des solutions forment un carré de centre $O$.
\end{exoresolu}

\begin{popfigure}
\begin{tikzpicture}[scale=1.3, >=Stealth]
\draw[popline] (-2.5,0) -- (2.5,0) node[right] {$\mathbb{R}$};
\draw[popline] (0,-2.5) -- (0,2.5) node[above] {$i\mathbb{R}$};
\draw[popblueL] (0,0) circle (2);
\foreach \k/\angle/\coord/\label in {
  0/30/{1.732,1}/$z_0=\sqrt3+i$,
  1/120/{-1,1.732}/$z_1=-1+\sqrt3 i$,
  2/210/{-1.732,-1}/$z_2=-\sqrt3-i$,
  3/300/{1,-1.732}/$z_3=1-\sqrt3 i$
} {
  \fill[popblue] (\coord) circle (2pt);
  \draw[popblue, thin] (0,0) -- (\coord);
  \node[popblue] at (1.3*\angle:2.3) {\small \label};
}
\draw[poporange, thick] (1.732,1) -- (-1.732,-1);
\draw[poporange, thick] (-1,1.732) -- (1,-1.732);
\node at (0, -2.2) {Carré inscrit dans le cercle de rayon 2};
\end{tikzpicture}
\legende{Les 4 racines quatrièmes de $-8+8\sqrt3 i$ : sommets d'un carré.}
\end{popfigure}

\courssub{Savoir-faire 6 : Lignes de niveau définies par une condition sur l'argument}

\begin{methode}[Lignes de niveau $\arg\dfrac{z-a}{z-b}=\alpha$]
Pour déterminer l'ensemble des points $M(z)$, distincts de $A(a)$ et $B(b)$, tels que $\arg\dfrac{z-a}{z-b}=\alpha\ [\pi]$ ou $[2\pi]$ :
\begin{enumerate}
\item \textbf{Traduire} : $\arg\dfrac{z-a}{z-b}=\arg(z-a)-\arg(z-b)=\arg(\overrightarrow{AM})-\arg(\overrightarrow{BM})=\left(\overrightarrow{BM}\,;\overrightarrow{AM}\right)=\left(\overrightarrow{MB}\,;\overrightarrow{MA}\right)$. Attention à l'\textbf{ordre} des points !
\item \textbf{Cas modulo $\pi$ :} $\left(\overrightarrow{MB}\,;\overrightarrow{MA}\right)=\alpha\ [\pi]$ signifie que $M$ voit le segment $[AB]$ sous un angle constant (au signe près) : c'est un \textbf{cercle passant par $A$ et $B$}, privé des points $A$ et $B$ (théorème de l'angle inscrit). Cas particulier : $\alpha=0\ [\pi]$ donne la droite $(AB)$ privée de $A$ et $B$ ; $\alpha=\dfrac{\pi}{2}\ [\pi]$ donne le cercle de diamètre $[AB]$ privé de $A$ et $B$.
\item \textbf{Cas modulo $2\pi$ :} l'angle orienté est fixé, donc $M$ appartient à un \textbf{arc de cercle} (un seul des deux arcs) d'extrémités $A$ et $B$, extrémités exclues. Pour savoir lequel, tester un point (par exemple le milieu d'un arc, ou un point d'ordonnée positive).
\item \textbf{Conclure} en précisant toujours que $A$ et $B$ sont exclus, et représenter l'ensemble.
\end{enumerate}
\end{methode}

\begin{exoresolu}[Ligne de niveau]
Le plan complexe est muni du repère orthonormé direct $(O\,;\vec u,\vec v)$. Soient $A$ et $B$ les points d'affixes $a=-1$ et $b=1$. Déterminer et représenter :
\begin{enumerate}
\item l'ensemble $(\mathcal{C})$ des points $M$ d'affixe $z$ tels que $\arg\dfrac{z+1}{z-1}=\dfrac{\pi}{2}\ [\pi]$ ;
\item l'ensemble $(\mathcal{A})$ des points $M$ tels que $\arg\dfrac{z+1}{z-1}=\dfrac{\pi}{2}\ [2\pi]$.
\end{enumerate}
\tcblower
Remarquons d'abord que $\arg\dfrac{z+1}{z-1}=\arg\dfrac{z-a}{z-b}=\left(\overrightarrow{MB}\,;\overrightarrow{MA}\right)$, avec $M\notin\{A,B\}$.

\textbf{1.} La condition $\left(\overrightarrow{MB}\,;\overrightarrow{MA}\right)=\dfrac{\pi}{2}\ [\pi]$ signifie que les droites $(MA)$ et $(MB)$ sont perpendiculaires, c'est-à-dire que le triangle $MAB$ est rectangle en $M$. D'après le théorème de l'angle inscrit (réciproque), $M$ appartient au \textbf{cercle de diamètre $[AB]$}.

Comme $A(-1)$ et $B(1)$, ce cercle a pour centre $O$ et pour rayon $1$ : c'est le \textbf{cercle trigonométrique}, privé des points $A$ et $B$ (en ces points le quotient n'est pas défini ou est nul, $\arg 0$ n'existant pas).

\textbf{2.} Modulo $2\pi$, l'angle orienté est exactement $\dfrac{\pi}{2}$ : il faut choisir l'un des deux demi-cercles. Testons le point $M(i)$ (affixe $i$, sur le cercle) :
\[ \frac{i+1}{i-1}=\frac{(1+i)(-1-i)}{(i-1)(-1-i)}=\frac{(1+i)(-1-i)}{2}=\frac{-1-i-i-i^2}{2}=\frac{-2i}{2}=-i. \]
Or $\arg(-i)=-\dfrac{\pi}{2}\ [2\pi]$, donc $M(i)$ vérifie la condition $-\dfrac{\pi}{2}$, pas $\dfrac{\pi}{2}$. Le point $M(-i)$ donne par conjugaison $\arg(i)=\dfrac{\pi}{2}$ : il convient.

$(\mathcal{A})$ est donc le \textbf{demi-cercle inférieur} (points d'ordonnée strictement négative) du cercle trigonométrique, privé de $A$ et $B$.

\textbf{Conclusion :} $(\mathcal{C})$ est le cercle de centre $O$ et de rayon $1$ privé de $A$ et $B$ ; $(\mathcal{A})$ est le demi-cercle situé au-dessous de l'axe réel, extrémités $A$ et $B$ exclues. La différence entre les conditions $[\pi]$ et $[2\pi]$ est essentielle : elle change un cercle en un demi-cercle.
\end{exoresolu}

\begin{popfigure}
\begin{tikzpicture}[scale=1.5, >=Stealth]
% Cercle complet (C)
\draw[popblue, thick] (0,0) circle (1);
\fill[popblue] (-1,0) circle (2pt) node[left] {$A(-1)$};
\fill[popblue] (1,0) circle (2pt) node[right] {$B(1)$};
\draw[popblue, dashed] (-1,0) -- (0,1) -- (1,0);
\draw[popblue, dashed] (-1,0) -- (0,-1) -- (1,0);
\node[popblue] at (0,1.3) {$(\mathcal{C})$ : cercle entier (mod $\pi$)};
% Demi-cercle (A) - on le décale pour la lisibilité ou on utilise des pointillés différents
% Ici on montre les deux sur la même figure avec styles différents
\begin{scope}[xshift=4cm]
\draw[popline] (-1.5,0) -- (1.5,0) node[right] {$\mathbb{R}$};
\draw[popline] (0,-1.5) -- (0,1.5) node[above] {$i\mathbb{R}$};
\draw[poporange, very thick] (1,0) arc (0:-180:1); % Demi-cercle inférieur
\fill[popblue] (-1,0) circle (2pt) node[left] {$A$};
\fill[popblue] (1,0) circle (2pt) node[right] {$B$};
\draw[poporange, thick] (-1,0) -- (1,0); % Diamètre exclu visuellement mais base
\node[poporange] at (0,-1.3) {$(\mathcal{A})$ : demi-cercle inf. (mod $2\pi$)};
\fill[poporange] (0,-1) circle (2pt) node[below] {$M(-i)$};
\draw[poporange, ->] (0,0) -- (0,-1) node[midway, right] {$1$};
\end{scope}
\end{tikzpicture}
\legende{À gauche : ensemble $(\mathcal{C})$ (cercle complet privé de $A,B$). À droite : ensemble $(\mathcal{A})$ (demi-cercle inférieur privé de $A,B$).}
\end{popfigure}

\begin{attention}[Les 5 erreurs qui coûtent des points au Bac]
\begin{enumerate}
\item \textbf{Oublier d'exclure les points $A$ et $B$} dans les ensembles définis par $\left|\dfrac{z-a}{z-b}\right|=\alpha$ ou $\arg\dfrac{z-a}{z-b}=\alpha$ : le quotient n'existe pas en $B$, et s'annule en $A$ (or $\arg 0$ n'existe pas). Toute conclusion doit mentionner « privé de $A$ et $B$ ».
\item \textbf{Se tromper de quadrant pour l'argument.} De $\cos\theta=\dfrac{a}{r}$ seul, on ne peut pas conclure : il faut utiliser à la fois le cosinus \emph{et} le sinus (ou les signes de $a$ et $b$). Par exemple, $\arg(-1-i)=-\dfrac{3\pi}{4}$ et non $\dfrac{\pi}{4}$.
\item \textbf{Inverser le sens de l'angle orienté.} $\arg\dfrac{c-a}{b-a}=\left(\overrightarrow{AB}\,;\overrightarrow{AC}\right)$ : le dénominateur correspond au premier vecteur. Inverser numérateur et dénominateur change le signe de l'angle et fait échouer toute la démonstration géométrique.
\item \textbf{Oublier la condition $r>0$ ou écrire un module négatif.} Un module est toujours positif ; si l'on trouve $z=-2(\cos\theta+i\sin\theta)$, il faut réécrire $z=2(\cos(\theta+\pi)+i\sin(\theta+\pi))$ pour obtenir une vraie forme trigonométrique.
\item \textbf{Donner une seule racine $n$-ième.} L'équation $z^n=a$ ($a\neq0$) admet exactement $n$ solutions distinctes, obtenues pour $k=0,1,\dots,n-1$. En oublier, ou prendre $k\in\mathbb{Z}$ sans se restreindre (ce qui redonne les mêmes valeurs), est sanctionné. De même, dans la formule de Moivre, ne pas oublier que $(\cos\theta+i\sin\theta)^n=\cos(n\theta)+i\sin(n\theta)$ exige que le module soit $1$.
\end{enumerate}
\end{attention}

\newpage

\coursec{Exercices}

\courssub{Je vérifie mes connaissances}

\begin{exercice}[QCM : une seule bonne réponse]
Pour chaque question, choisir la bonne réponse (aucune justification demandée).
\begin{enumerate}
\item Le point $M$ d'affixe $z = -2 + 3i$ a pour coordonnées :
\begin{enumerate}
\item $(-2\,;\,3)$ \quad \item $(3\,;\,-2)$ \quad \item $(-2\,;\,-3)$ \quad \item $(2\,;\,3)$
\end{enumerate}
\item Soit $z = 2e^{i\frac{\pi}{3}}$. L'argument principal de $z^3$ est :
\begin{enumerate}
\item $\dfrac{\pi}{3}$ \quad \item $\pi$ \quad \item $0$ \quad \item $\dfrac{\pi}{9}$
\end{enumerate}
\item L'ensemble des points $M(z)$ tels que $|z - 2| = 3$ est :
\begin{enumerate}
\item une droite \quad \item un segment \quad \item un cercle de centre $A(2)$ et de rayon $3$ \quad \item un disque
\end{enumerate}
\item Si $z = -1 - i\sqrt{3}$, alors $|z|$ vaut :
\begin{enumerate}
\item $1$ \quad \item $2$ \quad \item $4$ \quad \item $\sqrt{2}$
\end{enumerate}
\end{enumerate}
\end{exercice}

\begin{exercice}[Vrai ou faux — justifier]
Pour chaque affirmation, dire si elle est vraie ou fausse en justifiant soigneusement (une justification est exigée, comme au Baccalauréat).
\begin{enumerate}
\item « Si $|z| = 1$ alors $\bar{z} = \dfrac{1}{z}$. »
\item « L'ensemble des points $M(z)$ tels que $\left|\dfrac{z - i}{z + i}\right| = 1$ est un cercle. »
\item « Pour tout complexe non nul $z$, $\arg(z^2) = 2\arg(z) \;[2\pi]$. »
\item « L'ensemble $|z - 1 + i| \leq 2$ est un cercle. »
\end{enumerate}
\end{exercice}

\begin{exercice}[Définitions et formes]
\begin{enumerate}
\item Donner la définition de l'affixe d'un point et de l'affixe d'un vecteur dans le plan complexe.
\item Rappeler la forme trigonométrique et la forme exponentielle d'un nombre complexe non nul $z$ de module $r$ et d'argument $\theta$.
\item Soient $A$, $B$, $C$, $D$ quatre points d'affixes respectives $z_A$, $z_B$, $z_C$, $z_D$ (avec $A \neq B$ et $C \neq D$). Interpréter géométriquement le module et un argument du quotient $\dfrac{z_D - z_C}{z_B - z_A}$.
\end{enumerate}
\end{exercice}

\begin{exercice}[Calculs directs]
On donne $z_1 = 1 + i$ et $z_2 = \sqrt{3} - i$.
\begin{enumerate}
\item Calculer le module de $z_1$ et celui de $z_2$.
\item Déterminer un argument de $z_1$ et un argument de $z_2$ (on donnera les arguments principaux).
\item Écrire $z_1$ et $z_2$ sous forme exponentielle.
\item Calculer $z_1 z_2$ sous forme algébrique.
\end{enumerate}
\end{exercice}

\courssub{Je m'entraîne}

\begin{exercice}[Formes trigonométrique et exponentielle]
\begin{enumerate}
\item Mettre sous forme trigonométrique puis exponentielle chacun des nombres complexes :
\[ z_1 = 1 - i\sqrt{3} \;; \qquad z_2 = -5i \;; \qquad z_3 = 2\sqrt{3} + 2i \;; \qquad z_4 = -4. \]
\item On donne $z = 2e^{-i\frac{\pi}{6}}$. Écrire $z$ sous forme algébrique.
\item Soit $z = \dfrac{1 + i}{\sqrt{3} - i}$. Déterminer la forme exponentielle de $z$ en utilisant les formes exponentielles du numérateur et du dénominateur, puis en déduire $\cos\left(\dfrac{5\pi}{12}\right)$ et $\sin\left(\dfrac{5\pi}{12}\right)$.
\end{enumerate}
\end{exercice}

\begin{exercice}[Ensembles de points : conditions sur le module]
Le plan est muni d'un repère orthonormé direct $(O\,;\vec{u},\vec{v})$. Déterminer la nature géométrique de chacun des ensembles suivants, préciser ses éléments caractéristiques (centre, rayon, équation cartésienne, etc.), puis le représenter (on utilisera un repère différent pour chaque ensemble).
\begin{enumerate}
\item $\mathcal{E}_1$ : ensemble des points $M(z)$ tels que $|z - 1 + i| = 2$.
\item $\mathcal{E}_2$ : ensemble des points $M(z)$ tels que $|z - 1 + i| \leq 2$.
\item $\mathcal{E}_3$ : ensemble des points $M(z)$ tels que $\left|\dfrac{z - 3}{z + 1}\right| = 1$.
\end{enumerate}
\end{exercice}

\begin{popfigure}
\begin{tikzpicture}[scale=1.2, >=Stealth]
\draw[popline] (-3,0) -- (3,0) node[right] {$x$};
\draw[popline] (0,-3) -- (0,3) node[above] {$y$};
\foreach \x in {-2,-1,1,2} \draw (\x,0.05) -- (\x,-0.05) node[below] {\small $\x$};
\foreach \y in {-2,-1,1,2} \draw (0.05,\y) -- (-0.05,\y) node[left] {\small $\y$};
% Circle center (1,-1) radius 2
\draw[popblue, thick] (1,-1) circle (2);
\fill[popblue] (1,-1) circle (1.5pt) node[below right] {$\Omega(1,-1)$};
\draw[popblue, dashed] (1,-1) -- (3,-1) node[midway, below] {$R=2$};
\node[popblue] at (3,1) {$\mathcal{E}_1$};
\end{tikzpicture}
\legende{Représentation de $\mathcal{E}_1$ : cercle de centre $\Omega(1,-1)$ et de rayon $2$.}
\end{popfigure}

\begin{popfigure}
\begin{tikzpicture}[scale=1.2, >=Stealth]
\draw[popline] (-3,0) -- (3,0) node[right] {$x$};
\draw[popline] (0,-3) -- (0,3) node[above] {$y$};
\foreach \x in {-2,-1,1,2} \draw (\x,0.05) -- (\x,-0.05) node[below] {\small $\x$};
\foreach \y in {-2,-1,1,2} \draw (0.05,\y) -- (-0.05,\y) node[left] {\small $\y$};
% Disk center (1,-1) radius 2
\fill[popblue!15] (1,-1) circle (2);
\draw[popblue, thick] (1,-1) circle (2);
\fill[popblue] (1,-1) circle (1.5pt) node[below right] {$\Omega(1,-1)$};
\node[popblue] at (3,1) {$\mathcal{E}_2$};
\end{tikzpicture}
\legende{Représentation de $\mathcal{E}_2$ : disque fermé de centre $\Omega(1,-1)$ et de rayon $2$ (hachuré).}
\end{popfigure}

\begin{popfigure}
\begin{tikzpicture}[scale=1.2, >=Stealth]
\draw[popline] (-3,0) -- (3,0) node[right] {$x$};
\draw[popline] (0,-3) -- (0,3) node[above] {$y$};
\foreach \x in {-2,-1,1,2} \draw (\x,0.05) -- (\x,-0.05) node[below] {\small $\x$};
\foreach \y in {-2,-1,1,2} \draw (0.05,\y) -- (-0.05,\y) node[left] {\small $\y$};
% Line x=1
\draw[poporange, thick] (1,-2.5) -- (1,2.5);
\fill[poporange] (1,0) circle (1.5pt) node[above right] {$I(1,0)$};
\fill[poporange] (3,0) circle (1.5pt) node[below] {$A(3,0)$};
\fill[poporange] (-1,0) circle (1.5pt) node[below] {$B(-1,0)$};
\node[poporange] at (2.5,2) {$\mathcal{E}_3$};
\end{tikzpicture}
\legende{Représentation de $\mathcal{E}_3$ : médiatrice du segment $[AB]$ (droite $x=1$).}
\end{popfigure}

\begin{exercice}[Linéarisation et primitive]
\begin{enumerate}
\item À l'aide des formules d'Euler et de la formule du binôme de Newton, linéariser $\sin^5 x$, c'est-à-dire montrer que pour tout réel $x$ :
\[ \sin^5 x = \frac{1}{16}\left(\sin 5x - 5\sin 3x + 10\sin x\right). \]
\item En déduire une primitive sur $\mathbb{R}$ de la fonction $x \mapsto \sin^5 x$.
\end{enumerate}
\end{exercice}

\begin{exercice}[Racines quatrièmes et configuration]
\begin{enumerate}
\item Écrire le nombre complexe $-16$ sous forme exponentielle.
\item Résoudre dans $\mathbb{C}$ l'équation $z^4 = -16$. On donnera les solutions sous forme exponentielle puis sous forme algébrique.
\item On note $M_0$, $M_1$, $M_2$, $M_3$ les images des solutions dans le plan complexe, rangées par arguments croissants dans $[0\,;\,2\pi[$. Représenter ces points dans un repère orthonormé (unité : 1 cm).
\item Montrer que le quadrilatère $M_0M_1M_2M_3$ est un carré dont on précisera le centre et la longueur du côté.
\end{enumerate}
\end{exercice}

\begin{popfigure}
\begin{tikzpicture}[scale=1.5, >=Stealth]
\draw[popline] (-2.5,0) -- (2.5,0) node[right] {$x$};
\draw[popline] (0,-2.5) -- (0,2.5) node[above] {$y$};
\foreach \x in {-2,-1,1,2} \draw (\x,0.05) -- (\x,-0.05) node[below] {\small $\x$};
\foreach \y in {-2,-1,1,2} \draw (0.05,\y) -- (-0.05,\y) node[left] {\small $\y$};
% Circle radius 2
\draw[popmuted, dashed] (0,0) circle (2);
% Roots
\coordinate (M0) at (45:2);
\coordinate (M1) at (135:2);
\coordinate (M2) at (225:2);
\coordinate (M3) at (315:2);
\fill[popblue] (M0) circle (2pt) node[above right] {$M_0(\sqrt{2},\sqrt{2})$};
\fill[popblue] (M1) circle (2pt) node[above left] {$M_1(-\sqrt{2},\sqrt{2})$};
\fill[popblue] (M2) circle (2pt) node[below left] {$M_2(-\sqrt{2},-\sqrt{2})$};
\fill[popblue] (M3) circle (2pt) node[below right] {$M_3(\sqrt{2},-\sqrt{2})$};
% Square
\draw[popblue, thick] (M0) -- (M1) -- (M2) -- (M3) -- cycle;
\fill[popdark] (0,0) circle (2pt) node[above left] {$O$};
\end{tikzpicture}
\legende{Les racines quatrièmes de $-16$ forment un carré de centre $O$ et de côté $2\sqrt{2}$.}
\end{popfigure}

\courssub{Je me prépare au Bac}

\begin{exercice}[Type Bac — forme exponentielle et valeurs exactes]
Le plan complexe est muni d'un repère orthonormé direct $(O\,;\vec{u},\vec{v})$. On considère les nombres complexes :
\[ z_1 = 1 + i \qquad \text{et} \qquad z_2 = \sqrt{3} - i. \]
\begin{enumerate}
\item \begin{enumerate}
\item Déterminer le module et un argument de $z_1$, puis écrire $z_1$ sous forme exponentielle.
\item Déterminer le module et un argument de $z_2$, puis écrire $z_2$ sous forme exponentielle.
\end{enumerate}
\item \begin{enumerate}
\item Calculer $z_1 z_2$ sous forme algébrique.
\item À l'aide des résultats de la question 1, écrire $z_1 z_2$ sous forme exponentielle.
\item En déduire les valeurs exactes de $\cos\left(\dfrac{\pi}{12}\right)$ et de $\sin\left(\dfrac{\pi}{12}\right)$.
\end{enumerate}
\item On pose $Z = \dfrac{z_1}{z_2}$.
\begin{enumerate}
\item Écrire $Z$ sous forme exponentielle.
\item En déduire la forme algébrique de $Z^{12}$. Que peut-on dire de $Z^{12}$ ?
\end{enumerate}
\item Application : un technicien de CAMTEL modélise la transformation d'un signal par la multiplication par $z_1$ dans le plan complexe (similitude directe de rapport $\sqrt{2}$ et d'angle $\frac{\pi}{4}$). Déterminer l'image $M'$ du point $M$ d'affixe $z_2$ par cette transformation et placer $M$ et $M'$ dans un repère orthonormé (unité : 2 cm).
\end{enumerate}
\end{exercice}

\begin{exercice}[Type Bac — géométrie, quotient d'affixes et ligne de niveau]
Le plan complexe est rapporté à un repère orthonormé direct $(O\,;\vec{u},\vec{v})$ d'unité graphique 2 cm. On considère les points $A$, $B$ et $C$ d'affixes respectives :
\[ z_A = 1, \qquad z_B = i, \qquad z_C = -1 - i. \]
\begin{enumerate}
\item Placer les points $A$, $B$ et $C$ dans le repère.
\item \begin{enumerate}
\item Calculer le quotient $q = \dfrac{z_C - z_A}{z_B - z_A}$ sous forme algébrique.
\item Déterminer le module et un argument de $q$.
\item En déduire la nature exacte du triangle $ABC$ (on justifiera soigneusement).
\end{enumerate}
\item On considère l'ensemble $\Gamma$ des points $M$ d'affixe $z$, avec $z \neq -i$, tels que :
\[ \arg\left(\frac{z - 1}{z + i}\right) = \frac{\pi}{3} \;[\pi]. \]
\begin{enumerate}
\item Interpréter géométriquement cette condition à l'aide d'un angle de vecteurs.
\item Montrer que $\Gamma$ est un arc de cercle dont on précisera le cercle support (centre et rayon) et les points exclus.
\item Construire $\Gamma$ sur la figure.
\end{enumerate}
\item Déterminer l'ensemble des points $M$ d'affixe $z$ tels que $\left|\dfrac{z - 1}{z + i}\right| = 2$. On précisera la nature de cet ensemble.
\end{enumerate}
\end{exercice}

\begin{popfigure}
\begin{tikzpicture}[scale=1.5, >=Stealth]
\draw[popline] (-2.5,0) -- (2.5,0) node[right] {$x$};
\draw[popline] (0,-2.5) -- (0,2.5) node[above] {$y$};
\foreach \x in {-2,-1,1,2} \draw (\x,0.05) -- (\x,-0.05) node[below] {\small $\x$};
\foreach \y in {-2,-1,1,2} \draw (0.05,\y) -- (-0.05,\y) node[left] {\small $\y$};
% Points A, B, C
\fill[popblue] (1,0) circle (2pt) node[below right] {$A(1)$};
\fill[popblue] (0,1) circle (2pt) node[above left] {$B(i)$};
\fill[popblue] (-1,-1) circle (2pt) node[below left] {$C(-1-i)$};
% Triangle ABC
\draw[popblue, thick] (1,0) -- (0,1) -- (-1,-1) -- cycle;
% Circle for Gamma: arc of circle through A and B with angle pi/3
% Center calculation: chord AB length sqrt(2). Angle pi/3 => R = sqrt(2)/(2 sin(pi/3)) = sqrt(2)/sqrt(3) = sqrt(6)/3 approx 0.816.
% Midpoint of AB: (0.5, 0.5). Perp bisector slope -1.
% Distance from midpoint to center: d = R cos(pi/3) = R/2 = sqrt(6)/6 approx 0.408.
% Two centers on line y = -x + 1.
% Center 1 (above AB): (0.5 - d/sqrt(2), 0.5 + d/sqrt(2)) = (0.5 - 1/√3, 0.5 + 1/√3) approx (0.5-0.577, 0.5+0.577) = (-0.077, 1.077)
% Center 2 (below AB): (0.5 + 1/√3, 0.5 - 1/√3) approx (1.077, -0.077)
% arg((z-1)/(z+i)) = pi/3 [pi]. Vectors MA (1-z) and MB (i-z)? No: (z-1)/(z+i) = (z-A)/(z-B). arg = arg(z-A) - arg(z-B) = angle(BM, AM) = angle(MA, MB).
% Condition angle AMB = pi/3 or 2pi/3.
% Arc not containing the other points? Standard: locus is arc AB where angle AMB = pi/3 (on one side) and 2pi/3 (other side). Since [pi], it's the union of two arcs symmetric wrt AB.
% We draw the circle support.
\coordinate (O1) at ({0.5 - 1/sqrt(3)}, {0.5 + 1/sqrt(3)});
\coordinate (O2) at ({0.5 + 1/sqrt(3)}, {0.5 - 1/sqrt(3)});
\draw[poporange, thick] (O1) circle ({sqrt(max(0,6))/3});
\draw[poporange, thick] (O2) circle ({sqrt(max(0,6))/3});
\fill[poporange] (O1) circle (1.5pt) node[above left] {$\Omega_1$};
\fill[poporange] (O2) circle (1.5pt) node[below right] {$\Omega_2$};
\node[poporange] at (2,2) {$\Gamma$};
\end{tikzpicture}
\legende{Points $A$, $B$, $C$ et cercles supports de $\Gamma$ (arcs d'extrémités $A$ et $B$).}
\end{popfigure}

\begin{exercice}[Type Bac — problème de synthèse : rotations et triangles équilatéraux]
Le plan complexe est muni d'un repère orthonormé direct $(O\,;\vec{u},\vec{v})$. Une coopérative agricole de l'ouest du Cameroun modélise son terrain par un triangle $ABC$ dont les sommets ont pour affixes respectives :
\[ z_A = 0, \qquad z_B = 4, \qquad z_C = 1 + i\sqrt{3}. \]
On construit, à l'extérieur du triangle $ABC$, les triangles équilatéraux $ABD$ et $BCE$ de sens direct, de sommets respectifs $D$ et $E$.
\begin{enumerate}
\item \begin{enumerate}
\item Déterminer le module et un argument de $z_C - z_A$. Quelle est la nature du triangle $ABC$ ?
\item Placer les points $A$, $B$ et $C$ dans un repère (unité : 1 cm).
\end{enumerate}
\item On rappelle que la rotation de centre $\Omega$ d'affixe $\omega$ et d'angle $\dfrac{\pi}{3}$ transforme le point $M(z)$ en $M'(z')$ avec :
\[ z' - \omega = e^{i\frac{\pi}{3}}(z - \omega). \]
\begin{enumerate}
\item En utilisant la rotation de centre $A$ et d'angle $\dfrac{\pi}{3}$, déterminer l'affixe $z_D$ du point $D$.
\item En utilisant la rotation de centre $B$ et d'angle $\dfrac{\pi}{3}$, déterminer l'affixe $z_E$ du point $E$.
\item Placer $D$ et $E$ sur la figure.
\end{enumerate}
\item \begin{enumerate}
\item Calculer les affixes des vecteurs $\overrightarrow{AE}$ et $\overrightarrow{DC}$.
\item Montrer que $AE = DC$ et que les droites $(AE)$ et $(DC)$ forment un angle de $\dfrac{\pi}{3}$.
\end{enumerate}
\item Généralisation : on considère désormais un triangle $ABC$ quelconque et on construit extérieurement les triangles équilatéraux directs $ABD$ et $BCE$. À l'aide des rotations complexes, montrer que l'on a toujours $AE = DC$, quelle que soit la forme du triangle $ABC$.
\end{enumerate}
\end{exercice}

\begin{popfigure}
\begin{tikzpicture}[scale=1.2, >=Stealth]
\draw[popline] (-1,0) -- (5,0) node[right] {$x$};
\draw[popline] (0,-2) -- (0,4) node[above] {$y$};
\foreach \x in {1,2,3,4} \draw (\x,0.05) -- (\x,-0.05) node[below] {\small $\x$};
\foreach \y in {1,2,3} \draw (0.05,\y) -- (-0.05,\y) node[left] {\small $\y$};
% Triangle ABC
\coordinate (A) at (0,0);
\coordinate (B) at (4,0);
\coordinate (C) at (1,{sqrt(3)});
\fill[popblue] (A) circle (2pt) node[below left] {$A(0)$};
\fill[popblue] (B) circle (2pt) node[below right] {$B(4)$};
\fill[popblue] (C) circle (2pt) node[above] {$C(1,\sqrt{3})$};
\draw[popblue, thick] (A) -- (B) -- (C) -- cycle;
% Point D: rotation of B around A by +pi/3
\coordinate (D) at (60:4); % (2, 2sqrt(3))
\fill[popgreen] (D) circle (2pt) node[above right] {$D(2,2\sqrt{3})$};
% Point E: rotation of C around B by +pi/3
% C-B = (-3, sqrt(3)) = 2sqrt(3) e^{i 150deg}
% Rotate by 60deg -> 2sqrt(3) e^{i 210deg} = 2sqrt(3)(-sqrt(3)/2 - i/2) = (-3, -sqrt(3))
% E = B + (-3, -sqrt(3)) = (1, -sqrt(3))
\coordinate (E) at (1,{-sqrt(3)});
\fill[popgreen] (E) circle (2pt) node[below] {$E(1,-\sqrt{3})$};
% Equilateral triangles
\draw[popgreen, thick] (A) -- (B) -- (D) -- cycle;
\draw[popgreen, thick] (B) -- (C) -- (E) -- cycle;
% Segments AE and DC
\draw[poporange, thick, dashed] (A) -- (E);
\draw[poporange, thick, dashed] (D) -- (C);
\fill[poporange] (A) circle (1pt); \fill[poporange] (E) circle (1pt);
\fill[poporange] (D) circle (1pt); \fill[poporange] (C) circle (1pt);
\node[poporange] at (3,3) {$AE = DC = 2$};
\node[poporange] at (3,-1.5) {$\angle((AE),(DC)) = \frac{\pi}{3}$};
\end{tikzpicture}
\legende{Configuration du problème de synthèse : triangles équilatéraux extérieurs $ABD$ et $BCE$.}
\end{popfigure}

\newpage

\coursec{Corrigés des exercices}

\begin{corrige}[Exercice 1 — QCM : une seule bonne réponse]
\begin{enumerate}
\item \textbf{Réponse a : $(-2\,;\,3)$.} Par définition, le point $M$ d'affixe $z = x + iy$ a pour coordonnées $(x\,;\,y)$. Ici $x = -2$ et $y = 3$.
\item \textbf{Réponse b : $\pi$.} On a $z^3 = \left(2e^{i\frac{\pi}{3}}\right)^3 = 2^3 e^{i \times 3 \times \frac{\pi}{3}} = 8e^{i\pi}$. Comme $\pi \in \left]-\pi\,;\,\pi\right]$, l'argument principal de $z^3$ est $\pi$.
\item \textbf{Réponse c.} $|z - 2| = 3$ se traduit par $AM = 3$ où $A$ est le point d'affixe $2$ : c'est le cercle de centre $A(2\,;\,0)$ et de rayon $3$.
\item \textbf{Réponse b : $2$.} $|z| = \sqrt{(-1)^2 + (-\sqrt{3})^2} = \sqrt{1 + 3} = \sqrt{4} = 2$.
\end{enumerate}
\end{corrige}

\begin{corrige}[Exercice 2 — Vrai ou faux — justifier]
\begin{enumerate}
\item \textbf{Vrai.} Si $|z| = 1$, alors $z\bar{z} = |z|^2 = 1$, donc $z \neq 0$ et, en divisant par $z$, $\bar{z} = \dfrac{1}{z}$.
\item \textbf{Faux.} Posons $A(i)$ et $B(-i)$. La condition s'écrit $\dfrac{|z - i|}{|z + i|} = 1$, c'est-à-dire $AM = BM$. L'ensemble des points équidistants de $A$ et $B$ est la \textbf{médiatrice} du segment $[AB]$ (ici l'axe des réels), et non un cercle. Un cercle (cercle d'Apollonius) n'apparaît que lorsque le rapport vaut $\alpha \neq 1$.
\item \textbf{Vrai.} C'est la propriété de l'argument d'une puissance : pour $z \neq 0$ et $n \in \mathbb{N}^*$, $\arg(z^n) = n\arg(z)\;[2\pi]$, donc avec $n = 2$ : $\arg(z^2) = 2\arg(z)\;[2\pi]$.
\item \textbf{Faux.} $|z - (1 - i)| \leq 2$ se traduit par $\Omega M \leq 2$ avec $\Omega(1\,;\,-1)$ : c'est le \textbf{disque fermé} de centre $\Omega$ et de rayon $2$ (cercle et intérieur), et non le cercle, qui correspondrait à l'égalité $|z - 1 + i| = 2$.
\end{enumerate}
\end{corrige}

\begin{corrige}[Exercice 3 — Définitions et formes]
\begin{enumerate}
\item Dans le plan muni d'un repère orthonormé direct $(O\,;\vec{u},\vec{v})$, l'\cle{affixe} du point $M(x\,;\,y)$ est le nombre complexe $z_M = x + iy$. L'\cle{affixe} du vecteur $\vec{w} = a\vec{u} + b\vec{v}$ est le nombre complexe $z_{\vec{w}} = a + ib$. En particulier, l'affixe de $\overrightarrow{AB}$ est $z_B - z_A$.
\item Si $z$ est un complexe non nul de module $r$ et d'argument $\theta$ :
\[ \text{Forme trigonométrique : } z = r(\cos\theta + i\sin\theta) \;; \qquad \text{forme exponentielle : } z = re^{i\theta}. \]
\item $\left|\dfrac{z_D - z_C}{z_B - z_A}\right| = \dfrac{|z_D - z_C|}{|z_B - z_A|} = \dfrac{CD}{AB}$ : c'est le \textbf{rapport des longueurs} $CD$ et $AB$. Et :
\[ \arg\left(\frac{z_D - z_C}{z_B - z_A}\right) = \arg(z_D - z_C) - \arg(z_B - z_A) = \left(\overrightarrow{AB}\,,\,\overrightarrow{CD}\right)\;[2\pi] : \]
c'est une \textbf{mesure de l'angle orienté} des vecteurs $\overrightarrow{AB}$ et $\overrightarrow{CD}$. En particulier, le quotient est réel si et seulement si $(AB) \parallel (CD)$, et imaginaire pur si et seulement si $(AB) \perp (CD)$.
\end{enumerate}
\end{corrige}

\begin{corrige}[Exercice 4 — Calculs directs]
\begin{enumerate}
\item $|z_1| = \sqrt{1^2 + 1^2} = \sqrt{2}$ ; $|z_2| = \sqrt{(\sqrt{3})^2 + (-1)^2} = \sqrt{3 + 1} = 2$.
\item Pour $z_1$ : on factorise par le module : $z_1 = \sqrt{2}\left(\dfrac{\sqrt{2}}{2} + i\dfrac{\sqrt{2}}{2}\right)$, donc $\cos\theta_1 = \dfrac{\sqrt{2}}{2}$ et $\sin\theta_1 = \dfrac{\sqrt{2}}{2}$, d'où $\arg(z_1) = \dfrac{\pi}{4}\;[2\pi]$.

Pour $z_2$ : $z_2 = 2\left(\dfrac{\sqrt{3}}{2} - \dfrac{i}{2}\right)$, donc $\cos\theta_2 = \dfrac{\sqrt{3}}{2}$ et $\sin\theta_2 = -\dfrac{1}{2}$, d'où $\arg(z_2) = -\dfrac{\pi}{6}\;[2\pi]$.
\item $z_1 = \sqrt{2}\,e^{i\frac{\pi}{4}}$ et $z_2 = 2e^{-i\frac{\pi}{6}}$.
\item $z_1 z_2 = (1 + i)(\sqrt{3} - i) = \sqrt{3} - i + i\sqrt{3} - i^2 = (\sqrt{3} + 1) + i(\sqrt{3} - 1)$.
\end{enumerate}
\end{corrige}

\begin{corrige}[Exercice 5 — Formes trigonométrique et exponentielle]
\begin{enumerate}
\item \textbf{Pour $z_1 = 1 - i\sqrt{3}$ :} $|z_1| = \sqrt{1 + 3} = 2$ ; $z_1 = 2\left(\dfrac{1}{2} - i\dfrac{\sqrt{3}}{2}\right) = 2\left(\cos\left(-\dfrac{\pi}{3}\right) + i\sin\left(-\dfrac{\pi}{3}\right)\right) = 2e^{-i\frac{\pi}{3}}$.

\textbf{Pour $z_2 = -5i$ :} $|z_2| = 5$ ; $z_2 = 5\left(\cos\left(-\dfrac{\pi}{2}\right) + i\sin\left(-\dfrac{\pi}{2}\right)\right) = 5e^{-i\frac{\pi}{2}}$.

\textbf{Pour $z_3 = 2\sqrt{3} + 2i$ :} $|z_3| = \sqrt{12 + 4} = 4$ ; $z_3 = 4\left(\dfrac{\sqrt{3}}{2} + i\dfrac{1}{2}\right) = 4\left(\cos\dfrac{\pi}{6} + i\sin\dfrac{\pi}{6}\right) = 4e^{i\frac{\pi}{6}}$.

\textbf{Pour $z_4 = -4$ :} $|z_4| = 4$ ; $z_4 = 4(\cos\pi + i\sin\pi) = 4e^{i\pi}$.
\item $z = 2\left(\cos\left(-\dfrac{\pi}{6}\right) + i\sin\left(-\dfrac{\pi}{6}\right)\right) = 2\left(\dfrac{\sqrt{3}}{2} - \dfrac{i}{2}\right) = \sqrt{3} - i$.
\item On a $1 + i = \sqrt{2}\,e^{i\frac{\pi}{4}}$ et $\sqrt{3} - i = 2e^{-i\frac{\pi}{6}}$. Donc :
\[ z = \frac{\sqrt{2}\,e^{i\frac{\pi}{4}}}{2e^{-i\frac{\pi}{6}}} = \frac{\sqrt{2}}{2}\,e^{i\left(\frac{\pi}{4} + \frac{\pi}{6}\right)} = \frac{\sqrt{2}}{2}\,e^{i\frac{5\pi}{12}}. \]
D'autre part, sous forme algébrique, en multipliant par le conjugué du dénominateur :
\[ z = \frac{(1 + i)(\sqrt{3} + i)}{(\sqrt{3} - i)(\sqrt{3} + i)} = \frac{\sqrt{3} + i + i\sqrt{3} + i^2}{3 + 1} = \frac{\sqrt{3} - 1}{4} + i\,\frac{\sqrt{3} + 1}{4}. \]
Par identification des parties réelle et imaginaire avec $z = \dfrac{\sqrt{2}}{2}\cos\dfrac{5\pi}{12} + i\dfrac{\sqrt{2}}{2}\sin\dfrac{5\pi}{12}$ :
\[ \cos\frac{5\pi}{12} = \frac{\sqrt{3} - 1}{4} \times \frac{2}{\sqrt{2}} = \frac{\sqrt{3} - 1}{2\sqrt{2}} = \frac{\sqrt{6} - \sqrt{2}}{4} \;; \qquad \sin\frac{5\pi}{12} = \frac{\sqrt{3} + 1}{2\sqrt{2}} = \frac{\sqrt{6} + \sqrt{2}}{4}. \]
\end{enumerate}
\end{corrige}

\begin{corrige}[Exercice 6 — Ensembles de points : conditions sur le module]
\begin{enumerate}
\item $|z - (1 - i)| = 2$ s'écrit $\Omega M = 2$ avec $\Omega$ d'affixe $1 - i$, c'est-à-dire $\Omega(1\,;\,-1)$. $\mathcal{E}_1$ est le \textbf{cercle} de centre $\Omega(1\,;\,-1)$ et de rayon $2$. Équation cartésienne : $(x - 1)^2 + (y + 1)^2 = 4$. Voir la première figure de l'énoncé.
\item $\Omega M \leq 2$ : $\mathcal{E}_2$ est le \textbf{disque fermé} de centre $\Omega(1\,;\,-1)$ et de rayon $2$ (cercle et intérieur). Voir la deuxième figure.
\item Posons $A(3)$ et $B(-1)$. La condition s'écrit $\dfrac{|z - 3|}{|z + 1|} = 1$, soit $AM = BM$ (le point $B$ est exclu car le dénominateur s'y annule, mais il ne vérifie de toute façon pas l'égalité). $\mathcal{E}_3$ est la \textbf{médiatrice} du segment $[AB]$. Le milieu de $[AB]$ est $I(1\,;\,0)$ et $[AB]$ est horizontal, donc la médiatrice est la droite verticale d'équation $x = 1$. Vérification algébrique : $|z - 3|^2 = |z + 1|^2$ donne $(x-3)^2 + y^2 = (x+1)^2 + y^2$, soit $-6x + 9 = 2x + 1$, d'où $x = 1$. Voir la troisième figure.
\end{enumerate}
\end{corrige}

\begin{attention}
Le cas $\left|\dfrac{z - a}{z - b}\right| = 1$ donne toujours une \textbf{médiatrice} ; le cercle (de Apollonius) n'apparaît que pour un rapport $\alpha \neq 1$.
\end{attention}

\begin{corrige}[Exercice 7 — Linéarisation et primitive]
\begin{enumerate}
\item D'après les formules d'Euler, $\sin x = \dfrac{e^{ix} - e^{-ix}}{2i}$. Posons $u = e^{ix}$, de sorte que $e^{-ix} = \bar{u}$ et $u\bar{u} = |u|^2 = 1$. Alors, puisque $(2i)^5 = 32i^5 = 32i$ :
\[ \sin^5 x = \frac{(u - \bar{u})^5}{(2i)^5} = \frac{(u - \bar{u})^5}{32i}. \]
Par la formule du binôme de Newton :
\[ (u - \bar{u})^5 = u^5 - 5u^4\bar{u} + 10u^3\bar{u}^2 - 10u^2\bar{u}^3 + 5u\bar{u}^4 - \bar{u}^5. \]
Simplifions chaque terme grâce à $u\bar{u} = 1$ :
\[ u^4\bar{u} = u^3(u\bar{u}) = u^3 \;; \quad u^3\bar{u}^2 = u(u\bar{u})^2 = u \;; \quad u^2\bar{u}^3 = \bar{u}(u\bar{u})^2 = \bar{u} \;; \quad u\bar{u}^4 = \bar{u}^3(u\bar{u}) = \bar{u}^3. \]
En regroupant les termes conjugués deux à deux :
\[ (u - \bar{u})^5 = \left(u^5 - \bar{u}^5\right) - 5\left(u^3 - \bar{u}^3\right) + 10\left(u - \bar{u}\right). \]
Or, pour tout entier $k$, $u^k - \bar{u}^k = e^{ikx} - e^{-ikx} = 2i\sin(kx)$. Par suite :
\[ \sin^5 x = \frac{2i\sin 5x - 5 \times 2i\sin 3x + 10 \times 2i\sin x}{32i} = \frac{1}{16}\left(\sin 5x - 5\sin 3x + 10\sin x\right). \]
\item Une primitive de $x \mapsto \sin(kx)$ est $x \mapsto -\dfrac{\cos(kx)}{k}$. Donc une primitive de $x \mapsto \sin^5 x$ sur $\mathbb{R}$ est :
\[ F(x) = \frac{1}{16}\left(-\frac{\cos 5x}{5} + \frac{5\cos 3x}{3} - 10\cos x\right) = -\frac{\cos 5x}{80} + \frac{5\cos 3x}{48} - \frac{5\cos x}{8}. \]
\end{enumerate}
\end{corrige}

\begin{corrige}[Exercice 8 — Racines quatrièmes et configuration]
\begin{enumerate}
\item $-16$ est un réel négatif de module $16$, donc $-16 = 16e^{i\pi}$.
\item $z^4 = -16$ équivaut à $z^4 = 16e^{i\pi}$. Les solutions sont les nombres $z_k = 16^{1/4}e^{i\left(\frac{\pi}{4} + \frac{2k\pi}{4}\right)} = 2e^{i\left(\frac{\pi}{4} + \frac{k\pi}{2}\right)}$ pour $k \in \{0, 1, 2, 3\}$ :
\[ z_0 = 2e^{i\frac{\pi}{4}}, \quad z_1 = 2e^{i\frac{3\pi}{4}}, \quad z_2 = 2e^{i\frac{5\pi}{4}}, \quad z_3 = 2e^{i\frac{7\pi}{4}}. \]
Sous forme algébrique :
\[ z_0 = \sqrt{2} + i\sqrt{2}, \quad z_1 = -\sqrt{2} + i\sqrt{2}, \quad z_2 = -\sqrt{2} - i\sqrt{2}, \quad z_3 = \sqrt{2} - i\sqrt{2}. \]
\item Les quatre points $M_0, M_1, M_2, M_3$ sont sur le cercle de centre $O$ et de rayon $2$, aux arguments $\frac{\pi}{4}, \frac{3\pi}{4}, \frac{5\pi}{4}, \frac{7\pi}{4}$ : voir la figure de l'énoncé.
\item Les quatre points sont sur le cercle de centre $O$ et de rayon $2$, et leurs arguments diffèrent deux à deux de $\dfrac{\pi}{2}$ : ce sont les sommets d'un polygone régulier à 4 côtés, donc d'un \textbf{carré} de centre $O$. Calculons le côté :
\[ z_1 - z_0 = \left(-\sqrt{2} - \sqrt{2}\right) + i\left(\sqrt{2} - \sqrt{2}\right) = -2\sqrt{2}, \]
donc $M_0M_1 = |z_1 - z_0| = |-2\sqrt{2}| = 2\sqrt{2}$. (Vérification : le côté d'un carré inscrit dans un cercle de rayon $R$ vaut $R\sqrt{2} = 2\sqrt{2}$ : cohérent.)
\end{enumerate}
\end{corrige}

\begin{corrige}[Exercice 9 — Type Bac — forme exponentielle et valeurs exactes]
\textbf{Barème indicatif (sur 5 pts) :} Q1 : 1 pt ; Q2 : 2 pts ; Q3 : 1 pt ; Q4 : 1 pt.

\begin{enumerate}
\item \begin{enumerate}
\item $|z_1| = \sqrt{2}$ et $z_1 = \sqrt{2}\left(\dfrac{\sqrt{2}}{2} + i\dfrac{\sqrt{2}}{2}\right)$, donc $\arg(z_1) = \dfrac{\pi}{4}\;[2\pi]$ et $z_1 = \sqrt{2}\,e^{i\frac{\pi}{4}}$.
\item $|z_2| = 2$ et $z_2 = 2\left(\dfrac{\sqrt{3}}{2} - \dfrac{i}{2}\right)$, donc $\arg(z_2) = -\dfrac{\pi}{6}\;[2\pi]$ et $z_2 = 2e^{-i\frac{\pi}{6}}$.
\end{enumerate}
\item \begin{enumerate}
\item $z_1z_2 = (1 + i)(\sqrt{3} - i) = (\sqrt{3} + 1) + i(\sqrt{3} - 1)$.
\item $z_1z_2 = \sqrt{2} \times 2 \times e^{i\left(\frac{\pi}{4} - \frac{\pi}{6}\right)} = 2\sqrt{2}\,e^{i\frac{\pi}{12}}$.
\item On a donc $z_1z_2 = 2\sqrt{2}\cos\dfrac{\pi}{12} + 2i\sqrt{2}\sin\dfrac{\pi}{12}$. Par identification avec la forme algébrique :
\[ \cos\frac{\pi}{12} = \frac{\sqrt{3} + 1}{2\sqrt{2}} = \frac{\sqrt{6} + \sqrt{2}}{4} \;; \qquad \sin\frac{\pi}{12} = \frac{\sqrt{3} - 1}{2\sqrt{2}} = \frac{\sqrt{6} - \sqrt{2}}{4}. \]
\end{enumerate}
\item \begin{enumerate}
\item $Z = \dfrac{z_1}{z_2} = \dfrac{\sqrt{2}}{2}e^{i\left(\frac{\pi}{4} + \frac{\pi}{6}\right)} = \dfrac{\sqrt{2}}{2}e^{i\frac{5\pi}{12}}$.
\item $Z^{12} = \left(\dfrac{\sqrt{2}}{2}\right)^{12} e^{i \times 12 \times \frac{5\pi}{12}} = \dfrac{2^6}{2^{12}}e^{i5\pi} = \dfrac{1}{64}e^{i\pi} = -\dfrac{1}{64}$. $Z^{12}$ est un \textbf{réel strictement négatif}.
\end{enumerate}
\item La transformation est la multiplication par $z_1$ (similitude directe de centre $O$, de rapport $\sqrt{2}$ et d'angle $\frac{\pi}{4}$) : l'image $M'$ a pour affixe $z' = z_1 z_2 = (\sqrt{3} + 1) + i(\sqrt{3} - 1)$, soit $M'\left(\sqrt{3} + 1\,;\,\sqrt{3} - 1\right) \approx (2{,}73\,;\,0{,}73)$. On place $M(\sqrt{3}\,;\,-1) \approx (1{,}73\,;\,-1)$ et $M'$ dans le repère (unité 2 cm) ; on vérifie que $OM' = \sqrt{2}\,OM = 2\sqrt{2}$ et que $\left(\vec{u}, \overrightarrow{OM'}\right) = \dfrac{\pi}{12}\;[2\pi]$.
\end{enumerate}
\end{corrige}

\begin{attention}[Points de vigilance]
Citer les propriétés utilisées : module et argument d'un produit ($|z_1z_2| = |z_1||z_2|$, $\arg(z_1z_2) = \arg(z_1) + \arg(z_2)\;[2\pi]$) et d'un quotient. L'identification des parties réelle et imaginaire doit être explicitée pour obtenir les valeurs exactes.
\end{attention}

\begin{corrige}[Exercice 10 — Type Bac — géométrie, quotient d'affixes et ligne de niveau]
\textbf{Barème indicatif (sur 6 pts) :} Q1 : 0,5 pt ; Q2 : 2 pts ; Q3 : 2,5 pts ; Q4 : 1 pt.

\begin{enumerate}
\item $A(1\,;\,0)$, $B(0\,;\,1)$, $C(-1\,;\,-1)$ : voir la figure de l'énoncé.
\item \begin{enumerate}
\item $q = \dfrac{z_C - z_A}{z_B - z_A} = \dfrac{-1 - i - 1}{i - 1} = \dfrac{-2 - i}{-1 + i}$. On multiplie numérateur et dénominateur par le conjugué du dénominateur :
\[ q = \frac{(-2 - i)(-1 - i)}{(-1 + i)(-1 - i)} = \frac{2 + 2i + i + i^2}{1 + 1} = \frac{1 + 3i}{2} = \frac{1}{2} + \frac{3}{2}i. \]
\item $|q| = \sqrt{\dfrac{1}{4} + \dfrac{9}{4}} = \sqrt{\dfrac{10}{4}} = \dfrac{\sqrt{10}}{2}$. Vérification par les modules : $\dfrac{|z_C - z_A|}{|z_B - z_A|} = \dfrac{|-2 - i|}{|-1 + i|} = \dfrac{\sqrt{5}}{\sqrt{2}} = \dfrac{\sqrt{10}}{2}$. Un argument $\theta$ de $q$ vérifie $\cos\theta = \dfrac{1/2}{\sqrt{10}/2} = \dfrac{1}{\sqrt{10}}$ et $\sin\theta = \dfrac{3/2}{\sqrt{10}/2} = \dfrac{3}{\sqrt{10}}$ ; comme $\cos\theta > 0$ et $\sin\theta > 0$, $\theta$ est dans le premier quadrant et $\theta = \arctan 3$.
\item On a $\dfrac{AC}{AB} = |q| = \dfrac{\sqrt{10}}{2}$ et $\left(\overrightarrow{AB}, \overrightarrow{AC}\right) = \arg(q) = \arctan 3 \approx 71{,}6°\;[2\pi]$. Calculons les trois côtés : $AB = |z_B - z_A| = \sqrt{2}$, $AC = |z_C - z_A| = \sqrt{5}$, $BC = |z_C - z_B| = |-1 - 2i| = \sqrt{5}$. Comme $AC = BC = \sqrt{5}$, le triangle $ABC$ est \textbf{isocèle en $C$}.
\end{enumerate}
\item \begin{enumerate}
\item Pour $M(z)$ distinct de $A$ et de $B$ :
\[ \arg\left(\frac{z - 1}{z + i}\right) = \arg\left(\frac{z - z_A}{z - z_B}\right) = \left(\overrightarrow{MB}, \overrightarrow{MA}\right)\;[2\pi]. \]
La condition $\arg\left(\dfrac{z - 1}{z + i}\right) = \dfrac{\pi}{3}\;[\pi]$ équivaut donc à $\left(\overrightarrow{MB}, \overrightarrow{MA}\right) = \dfrac{\pi}{3}\;[\pi]$, c'est-à-dire à une mesure de l'angle \textbf{de droites} $\left((MB), (MA)\right)$ égale à $\dfrac{\pi}{3}\;[\pi]$.
\item L'ensemble des points $M$ d'où l'on voit le segment $[AB]$ sous un angle de droites égal à $\dfrac{\pi}{3}\;[\pi]$ est la réunion de deux \textbf{arcs de cercle} d'extrémités $A$ et $B$ (arcs capables), symétriques par rapport à la droite $(AB)$, \textbf{privés des points $A$ et $B$}.

Déterminons les cercles supports. La corde $[AB]$ a pour longueur $AB = \sqrt{2}$ ; pour un angle inscrit de $\dfrac{\pi}{3}$, la relation $AB = 2R\sin\dfrac{\pi}{3}$ donne :
\[ R = \frac{\sqrt{2}}{2 \times \frac{\sqrt{3}}{2}} = \frac{\sqrt{2}}{\sqrt{3}} = \frac{\sqrt{6}}{3}. \]
Le centre $\Omega$ est sur la médiatrice de $[AB]$. Comme $A(1\,;\,0)$ et $B(0\,;\,1)$, le milieu est $I\left(\dfrac{1}{2}\,;\,\dfrac{1}{2}\right)$ et la médiatrice est la droite d'équation $y = x$. La distance du centre au milieu vaut :
\[ d = R\cos\frac{\pi}{3} = \frac{\sqrt{6}}{3} \times \frac{1}{2} = \frac{\sqrt{6}}{6}. \]
Un vecteur unitaire directeur de la médiatrice est $\dfrac{1}{\sqrt{2}}(1\,;\,1)$, donc les deux centres sont $I \pm \dfrac{d}{\sqrt{2}}(1\,;\,1)$ avec $\dfrac{d}{\sqrt{2}} = \dfrac{\sqrt{6}}{6\sqrt{2}} = \dfrac{\sqrt{3}}{6}$ :
\[ \Omega_1\left(\frac{1}{2} - \frac{\sqrt{3}}{6}\,;\,\frac{1}{2} - \frac{\sqrt{3}}{6}\right) \quad \text{et} \quad \Omega_2\left(\frac{1}{2} + \frac{\sqrt{3}}{6}\,;\,\frac{1}{2} + \frac{\sqrt{3}}{6}\right). \]
$\Gamma$ est la réunion des deux arcs d'extrémités $A$ et $B$ portés par ces cercles, \textbf{points $A$ et $B$ exclus} (le quotient n'est pas défini en $B$, et l'argument n'a pas de sens en $A$).
\item Construction : tracer les deux cercles de centres $\Omega_1$, $\Omega_2$ et de rayon $\dfrac{\sqrt{6}}{3} \approx 0{,}82$ (soit environ 1,6 cm à l'échelle 2 cm), et ne retenir que les arcs d'extrémités $A$ et $B$ (voir figure).
\end{enumerate}
\item $\left|\dfrac{z - 1}{z + i}\right| = 2 \iff AM = 2\,BM$ avec $M \neq B$. Élevons au carré avec $M(x\,;\,y)$ :
\[ (x-1)^2 + y^2 = 4\left(x^2 + (y+1)^2\right) \iff 3x^2 + 3y^2 + 2x + 8y + 3 = 0 \iff x^2 + y^2 + \frac{2}{3}x + \frac{8}{3}y + 1 = 0. \]
Complétons les carrés : $\left(x + \dfrac{1}{3}\right)^2 + \left(y + \dfrac{4}{3}\right)^2 = \dfrac{1}{9} + \dfrac{16}{9} - 1 = \dfrac{8}{9}$. C'est le \textbf{cercle} (cercle d'Apollonius) de centre $\Omega\left(-\dfrac{1}{3}\,;\,-\dfrac{4}{3}\right)$ et de rayon $\sqrt{\dfrac{8}{9}} = \dfrac{2\sqrt{2}}{3}$.
\end{enumerate}
\end{corrige}

\begin{attention}[Points de vigilance]
Pour la ligne de niveau, exiger : l'interprétation en angle de vecteurs, la condition $z \neq -i$ (et $z \neq 1$ pour que l'argument existe), la mention que $A$ et $B$ sont exclus de $\Gamma$, et la distinction entre le cercle support et l'arc. Pour le rapport de modules égal à 2 ($\neq 1$), ne pas conclure à tort à une médiatrice.
\end{attention}

\begin{corrige}[Exercice 11 — Type Bac — problème de synthèse : rotations et triangles équilatéraux]
\textbf{Barème indicatif (sur 6 pts) :} Q1 : 1 pt ; Q2 : 2 pts ; Q3 : 2 pts ; Q4 : 1 pt.

\begin{enumerate}
\item \begin{enumerate}
\item $z_C - z_A = 1 + i\sqrt{3}$ ; $|z_C - z_A| = \sqrt{1 + 3} = 2$ et $1 + i\sqrt{3} = 2\left(\dfrac{1}{2} + i\dfrac{\sqrt{3}}{2}\right) = 2e^{i\frac{\pi}{3}}$, donc $\arg(z_C - z_A) = \dfrac{\pi}{3}\;[2\pi]$. Ainsi $AC = 2$, $AB = |z_B - z_A| = 4$, et $\left(\overrightarrow{AB}, \overrightarrow{AC}\right) = \dfrac{\pi}{3}\;[2\pi]$. De plus $BC = |z_C - z_B| = |-3 + i\sqrt{3}| = \sqrt{9 + 3} = 2\sqrt{3}$. On constate que $AB^2 = 16 = 4 + 12 = AC^2 + BC^2$ : par la réciproque du théorème de Pythagore, le triangle $ABC$ est \textbf{rectangle en $C$} (c'est un demi-triangle équilatéral : angles $\frac{\pi}{6}, \frac{\pi}{3}, \frac{\pi}{2}$).
\item Voir la figure de l'énoncé.
\end{enumerate}
\item \begin{enumerate}
\item $D$ est l'image de $B$ par la rotation de centre $A(0)$ et d'angle $\dfrac{\pi}{3}$ :
\[ z_D - z_A = e^{i\frac{\pi}{3}}(z_B - z_A) \implies z_D = \left(\frac{1}{2} + i\frac{\sqrt{3}}{2}\right) \times 4 = 2 + 2i\sqrt{3}. \]
\item $E$ est l'image de $C$ par la rotation de centre $B(4)$ et d'angle $\dfrac{\pi}{3}$ :
\[ z_E - 4 = e^{i\frac{\pi}{3}}(z_C - 4) = \left(\frac{1}{2} + i\frac{\sqrt{3}}{2}\right)(-3 + i\sqrt{3}). \]
Développons : $\dfrac{1}{2} \times (-3) + \dfrac{1}{2} \times i\sqrt{3} + i\dfrac{\sqrt{3}}{2} \times (-3) + i\dfrac{\sqrt{3}}{2} \times i\sqrt{3} = -\dfrac{3}{2} - \dfrac{3}{2} + i\left(\dfrac{\sqrt{3}}{2} - \dfrac{3\sqrt{3}}{2}\right) = -3 - i\sqrt{3}$. Donc $z_E = 4 - 3 - i\sqrt{3} = 1 - i\sqrt{3}$.
\item Voir la figure : $D(2\,;\,2\sqrt{3})$ et $E(1\,;\,-\sqrt{3})$.
\end{enumerate}
\item \begin{enumerate}
\item Affixe de $\overrightarrow{AE}$ : $z_E - z_A = 1 - i\sqrt{3}$. Affixe de $\overrightarrow{DC}$ : $z_C - z_D = (1 + i\sqrt{3}) - (2 + 2i\sqrt{3}) = -1 - i\sqrt{3}$.
\item $AE = |1 - i\sqrt{3}| = \sqrt{1 + 3} = 2$ et $DC = |-1 - i\sqrt{3}| = 2$, donc $AE = DC$. De plus :
\[ \left(\overrightarrow{AE}, \overrightarrow{DC}\right) = \arg\left(\frac{-1 - i\sqrt{3}}{1 - i\sqrt{3}}\right)\;[2\pi]. \]
Multiplions par le conjugué : $\dfrac{(-1 - i\sqrt{3})(1 + i\sqrt{3})}{(1 - i\sqrt{3})(1 + i\sqrt{3})} = \dfrac{-1 - i\sqrt{3} - i\sqrt{3} - i^2 \times 3}{1 + 3} = \dfrac{2 - 2i\sqrt{3}}{4} = \dfrac{1}{2} - i\dfrac{\sqrt{3}}{2} = e^{-i\frac{\pi}{3}}$. Donc :
\[ \left(\overrightarrow{AE}, \overrightarrow{DC}\right) = -\frac{\pi}{3}\;[2\pi], \]
et l'angle des droites $(AE)$ et $(DC)$ vaut $\dfrac{\pi}{3}$.
\end{enumerate}
\item \textbf{Généralisation.} Soient $a$, $b$, $c$ les affixes de $A$, $B$, $C$. Posons $\omega = e^{i\frac{\pi}{3}}$. Les triangles $ABD$ et $BCE$ étant équilatéraux directs, construits à l'extérieur de $ABC$ :
\[ d = a + \omega(b - a) \qquad \text{et} \qquad e = b + \omega(c - b). \]
Nous allons établir la relation $c - d = e^{-i\frac{\pi}{3}}(e - a)$, c'est-à-dire $c - d = \bar{\omega}(e - a)$. Deux identités seront utiles :
\[ \omega + \bar{\omega} = 2\cos\frac{\pi}{3} = 1 \quad \text{donc} \quad \bar{\omega} = 1 - \omega \;; \qquad \omega^2 = e^{i\frac{2\pi}{3}} = \omega - 1 \]
(la seconde se vérifie : $\omega^2 - \omega + 1 = \left(-\dfrac{1}{2} + i\dfrac{\sqrt{3}}{2}\right) - \left(\dfrac{1}{2} + i\dfrac{\sqrt{3}}{2}\right) + 1 = 0$).

Calculons alors :
\begin{align*}
\bar{\omega}(e - a) &= (1 - \omega)\left[(b - a) + \omega(c - b)\right] \\
&= (b - a) + \omega(c - b) - \omega(b - a) - \omega^2(c - b).
\end{align*}
Or $-\omega^2(c - b) = -(\omega - 1)(c - b) = -\omega(c - b) + (c - b)$. Les termes en $\omega(c - b)$ se simplifient et il reste :
\[ \bar{\omega}(e - a) = (b - a) + (c - b) - \omega(b - a) = (c - a) - \omega(b - a) = c - d. \]
La relation $c - d = e^{-i\frac{\pi}{3}}(e - a)$ est donc établie pour \textbf{tout} triangle $ABC$. Conséquences :
\begin{itemize}
\item en prenant les modules : $|c - d| = \left|e^{-i\frac{\pi}{3}}\right| \times |e - a| = |e - a|$, c'est-à-dire $DC = AE$ ;
\item en prenant les arguments : $\left(\overrightarrow{AE}, \overrightarrow{DC}\right) = \arg(c - d) - \arg(e - a) = -\dfrac{\pi}{3}\;[2\pi]$.
\end{itemize}
Ainsi, quelle que soit la forme du triangle initial, les segments $[AE]$ et $[DC]$ ont \textbf{même longueur} et leurs supports font un \textbf{angle de $\dfrac{\pi}{3}$}.
\end{enumerate}
\end{corrige}

\begin{attention}[Points de vigilance]
Bien justifier le sens « direct » des triangles équilatéraux par le choix de l'angle $+\frac{\pi}{3}$ (et vérifier sur la figure que $D$ et $E$ sont bien à l'extérieur du triangle $ABC$). Dans la généralisation, les identités $\omega + \bar{\omega} = 1$ et $\omega^2 = \omega - 1$ sont les clés du calcul : les citer explicitement et les démontrer brièvement.
\end{attention}

\newpage

\coursec{Fiche bilan}

\begin{popfigure}
\begin{tikzpicture}[mindmap, grow cyclic, every node/.style={concept, font=\small},
  root concept/.append style={concept color=popblue, font=\bfseries},
  level 1/.append style={level distance=3.4cm, sibling angle=60},
  level 2/.append style={level distance=2.2cm, font=\scriptsize}]
\node[root concept]{Nombres complexes : approche géométrique}
  child[concept color=poporange]{ node {Plan complexe, affixes}
    child { node {$z_{\vec{u}}=z_B-z_A$} }
    child { node {$|z|=OM$} }
  }
  child[concept color=popgreen]{ node {Modules et ensembles}
    child { node {$|z-a|=\alpha$ : cercle} }
    child { node {$\left|\frac{z-a}{z-b}\right|=\alpha$} }
  }
  child[concept color=poppurple]{ node {Argument, formes}
    child { node {$z=r(\cos\theta+i\sin\theta)$} }
    child { node {$z=re^{i\theta}$} }
  }
  child[concept color=poppink]{ node {Opérations sur arguments}
    child { node {$\arg(zz')$, $\arg(z/z')$} }
    child { node {Alignement, orthogonalité} }
  }
  child[concept color=popgold]{ node {Moivre, racines}
    child { node {$(\cos\theta+i\sin\theta)^n$} }
    child { node {Racines $n$-ièmes} }
    child { node {Linéarisation} }
  }
  child[concept color=popblue!70]{ node {Lignes de niveau}
    child { node {$\arg\frac{z-a}{z-b}=\alpha$} }
    child { node {Arcs de cercle} }
  };
\end{tikzpicture}
\legende{Carte mentale de la leçon : les six compétences clés autour de l'interprétation géométrique des nombres complexes.}
\end{popfigure}

\begin{aretenir}[Bilan --- Définitions clés]
\begin{itemize}
  \item \cle{Affixe} : au point $M(x\,;y)$ du plan muni d'un repère orthonormé direct $(O\,;\vec{u},\vec{v})$, on associe $z_M=x+iy$ ; au vecteur $\vec{w}(a\,;b)$, on associe $z_{\vec{w}}=a+ib$. On a $z_{\overrightarrow{AB}}=z_B-z_A$.
  \item \cle{Module} : $|z|=\sqrt{a^2+b^2}=OM$ pour $z=a+ib$ ; plus généralement $|z_B-z_A|=AB$ (distance entre deux points).
  \item \cle{Argument} : pour $z\neq 0$, tout réel $\theta$ tel que $(\vec{u}\,;\overrightarrow{OM})=\theta\ [2\pi]$. L'\cle{argument principal} est l'unique argument appartenant à $]-\pi\,;\pi]$.
  \item \cle{Forme trigonométrique} : $z=r(\cos\theta+i\sin\theta)$ avec $r=|z|>0$ ; \cle{forme exponentielle} : $z=re^{i\theta}$.
\end{itemize}
\end{aretenir}

\begin{aretenir}[Bilan --- Formules et théorèmes à connaître par c\oe ur]
\begin{center}
\begin{tabularx}{\linewidth}{|l|X|}
\hline
\rowcolor{popblueL} \textbf{Objet} & \textbf{Formule} \\
\hline
Module d'un produit / quotient & $|zz'|=|z|\,|z'|$ ; $\left|\dfrac{z}{z'}\right|=\dfrac{|z|}{|z'|}$ ($z'\neq 0$) \\
\hline
Arguments & $\arg(zz')=\arg z+\arg z'\ [2\pi]$ ; $\arg\dfrac{z}{z'}=\arg z-\arg z'\ [2\pi]$ ($z\neq 0$, $z'\neq 0$) ; $\arg(z^n)=n\arg z\ [2\pi]$ ($z\neq 0$, $n\in\mathbb{Z}$) \\
\hline
Conjugué, opposé & $\arg\bar z=-\arg z\ [2\pi]$ ; $\arg(-z)=\arg z+\pi\ [2\pi]$ ($z\neq 0$) \\
\hline
Moivre & $(\cos\theta+i\sin\theta)^n=\cos(n\theta)+i\sin(n\theta)$, $n\in\mathbb{Z}$ \\
\hline
Euler & $\cos\theta=\dfrac{e^{i\theta}+e^{-i\theta}}{2}$ ; $\sin\theta=\dfrac{e^{i\theta}-e^{-i\theta}}{2i}$ \\
\hline
Racines $n$-ièmes de $re^{i\theta}$ & $z_k=\sqrt[n]{r}\,e^{i\left(\frac{\theta}{n}+\frac{2k\pi}{n}\right)}$, $k\in\{0,1,\dots,n-1\}$ \\
\hline
Angle de vecteurs & $(\overrightarrow{AB}\,;\overrightarrow{CD})=\arg\dfrac{z_D-z_C}{z_B-z_A}\ [2\pi]$ ($A\neq B$, $C\neq D$) \\
\hline
\end{tabularx}
\end{center}
\end{aretenir}

\begin{aretenir}[Bilan --- Ensembles de points et lignes de niveau]
\begin{itemize}
  \item $|z-a|=\alpha$ ($\alpha>0$) : \cle{cercle} de centre $A(a)$ et de rayon $\alpha$ ; $|z-a|\leq\alpha$ : \cle{disque fermé} de centre $A$ et de rayon $\alpha$.
  \item $|z-a|=|z-b|$ ($A\neq B$) : \cle{médiatrice} du segment $[AB]$.
  \item $\left|\dfrac{z-a}{z-b}\right|=\alpha$ ($\alpha>0$, $A\neq B$) : si $\alpha=1$, médiatrice de $[AB]$ ; si $\alpha\neq 1$, \cle{cercle} (cercle d'Apollonius) — chercher les deux points de la droite $(AB)$ vérifiant la condition : ils forment un diamètre du cercle.
  \item $\arg\dfrac{z-a}{z-b}=\alpha\ [2\pi]$ : \cle{arc de cercle} (ouvert) d'extrémités $A$ et $B$ ; si $\alpha=\dfrac{\pi}{2}\ [\pi]$, cercle de diamètre $[AB]$ privé de $A$ et $B$ ; si $\alpha=0\ [\pi]$, droite $(AB)$ privée de $A$ et $B$.
\end{itemize}
\end{aretenir}

\begin{methode}[Bilan --- Méthodes express]
\begin{itemize}
  \item \emph{Algébrique $\to$ trigonométrique} : calculer $r=|z|$, factoriser $r$ dans $z$, identifier $\cos\theta$ et $\sin\theta$ à l'aide des valeurs remarquables, conclure avec un angle de référence.
  \item \emph{Linéariser} $\cos^n x$ ou $\sin^n x$ : remplacer par les formules d'Euler, développer avec le binôme de Newton, regrouper les termes deux à deux via $e^{ikx}+e^{-ikx}=2\cos(kx)$.
  \item \emph{Racines $n$-ièmes} : écrire le nombre sous forme exponentielle, appliquer la formule ; les images forment un \cle{polygone régulier} à $n$ côtés inscrit dans le cercle de centre $O$ et de rayon $\sqrt[n]{r}$.
  \item \emph{Géométrie via les complexes} : traduire les longueurs par des modules et les angles par des arguments ; $A$, $B$, $C$ alignés $\Leftrightarrow \dfrac{z_C-z_A}{z_B-z_A}\in\mathbb{R}$ ; $(AB)\perp(CD) \Leftrightarrow \dfrac{z_D-z_C}{z_B-z_A}$ est imaginaire pur.
\end{itemize}
\end{methode}

\begin{attention}[Pièges à éviter absolument]
\begin{itemize}
  \item L'argument n'est défini que pour un complexe \textbf{non nul} : avant d'écrire $\arg\dfrac{z-a}{z-b}$, vérifier $z\neq a$ et $z\neq b$ ; c'est pourquoi $A$ et $B$ sont toujours exclus des lignes de niveau.
  \item Ne pas confondre $|z-a|=\alpha$ (cercle, une courbe) et $|z-a|\leq\alpha$ (disque, une surface).
  \item Les égalités entre arguments sont toujours données \textbf{modulo $2\pi$} (ou modulo $\pi$ pour les angles de droites) : ne jamais les écrire comme des égalités de réels.
  \item Dans la forme trigonométrique $z=r(\cos\theta+i\sin\theta)$, le même angle $\theta$ doit apparaître dans le cosinus \emph{et} le sinus, avec $r>0$ : $z=-2\left(\cos\frac{\pi}{3}+i\sin\frac{\pi}{3}\right)$ n'est \textbf{pas} sous forme trigonométrique (il faut écrire $z=2\left(\cos\frac{4\pi}{3}+i\sin\frac{4\pi}{3}\right)$).
  \item Pour les racines $n$-ièmes, il y a exactement $n$ racines distinctes : ne pas oublier les valeurs $k=1,2,\dots,n-1$.
\end{itemize}
\end{attention}

\begin{aretenir}[Checklist avant le Bac]
\begin{itemize}
  \item[$\square$] Je sais lire et placer l'affixe d'un point et d'un vecteur dans le plan complexe.
  \item[$\square$] Je sais que $|z_B-z_A|=AB$ et je traduis une condition sur les modules en ensemble de points.
  \item[$\square$] Je distingue cercle ($|z-a|=\alpha$), disque ($|z-a|\leq\alpha$) et médiatrice ($|z-a|=|z-b|$).
  \item[$\square$] Je sais passer de la forme algébrique à la forme trigonométrique ou exponentielle, et réciproquement.
  \item[$\square$] Je connais les valeurs remarquables de $\cos$ et $\sin$ ($\frac{\pi}{6}$, $\frac{\pi}{4}$, $\frac{\pi}{3}$, \dots) sans hésiter.
  \item[$\square$] J'applique correctement $\arg(zz')$, $\arg(z/z')$, $\arg(z^n)$ modulo $2\pi$, avec $z\neq 0$ et $z'\neq 0$ pour le quotient.
  \item[$\square$] Je sais énoncer et utiliser la formule de Moivre et les formules d'Euler pour linéariser $\cos^n x$ et $\sin^n x$.
  \item[$\square$] Je sais déterminer les racines $n$-ièmes d'un complexe et représenter le polygone régulier associé.
  \item[$\square$] Je sais caractériser la ligne de niveau $\arg\dfrac{z-a}{z-b}=\alpha$ (arc de cercle, cercle, droite) sans oublier d'exclure $A$ et $B$.
  \item[$\square$] Je sais prouver alignement, orthogonalité ou nature d'un triangle (isocèle, équilatéral, rectangle) par le calcul complexe.
  \item[$\square$] Je vérifie toujours que les complexes sous $\arg$ sont \textbf{non nuls} avant de manipuler les arguments.
  \item[$\square$] Je rédige proprement : module d'abord, argument ensuite, conclusion géométrique explicite.
\end{itemize}
\end{aretenir}

\findecours

\end{document}
